\documentclass[12pt,a4paper]{article}

% ==== Русский язык и кодировки (pdfLaTeX) ====
\usepackage[T2A]{fontenc}
\usepackage[utf8]{inputenc}
\usepackage[russian]{babel}

% ==== Геометрия ====
\usepackage{geometry}
\geometry{margin=2cm}

% ==== Встроенные шрифты TeX-дистрибутива ====

% ==== Внешний вид абзацев ====

% ==== Математика ====
\usepackage{amsmath, amssymb, amsthm}
\numberwithin{equation}{section}
\usepackage{mathtools}
\newcommand{\xleftrightarrows}[2][]{%
  \overset{#2}{\underset{#1}{\rightleftarrows}}%
}

% ==== Графика/рисунки ====
\usepackage{graphicx}
\graphicspath{{images/}}
\usepackage{float}
\usepackage[export]{adjustbox}

% ==== TikZ/PGFPlots ====
\usepackage{tikz}
\usetikzlibrary{arrows.meta,angles,quotes,calc}
\usepackage{pgfplots}
\pgfplotsset{compat=1.18}

% ==== Остальные пакеты ====
\usepackage{xcolor}
\usepackage{mdframed}
\usepackage{fancyhdr}
\usepackage{titlesec}
% Заголовок не остаётся один внизу страницы: если места меньше 10 строк,
% он переносится на следующую страницу вместе с первой коробочкой.
\usepackage{needspace}
\newcommand{\sectionbreak}{\Needspace{10\baselineskip}}
\newcommand{\subsectionbreak}{\Needspace{10\baselineskip}}
\newcommand{\subsubsectionbreak}{\Needspace{10\baselineskip}}
\usepackage{hyperref}
\usepackage[most]{tcolorbox}
\usepackage{enumitem}
\usepackage{cancel}
\usepackage{parskip}

% ==== Цвета ====
\definecolor{HSEblue}{HTML}{004C97}
\definecolor{HSElightblue}{HTML}{5B8CDE}
\definecolor{HSEdarkblue}{HTML}{00315E}

\definecolor{HSEgreen}{HTML}{007F5F}
\definecolor{HSElightgreen}{HTML}{00A47A}
\definecolor{HSEred}{HTML}{B80C09}
\definecolor{HSEorange}{HTML}{F26D00}
\definecolor{HSEpurple}{HTML}{6E56DB}
\definecolor{HSEteal}{HTML}{008B8B}
\definecolor{HSEgold}{HTML}{B8860B}
\definecolor{HSEmagenta}{HTML}{D63384}

\definecolor{HSEblack}{HTML}{000000}
\definecolor{HSEdarkgray}{HTML}{333333}
\definecolor{HSEgray}{HTML}{666666}
\definecolor{HSElightgray}{HTML}{999999}
\definecolor{HSEwhite}{HTML}{FFFFFF}

\definecolor{HSEbglight}{HTML}{F8F9FA}
\definecolor{HSEbglighter}{HTML}{F0F2F5}

% ==== Замены твоих "HSE-шрифтов" на стандартные переключатели ====
\newcommand{\thintext}[1]{{\mdseries #1}}
\newcommand{\blacktext}[1]{{\bfseries #1}}

% ==== Полезные обозначения из твоего исходника ====
\newcommand{\Od}[2]{O_{#1}\!\left(#2\right)}
\newcommand{\Odp}[2]{\dot O_{#1}\!\left(#2\right)}

% операторы (у babel/russian часто заняты \tg,\sh,\ch — освобождаем)
\let\sh\relax
\let\ch\relax
\DeclareMathOperator{\sgn}{sgn}
\DeclareMathOperator{\arcsh}{arcsh}
\DeclareMathOperator{\Lip}{Lip}
\renewcommand{\Re}{\operatorname{Re}}
\renewcommand{\Im}{\operatorname{Im}}
\providecommand{\NOD}{\operatorname{НОД}}
\providecommand{\NOK}{\operatorname{НОК}}

% множества функций (чтобы \UC, \Cont, \Bnd компилились)
\DeclareMathOperator{\UC}{UC}
\DeclareMathOperator{\Cont}{C}
\DeclareMathOperator{\Bnd}{Bnd}

% дополнительные обозначения для диффуров
\newcommand{\R}{\mathbb{R}}
\newcommand{\Z}{\mathbb{Z}}
\newcommand{\vect}[1]{\vec{#1}}
\newcommand{\e}{\mathrm{e}}
\newcommand{\tr}{\operatorname{tr}}
\newcommand{\dd}{\mathop{}\!\mathrm{d}}

% ==== Заголовки ====
\titleformat{\section}
  {\Large\bfseries\color{HSEblue}}
  {\thesection.}{0.5em}{}

\titleformat{\subsection}
  {\large\bfseries\color{HSEblue}}
  {\thesubsection.}{0.5em}{}

% ==== Колонтитулы ====
\pagestyle{fancy}
\fancyhf{}
\fancyhead[L]{\bfseries\nouppercase{\leftmark}}
\fancyhead[R]{\thepage}
\setlength{\headheight}{15pt}

\renewcommand{\headrulewidth}{0.4pt}

% ==== Ссылки ====
\hypersetup{
  colorlinks=true,
  linkcolor=HSEblue,
  urlcolor=HSEblue
}

% ==== tcolorbox базовый стиль ====
\newcommand{\sqbullet}{\raisebox{0.15ex}{\scalebox{0.4}{$\blacksquare$}}}

\tcbset{
  enhanced,
  breakable,
  boxrule=0.8pt,
  arc=6pt,
  boxsep=1.5pt,
  before upper=\par\vspace{6pt}\noindent,
  top=2pt,
  fonttitle=\bfseries,
  fontupper=\normalsize,
  attach boxed title to top left={yshift=-2mm},
  boxed title style={
    boxrule=0pt,
    frame hidden,
    interior style={fill opacity=1}
  }
}

% ===== Коробочки без нумерации =====
\tcbset{
  hsebox/.style={
    enhanced,
    breakable,
    boxrule=0.8pt,
    arc=6pt,
    boxsep=1.5pt,
    top=2pt,
    fonttitle=\bfseries,
    coltitle=white,
    attach boxed title to top left={yshift=-2mm,xshift=2mm},
    boxed title style={
      boxrule=0pt,
      frame hidden,
      interior style={fill opacity=1}
    },
    before upper={
      \setlength{\parskip}{0.7em}%
      \setlength{\parindent}{0pt}%
      \par\vspace{6pt}\ignorespaces
    }
  }
}

\newcommand{\DeclareHSEBox}[5]{%
  \newtcolorbox{#1}[1][#2]{%
    hsebox,
    colframe=#3,
    colback=HSEgray!10,
    colbacktitle=#4,
    title={##1},
    before upper app={%
      \setlist[enumerate]{label=\color{#5}\arabic*., font=\color{#5}, itemsep=0.6em, topsep=0.3em}%
      \setlist[itemize]{label=\sqbullet, leftmargin=1.2em, itemsep=0.6em, topsep=0.3em}%
    }
  }%
}

\DeclareHSEBox{definitionHSE}{Определение}{HSEblue}{HSEblue!90!black}{HSEblue}
\DeclareHSEBox{theoremHSE}{Теорема}{HSEgreen!70!black}{HSEgreen!80!black}{HSEgreen}
\DeclareHSEBox{lemmaHSE}{Лемма}{HSEteal!80!black}{HSEteal!85!black}{HSEteal}
\DeclareHSEBox{corollaryHSE}{Следствие}{HSEpurple!60!black}{HSEpurple!70!black}{HSEpurple}
\DeclareHSEBox{exampleHSE}{Пример}{HSEred!80!black}{HSEred!90!black}{HSEred}
\DeclareHSEBox{taskHSE}{Задача}{HSEmagenta!80!black}{HSEmagenta!90!black}{HSEmagenta}
\DeclareHSEBox{remarkHSE}{Замечание}{HSEorange!70!black}{HSEorange!80!black}{HSEorange}

% --- Доказательство ---
\renewenvironment{proof}[1][Доказательство]{\begin{center}\textbf{#1}\end{center}}{\hfill$\blacksquare$\par}

% ==== Единый стиль графиков ====
\pgfplotsset{
  hseaxis/.style={
    axis lines=middle,
    axis line style={draw=HSEdarkgray},
    tick style={draw=HSEdarkgray},
    tick label style={text=HSEdarkgray},
    label style={text=HSEdarkgray},
    grid style={draw=HSElightgray!45},
    minor grid style={draw=HSElightgray!25}
  },
  hseoriginal/.style={samples=201,draw=HSEgray,thick},
  hsetransformed/.style={samples=201,very thick,draw=HSEblue},
  hseguide/.style={draw=HSEgray,dashed,semithick,no markers},
  hsepoint/.style={only marks,mark=*,mark size=2pt,draw=HSEblue,fill=HSEblue}
}

% ==== Единый стиль геометрических рисунков TikZ ====
\tikzset{
  hsetriangle/.style={draw=HSEdarkgray,thick},
  hseangle/.style={draw=HSEblue,fill=HSElightblue!25,thick}
}

% === МЕТАДАННЫЕ ===
\title{\textbf{Диффуры}\\[0.3em]\large Вопросник к экзамену и материалы с лекций}
\author{Медведев Максим --- \href{https://t.me/krols}{@krols}}
\date{}

\begin{document}
\maketitle
\tableofcontents
\clearpage

\section{Базовые определения ОДУ}

\begin{definitionHSE}[Опредение ОДУ и его порядка]
Обыкновенным дифференциальным уравнением (ОДУ) называется \underline{соотношение}
\begin{equation}
F\bigl(x,y,y',\ldots,y^{(n)}\bigr)=0,
\label{eq:odu}
\end{equation}
связывающее значения независимой переменной $x\in\R$, неизвестной функции $y=y(x)$ и её производных.

Порядок старшей производной, входящей в уравнение \eqref{eq:odu}, называется \textbf{порядком ОДУ}.
\end{definitionHSE}

Слово «обыкновенное» означает, что переменная одна: неизвестная функция зависит от одной независимой переменной. Если бы их было несколько, и стояли частные производные, уравнение называлось бы уравнением в частных производных, а это отдельный курс (уравнения математической физики).

\begin{definitionHSE}[Определение решения ОДУ]
Решением ОДУ \eqref{eq:odu} на множестве $M\subset\R$ называется \underline{функция} $y=\varphi(x)$, которая \textbf{на этом множестве}:
\begin{enumerate}
  \item определена и имеет производные \(\varphi',\varphi'',\ldots,\varphi^{(n)}\) до порядка $n$ включительно;
  \item при подстановке обращает уравнение \eqref{eq:odu} в тождество:
  \[
  F\bigl(x,\varphi(x),\varphi'(x),\ldots,\varphi^{(n)}(x)\bigr)=0,
  \qquad \forall x\in M.
  \]
\end{enumerate}
\end{definitionHSE}

Решение \textbf{всегда} задаётся вместе с множеством, на котором оно рассматривается.

\begin{definitionHSE}[Определение ОДУ, разрешённого относительно старшей производной]
ОДУ порядка \(n\) называется \textit{разрешённым относительно старшей
производной}, если оно записано в виде
\[
y^{(n)}=f\bigl(x,y,y',\ldots,y^{(n-1)}\bigr)
\]
\end{definitionHSE}
Для первого порядка это \(y'=f(x,y)\). Уравнение в иной форме записи называют уравнением \textit{неразрешённого вида}.

\section[Геометрическая интерпретация ОДУ]{Геометрическая интерпретация ОДУ}

Рассмотрим ДУ первого порядка, разрешённое относительно производной:
\begin{equation}
 y'=f(x,y),
 \label{eq:first-order}
\end{equation}
где $f$ определена в некоторой области $G\subset\mathbb{R}^2$.

\begin{remarkHSE}
Под \textit{областью} мы понимаем открытое связное множество.
\end{remarkHSE}

Если график решения $y=\varphi(x)$ ДУ \eqref{eq:first-order} на множестве \(M\) проходит через $P_0=(x_0,y_0)\in G$ (то есть \(\varphi(x_0)=y_0\)), то его касательная в этой точке имеет наклон
\begin{gather*}
\tan\alpha=\varphi'(x_0)=f(x_0,\varphi(x_0))=f(x_0,y_0),\\
\varphi'(x)=f(x,\varphi(x))\quad\forall x\in M,
\end{gather*}
где $\alpha$ --- угол касательной к графику решения с положительным направлением оси $Ox$.

Поэтому уравнение \eqref{eq:first-order} задаёт в каждой точке области \(G\) направление касательной к решению \(\varphi(x)\).

\begin{center}
\begin{tikzpicture}[
  >=Stealth, line cap=round, line join=round,
  every node/.style={font=\normalsize}
]
  \definecolor{solution}{HTML}{C84949}
  \definecolor{tangent}{HTML}{2787BD}

  % Точка M_0 и наклон касательной.
  \def\xzero{3}
  \def\yzero{2.1}
  \def\slope{1.2}
  \pgfmathsetmacro{\xintercept}{\xzero-\yzero/\slope}
  \pgfmathsetmacro{\ang}{atan(\slope)}
  \coordinate (M) at (\xzero,\yzero);
  \coordinate (A) at (\xintercept,0);

  % Область G с выраженной волнистой границей и плавными впадинами.
  % Область не предполагается выпуклой.
  \begin{scope}[shift={(1.8,1.45)}, rotate=18]
    \path[draw=black!65, fill=black!2, line width=0.9pt]
      plot[smooth cycle, domain=0:359, samples=360, variable=\t]
      ({3.85*(1+0.08*sin(7*\t+225)+0.025*sin(11*\t-20))*cos(\t)},
       {2.60*(1+0.08*sin(7*\t+225)+0.025*sin(11*\t-20))*sin(\t)})
      -- cycle;
  \end{scope}
  \node at (-0.85,-0.30) {$G$};

  % Оси координат.
  \draw[->, thick] (-2.35,0) -- (6,0) node[below] {$x$};
  \draw[->, thick] (0,-1.5) -- (0,4.65) node[left] {$y$};
  \node[below left, inner sep=3pt] at (0,0) {$O$};

  % Проекции M_0 на оси.
  \draw[densely dashed, black!65]
    (0,\yzero) node[left, text=black] {$y_0$}
    -- (M) -- (\xzero,0) node[below, text=black] {$x_0$};

  % Касательная: y = y_0 + slope*(x-x_0).
  \draw[tangent, line width=1.2pt]
    (A) -- (4.3,{\yzero+\slope*(4.3-\xzero)})
    node[pos=0.92, left, inner sep=5pt] {$\ell$};

  % Кривая проходит через M_0 и имеет там производную slope.
  \draw[solution, line width=1.4pt, domain=1.25:4.65, samples=101]
    plot (\x,{\yzero+\slope*(\x-\xzero)-0.25*(\x-\xzero)^3});
  \node[solution] at (1.55,1.65) {$y=\varphi(x)$};

  % Угол касательной с положительным направлением оси Ox.
  \begin{scope}[shift={(A)}]
    \draw[tangent, line width=0.8pt]
      (0.65,0) arc[start angle=0,end angle=\ang,radius=0.65];
    \node[tangent] at ({0.95*cos(\ang/2)},{0.95*sin(\ang/2)})
      {$\alpha$};
  \end{scope}

  \fill (M) circle (2pt);
  \node[below right, inner sep=6pt, font=\normalsize] at (M)
    {$P_0=(x_0,y_0)$};

  \node at (1.8,-2.10)
    {$\tan\alpha=\varphi'(x_0)=f(x_0,y_0)$};
\end{tikzpicture}
\end{center}

\begin{remarkHSE}[Мощное свойство]
Дифференциальное уравнение \eqref{eq:first-order} позволяет узнать наклон касательной к графику решения, даже \textbf{когда само решение \(\varphi(x)\) ещё не найдено}.
\end{remarkHSE}

В качестве примера рассмотрим поле направлений \(y'=y-x^2\):
\begin{center}
\begin{tikzpicture}[
    x=0.82cm, y=0.82cm,
    >=Stealth,
    field/.style={gray!65, line width=0.45pt, line cap=round},
    horizontal/.style={black!75, line width=0.8pt, line cap=round},
    solution/.style={blue!65!black, line width=0.95pt},
    isocline/.style={line width=1.05pt, dashed}
]
    % ===== Left: direction field without isoclines =====
    \begin{scope}
        \draw[step=0.5, gray!15, very thin] (-3,-2) grid (3,5);

        \begin{scope}
            \clip (-3.15,-2.15) rectangle (3.15,5.15);
            \foreach \x in {-3,-2.5,...,3} {
                \foreach \y in {-2,-1.5,...,5} {
                    \pgfmathsetmacro{\m}{\y-\x*\x}
                    \pgfmathsetmacro{\dx}{0.16/sqrt(1+\m*\m)}
                    \pgfmathsetmacro{\dy}{\m*\dx}
                    \draw[field]
                        ({\x-\dx},{\y-\dy}) -- ({\x+\dx},{\y+\dy});
                }
            }

            % The zero isocline y=x^2.
            \draw[isocline, orange!85!black]
                plot[domain=-2.3:2.3, samples=121, smooth]
                (\x,{\x*\x});

            % Explicit horizontal direction marks on y=x^2.
            \foreach \x in {-2,-1.75,...,2} {
                \pgfmathsetmacro{\yy}{\x*\x}
                \draw[horizontal]
                    ({\x-0.15},\yy) -- ({\x+0.15},\yy);
            }
        \end{scope}

        \draw[->, semithick] (-3.35,0) -- (3.45,0) node[right] {$x$};
        \draw[->, semithick] (0,-2.25) -- (0,5.35);
        \node[anchor=west] at (0.12,5.13) {$y$};
        \foreach \t in {-3,-2,-1,1,2,3} {
            \draw (\t,0.055) -- (\t,-0.055)
                node[below, fill=white, inner sep=1pt, font=\scriptsize] {$\t$};
        }
        \foreach \t in {-2,-1,1,2,3,4,5} {
            \draw (0.055,\t) -- (-0.055,\t)
                node[left, fill=white, inner sep=1pt, font=\scriptsize] {$\t$};
        }
        \node[below left, fill=white, inner sep=1pt, font=\scriptsize]
            at (0,0) {$0$};
    \end{scope}

    % ===== Right: the previous view with solution curves and y=x^2 =====
    \begin{scope}[xshift=6.5cm]
        \draw[step=0.5, gray!15, very thin] (-3,-2) grid (3,5);

        \begin{scope}
            \clip (-3.15,-2.15) rectangle (3.15,5.15);
            \foreach \x in {-3,-2.5,...,3} {
                \foreach \y in {-2,-1.5,...,5} {
                    \pgfmathsetmacro{\m}{\y-\x*\x}
                    \pgfmathsetmacro{\dx}{0.16/sqrt(1+\m*\m)}
                    \pgfmathsetmacro{\dy}{\m*\dx}
                    \draw[field]
                        ({\x-\dx},{\y-\dy}) -- ({\x+\dx},{\y+\dy});
                }
            }

            % Exact solutions: y=x^2+2x+2+Ce^x.
            \foreach \C in {-4,-2,-1,0,1} {
                \draw[solution]
                    plot[domain=-3.18:3.18, samples=241, smooth]
                    (\x,{\x*\x+2*\x+2+\C*exp(\x)});
            }

            % Zero isocline: y'=0 along y=x^2.
            \draw[isocline, orange!85!black]
                plot[domain=-2.3:2.3, samples=121, smooth]
                (\x,{\x*\x});
        \end{scope}

        \draw[->, semithick] (-3.35,0) -- (3.45,0) node[right] {$x$};
        \draw[->, semithick] (0,-2.25) -- (0,5.35);
        \node[anchor=west] at (0.12,5.13) {$y$};
        \foreach \t in {-3,-2,-1,1,2,3} {
            \draw (\t,0.055) -- (\t,-0.055)
                node[below, fill=white, inner sep=1pt, font=\scriptsize] {$\t$};
        }
        \foreach \t in {-2,-1,1,2,3,4,5} {
            \draw (0.055,\t) -- (-0.055,\t)
                node[left, fill=white, inner sep=1pt, font=\scriptsize] {$\t$};
        }
        \node[below left, fill=white, inner sep=1pt, font=\scriptsize]
            at (0,0) {$0$};

    \end{scope}

    % One centered legend for both panels.
    \node at (4,-2.75) {%
        \tikz[baseline=-0.55ex]
            \draw[blue!65!black, line width=0.95pt] (0,0) -- (0.75,0);
        \quad $y=x^2+2x+2+Ce^x$
        \qquad
        \tikz[baseline=-0.55ex]
            \draw[orange!85!black, dashed, line width=1.05pt]
                (0,0) -- (0.75,0);
        \quad $y=x^2\ (y'=0)$%
    };
\end{tikzpicture}
\end{center}

В каждой точке области \(G\) правая часть уравнения \eqref{eq:first-order} определяет направление касательной к решению, проходящему через эту точку. Множество всех подобных отрезков касательных к решениям в области \(G\) называется \textit{полем направлений}.

\begin{definitionHSE}[Определение поля направлений]
\textit{Полем направлений} уравнения \eqref{eq:first-order} называется множество всех отрезков касательных с угловым коэффициентом \(k=f(x,y)\), проведённых через точки области \(G\).
\end{definitionHSE}

\begin{remarkHSE}
Для уравнений высших порядков классическое двумерное поле направлений (на плоскости \(xOy\)) не строится.

Для уравнения второго порядка, например $y''=f(x,y,y')$, по одним только \(x\) и \(y\) направление уже определить нельзя: нужно ещё знать \(y'\). В одной и той же точке \((x,y)\) разные решения могут иметь разные наклоны.
\end{remarkHSE}

Формула для \(f(x,y)\) может не иметь смысла в отдельных точках, хотя предел там может существовать:
\[
f(x,y)=\frac{x^2-y^2}{x-y}\;(x\neq y),\;\lim_{y\rightarrow x}\frac{x^2-y^2}{x-y}=2x
\]
В таким местах функцию доопределяют по непрерывности, из-за чего область \(G\) становится шире, а поле направлений полнее. Но могут существовать разрывы, которые не устраняются подобным образом:
\begin{gather*}
f(x,y)=\frac{1}{x-y}\;(x\neq y),\;\lim_{y\rightarrow x}\frac{1}{x-y}=\infty
\end{gather*}
Если касасательная вертикальна (\(\alpha=\pi/2\)), то тангенс бесконечен, и функция в этой точке \textbf{не определена}. А кривые с такими касательными хотелось бы включить в рассмотрение.

Мы можем перевернуть уравнение и мерить угол не от оси \(Ox\), а \(Oy\):
\[
\frac{\dd y}{\dd x}=f(x,y)\iff\frac{\dd x}{\dd y}=\frac{1}{f(x,y)}
\]
Рассматриваем уравнение в тех точках, где его правая часть определена. Объединяя обе области определения (исходного уравнения и перевёрнутого), получают область определения \(\tilde{G} = G \cup G'\) \textit{расширенного поля направлений}.

На этой области исходное и перевёрнутое уравнения совместно задают расширенное поле направлений.

\begin{definitionHSE}[Определение изоклины]
Изоклина --- это геометрическое место точек, в которых поле направлений имеет один и тот же наклон.
\end{definitionHSE}

Её также называют линией равного наклона касательных.

Для уравнения \(y' = f(x, y)\) изоклины определяются уравнениями
\[
f(x,y)=k=\text{const}.
\]

\begin{remarkHSE}[Для чего нужны изоклины]
Нанеся несколько изоклин и короткие отрезки соответствующего наклона, можно приближённо восстановить вид интегральных кривых.
\end{remarkHSE}

Геометрически все штрихи поля направлений на этой линии параллельны.

\begin{exampleHSE}
«Изоклина двух» --- это линия, на которой тангенс угла наклона касательной к решению равен двум: \(\tg\alpha = 2\).
\end{exampleHSE}

\begin{definitionHSE}[Определение изолированной особой точки]
Точка \(P=(x_0,y_0)\) -- изолированная особая, если:
\begin{itemize}
\item в ней \textbf{не} определено расширенное поле направлений;
\item \(\exists\delta>0\): во всех точках \(\Odp{\delta}{x_0,y_0}\) расширенное поле определено.
\end{itemize} 
\end{definitionHSE}

Например, для \(y' = y/x\) исходная запись не определена при \(x = 0\), перевёрнутая
\({\dd x}/{\dd y} = x/y\) не определена при \(y = 0\), но их объединение задаёт направление всюду,
кроме \((0, 0)\). Поэтому \((0, 0)\) – изолированная особая точка.

\begin{definitionHSE}[Определение интегральной кривой]
Кривая, имеющая в каждой точке направление касательной, определяемое расширенным полем направлений уравнения \eqref{eq:first-order}, называется \textit{интегральной кривой}.
\end{definitionHSE}

Решение имеет вид \(y=\varphi(x)\), а интегральная кривая — его график, геометрическое место точек. Поэтому вертикальная прямая \(x=0\) не может быть интегральной кривой такого уравнения: она не задаёт \(y\) как функцию от \(x\).

\begin{definitionHSE}[Определение интеграла уравнения]
Соотношение \(\Phi(x,y)=0\), задающее некоторую интегральную кривую, называется \textit{интегралом уравнения}.
\end{definitionHSE}

\begin{definitionHSE}[Общее решение уравнения]
Однопараметрическое семейство решений уравнения \eqref{eq:first-order}
\[
y=\varphi(x,C),\qquad C\in E\subset\mathbb{R},
\]
называется \textit{общим решением уравнения} \eqref{eq:first-order} на множестве $M$, если
\begin{itemize}
\item при каждом $C_0\in E$ функция \(y=\varphi(x,C_0)\) является решением уравнения \eqref{eq:first-order} на \(M\);
\item для любого решения уравнения \eqref{eq:first-order} на \(M\) найдётся \(\tilde{C}\in E\), задающее это решение.
\end{itemize}
\end{definitionHSE}

Первое условие нужно для того, чтобы удостовериться в том, что мы не придумали лишних решений, а второе --- что мы никакие решения не потеряли.

\begin{definitionHSE}[Общий интеграл]
Соотношение
\[
\Phi(x,y,C)=0,\qquad C\in E\subset\mathbb{R},
\]
называется \textit{общим интегралом} уравнения \eqref{eq:first-order} на области \(D\subseteq\tilde{G}\), если
\begin{itemize}
\item при каждом $C_0\in E$ уравнение \(\Phi(x,y,C_0)=0\) задаёт интеграл уравнения \eqref{eq:first-order}, интегральная кривая которого лежит в \(D\);
\item для любой интегральной кривой уравнения \eqref{eq:first-order}, лежащей в \(D\), найдётся $\tilde{C}\in E$ такое, что уравнение \(\Phi(x,y,\tilde{C})=0\) задаёт эту интегральную кривую.
\end{itemize}
\end{definitionHSE}

\begin{remarkHSE}[Спор о выборе множества]
\end{remarkHSE}

Главное отличие \textit{общего интеграла} от \textit{общего решения} заключается в том, что он задаётся неявным уравнением, которое не всегда удаётся разрешить относительно \(y\).

Пример: \(x^2+y^2=C\) --- семейство концетрических окружностей. Окружность не является графиком функции.

Теорема о неявной функции утверждает: если функция \(F\) непрерывно дифференцируема в окрестности точки \((x_0,y_0)\), причём \(F(x_0,y_0)=0\) и \(F'_y(x_0,y_0)\neq 0\), то в некоторой окрестности этой точки уравнение \(F(x,y)=0\) однозначно разрешимо относительно \(y\), то есть задаёт \(y\) как функцию от \(x\).

\begin{exampleHSE}[Уравнение $y'=-x/y$]
\[
\frac{\dd y}{\dd x}=-\frac{x}{y},\qquad
\frac{\dd x}{\dd y}=-\frac{y}{x}.
\]
Область определения правой части исходного уравнения:
\[G=\mathbb{R}^2\setminus\{y=0\}\]
Область определения правой части перевёрнутого уравнения:
\[G'=\mathbb{R}^2\setminus\{y=0\}\]
Область определения расширенного поля направлений:
\[
\widetilde G=G\cup G'=\R^2\setminus\{(0,0)\},
\]
где $(0,0)$ --- изолированная особая точка. Разделяя переменные,
\[
y\dd y=-x\dd x,
\]
получаем
\[
\frac{x^2+y^2}{2}=C,\quad\boxed{\ x^2+y^2=C_1\ }\quad
(C,C_1\in\R)
\]
то есть интегральные кривые --- концетрические окружности с центром в начале координат.

% TODO: рисунок поля направлений и окружностей (страница 8 исходника).
\end{exampleHSE}

\begin{exampleHSE}[Уравнение $y'=y/x$ (дикритический узел)]
\[
\frac{\dd y}{\dd x}=\frac{y}{x},\qquad x\neq0.
\]
Область определения расширенного поля направлений аналогична предыдущему примеру:
\[
\widetilde G=G\cup G'=\R^2\setminus\{(0,0)\}
\]
Нестационарные решения:
\begin{gather*}
\frac{\dd y}{y}=\frac{\dd x}{x}
\ \Longrightarrow\ 
\ln|y|=\ln|x|+\ln C\ (C>0)
\ \Longrightarrow\ |y|=C|x|\ \Longrightarrow\ y=\pm Cx
\end{gather*}
Объединяем со стационарным решением \(y=0\):
\[y=C_1 x,\ (C_1\in\R)\]
В расширенном поле направлений дополнительно учитывается вертикальная интегральная прямая $x=0$.

Интегральными кривыми являются всевозможные лучи, выходящие из проколотого начала координат.

% TODO: рисунок изоклин и интегральных прямых (страница 9 исходника).
\end{exampleHSE}

\section{Задача Коши. Теорема Коши--Липшица}

\begin{definitionHSE}[Задача Коши]
Для уравнения первого порядка
\[
y'=f(x,y)
\]
задачей Коши называется система
\begin{equation}
\begin{cases}
 y'=f(x,y),\\
 y(x_0)=y_0.
\end{cases}
\label{eq:cauchy}
\end{equation}
Условие $y(x_0)=y_0$ называется начальным условием.
\end{definitionHSE}

Геометрически это означает, что среди всех возможных интегральных кривых уравнения \eqref{eq:first-order} необходимо найти ту, которая проходит через заданную точку \(P_0=(x_0,y_0)\).

\begin{theoremHSE}[Коши--Липшица]
Достаточные условия существования и единственности решения задачи Кошу ОДУ 1-го порядка, разрешённого относительно производной.
\end{theoremHSE}

\begin{center}
\begin{tikzpicture}[
  >=Stealth,
  line cap=round,
  line join=round,
  every node/.style={font=\normalsize}
]

  \definecolor{solution}{HTML}{C84949}
  \definecolor{domain}{HTML}{444444}

  % Основные параметры.
  \def\xzero{3.4}
  \def\yzero{2.4}
  \def\a{1.4}
  \def\b{1.1}
  \def\h{0.85}

  \pgfmathsetmacro{\xleft}{\xzero-\a}
  \pgfmathsetmacro{\xright}{\xzero+\a}
  \pgfmathsetmacro{\ybottom}{\yzero-\b}
  \pgfmathsetmacro{\ytop}{\yzero+\b}

  \pgfmathsetmacro{\jhleft}{\xzero-\h}
  \pgfmathsetmacro{\jhright}{\xzero+\h}

  \coordinate (M) at (\xzero,\yzero);

% Область D: асимметричная, округлая,
% левая часть более волнистая, но без резких изломов.
\path[
  draw=black!65,
  fill=black!2,
  line width=0.9pt
]
  (1.00,0.80)
  .. controls (0.78,0.95) and (0.55,1.20) ..
  (0.62,1.55)
  .. controls (0.70,1.90) and (0.95,2.05) ..
  (0.90,2.35)
  .. controls (0.84,2.65) and (0.62,2.95) ..
  (0.82,3.25)
  .. controls (1.02,3.55) and (1.45,3.85) ..
  (1.78,4.05)
  .. controls (2.10,4.25) and (2.40,4.42) ..
  (2.85,4.55)
  .. controls (3.45,4.80) and (3.85,4.30) ..
  (4.35,4.20)
  .. controls (4.90,4.10) and (5.25,4.45) ..
  (5.75,4.25)
  .. controls (6.20,4.05) and (6.20,3.55) ..
  (6.35,3.15)
  .. controls (6.55,2.70) and (6.95,2.55) ..
  (6.70,2.10)
  .. controls (6.50,1.75) and (6.05,1.70) ..
  (5.85,1.30)
  .. controls (5.60,0.85) and (5.15,0.70) ..
  (4.55,0.65)
  .. controls (4.00,0.60) and (3.55,0.30) ..
  (3.00,0.35)
  .. controls (2.35,0.40) and (1.55,0.35) ..
  (1.00,0.80)
  -- cycle;

\node at (6.35,4.15) {$D$};

  % Оси координат.
  \draw[->, thick]
    (-0.45,0) -- (6.8,0)
    node[below] {$x$};

  \draw[->, thick]
    (0,-1.05) -- (0,5.15)
    node[left] {$y$};

  \node[below left, inner sep=3pt] at (0,0) {$O$};

  % Прямоугольник Pi.
  \draw[line width=0.9pt]
    (\xleft,\ybottom)
    rectangle
    (\xright,\ytop);

  \node[anchor=north east, inner sep=5pt]
    at (\xright,\ytop)
    {$\Pi$};

  % Проекции границ прямоугольника на оси.
  \draw[densely dashed, black!60]
    (0,\ytop)
      node[left, text=black] {$y_0+b$}
    -- (\xleft,\ytop);

  \draw[densely dashed, black!60]
    (0,\yzero)
      node[left, text=black] {$y_0$}
    -- (M);

  \draw[densely dashed, black!60]
    (0,\ybottom)
      node[left, text=black] {$y_0-b$}
    -- (\xleft,\ybottom);

  \draw[densely dashed, black!60]
    (\xleft,0) -- (\xleft,\ybottom);

  \draw[densely dashed, black!60]
    (\xright,0) -- (\xright,\ybottom);

  \draw[densely dashed, black!60]
    (M) -- (\xzero,0);

  % Подписи x_0-a, x_0, x_0+a.
  \node[below=4pt] at (\xleft,0)
    {$x_0-a$};

  \node[below=4pt] at (\xzero,0)
    {$x_0$};

  \node[below=4pt] at (\xright,0)
    {$x_0+a$};

  % Интервал J.
  \draw[line width=1.3pt]
    (\jhleft,-0.13) -- (\jhright,-0.13);

  \fill (\jhleft,-0.13) circle (1.5pt);
  \fill (\jhright,-0.13) circle (1.5pt);

  \node[below=18pt] at (\jhleft,0)
    {$x_0-h$};

  \node[below=18pt] at (\jhright,0)
    {$x_0+h$};

  \node[below=36pt] at (\xzero,0)
    {$J$};

  % Решение y = phi(x), проходящее через M_0.
  \draw[
    solution,
    line width=1.4pt,
    domain=\jhleft:\jhright,
    samples=101
  ]
    plot
    (
      \x,
      {
        \yzero
        +0.60*(\x-\xzero)
        -0.16*(\x-\xzero)^3
      }
    );

  % Точка M_0.
  \fill (M) circle (2pt);

  % Подпись точки — близко к M_0, но не на кривой.
  \node[
    anchor=south east,
    inner sep=4pt
  ] at (3.30,2.53)
    {$P_0(x_0,y_0)$};

  % Подпись решения.
  \node[
    solution,
    anchor=west
  ] at (4.18,2.87)
    {$y=\varphi(x)$};

\end{tikzpicture}
\end{center}
\begin{proof}[Условия]
\textbf{Первое условие}. $\boxed{\ f\in C(D)$, $(x_0,y_0)\in D,\ D\subset\mathbb{R}^2.\ }$

Обозначим через \(\Pi\) замкнутый прямоугольник, целиком лежащий в области \(G\):
\[
\Pi=\{(x,y):|x-x_0|\le a,\ |y-y_0|\le b\}\subset D.
\]
Положим
\[
M=\max|f(x,y)|,\quad (x,y)\in\Pi,
\]
предполагаем, что $M\neq0$, случай нулевого максимума разберём отдельно.

Поскольку
\[|\varphi'(x)|=|f(x,\varphi(x))|\le M,\ (x,\varphi(x))\in\Pi,\]
получаем
\begin{gather*}
\varphi(x)=y_0+\int_{x_0}^{x}f(t,\varphi(t))\dd t\\[1em]
|\varphi(x)-y_0|=\left|\int_{x_0}^{x}f(t,\varphi(t))\dd t\right|\le\int_{x_0}^{x}|f(t,\varphi(t))|\dd t\le\int_{x_0}^{x} M\dd t=M|x-x_0|,\\[1em]
|\varphi(x)-y_0|\le M|x-x_0|
\end{gather*}
График решения \underline{на прямоугольнике} находится \textbf{внутри} «бабочки» с наклонами \(\pm M\).

Ограничиваем движение по \(x\), чтобы точка \((x,\varphi(x))\) не оказалась за пределами верхней или нижней границы \(\Pi\).
\[|\varphi(x)-y_0|\le M|x-x_0|\le b\quad\Longrightarrow\quad |x-x_0|\le\frac{b}{M}.\]
Оба ограничения на \(x\) должны выполняться одновременно, поэтому берём минимум:
\[J=\{x:|x-x_0|\le h\},\quad h=\min\left(a,\frac{b}{M}\right).\]
Если \(M=0\), полагают \(h=a\).

\textbf{Второе условие}. Если $f$ удовлетворяет условию Липшица по $y$ в $\Pi\in D$, то есть существует $L\ge0$ такое, что
\[
|f(x,y_1)-f(x,y_2)|\le L|y_1-y_2|,\quad\forall(x,y_1),(x,y_2)\in\Pi.
\]
\end{proof}
Если выполнены оба условия, то задача Коши \eqref{eq:cauchy} имеет притом единственное решение, определённое на основном отрезке $J$.

Непрерывность функции \(f\) существенна для \textbf{существования} решения.

Условие Липшица существенно для \textbf{единственности} решения.

Грубо говоря, существенность означает, что если условие есть, то всё хорошо. А если его нет, то может быть как хорошо, так и плохо.

\begin{center}
\begin{tikzpicture}[
    >=Stealth,
    scale=1.15,
    axis/.style={->, thick},
    bound/.style={red!70!black, dashed, thick},
    solution/.style={blue!70!black, very thick}
]
 
% Основные точки и размеры
\coordinate (M0) at (0,3);
\coordinate (BL) at (-3,1);
\coordinate (TR) at (3,5);
 
% Область допустимых наклонов
\fill[red!8]
    (M0) -- (2.4,5) -- (2.4,1) -- cycle;
 
\fill[red!8]
    (M0) -- (-2.4,5) -- (-2.4,1) -- cycle;
 
% Оси
\draw[axis] (-4,0) -- (4.3,0) node[right] {$x$};
\draw[axis] (-3.7,-0.3) -- (-3.7,5.8) node[above] {$y$};
 
% Прямоугольник Pi
\draw[thick]
    (BL) rectangle (TR);
 
\node[above right] at (3,5) {$\Pi$};
 
% Уровни y_0-b, y_0, y_0+b
\draw[dashed, gray] (-3.7,1) -- (3,1);
\draw[dashed, gray] (-3.7,3) -- (0,3);
\draw[dashed, gray] (-3.7,5) -- (3,5);
 
\node[left] at (-3.7,1) {$y_0-b$};
\node[left] at (-3.7,3) {$y_0$};
\node[left] at (-3.7,5) {$y_0+b$};
 
% Границы по x
\draw[dashed, gray] (-3,0) -- (-3,1);
\draw[dashed, gray] (3,0) -- (3,1);
 
\node[below] at (-3,0) {$x_0-a$};
\node[below] at (3,0) {$x_0+a$};
\node[below] at (0,0) {$x_0$};
 
% Предельные прямые с наклонами +M и -M
\draw[bound] (M0) -- (2.4,5)
    node[pos=0.72, above, sloped] {$y-y_0=M(x-x_0)$};
 
\draw[bound] (M0) -- (2.4,1)
    node[pos=0.72, below, sloped] {$y-y_0=-M(x-x_0)$};
 
\draw[bound] (M0) -- (-2.4,5);
\draw[bound] (M0) -- (-2.4,1);
 
% Плавная кубическая интегральная кривая
% \varphi(x)=3+0.18x+0.04x^3
\draw[solution, domain=-2.1:2.1, samples=150]
    plot (\x,{3 + 0.18*\x + 0.04*\x^3});

\node[blue!70!black, above right] at (1.45,3.48)
    {$y=\varphi(x)$};

% Начальная точка
\fill (M0) circle (2pt);
\node[above left] at (M0) {$M_0(x_0,y_0)$};

% Касательная в начальной точке:
% \varphi'(0)=0.18
\draw[orange!80!black, thick]
    (-1.2,{3-1.2*0.18}) --
    (1.2,{3+1.2*0.18});

\node[orange!80!black, below right] at (0.85,3.15)
    {$\varphi'(x_0)=f(x_0,y_0)$};
\end{tikzpicture}
\end{center}

\section{6 замечаний к теореме Коши–Липшица}

\begin{remarkHSE}[Достаточные условия Липшица]
Если $\dfrac{\partial f}{\partial y}$ существует и \(\left|\dfrac{\partial f}{\partial y}(x,y)\right|\le L\) на $\Pi$, то \(f\) липшицева по \(y\) с той же константой.
\begin{proof}
По теореме Лагранжа, если $g\in C[a;b]\cap D(a;b)$, то
\[
  \exists\, c\in(a;b)\colon\quad g(b)-g(a)=g'(c)\,(b-a).
\]
Возьмём $g(y)=f(\tilde x,y)$, где $\tilde x$ фиксировано. Тогда
\[
  g(y_1)-g(y_2)=\frac{\partial f}{\partial y}(\tilde x,c)\cdot(y_1-y_2),
\]
где $c$ лежит между $y_1$ и $y_2$. Поскольку
\(\left|\dfrac{\partial f}{\partial y}(\tilde x,c)\right|\le L\), получаем
\[
  |g(y_1)-g(y_2)|\le L\,|y_1-y_2|\qquad\forall\, y_1,y_2\in[y_0-b;\,y_0+b],
\]
т.\,е.
\[
  |f(\tilde x,y_1)-f(\tilde x,y_2)|\le L\,|y_1-y_2|\qquad\forall\,(\tilde x,y_1),\,(\tilde x,y_2)\in\Pi.
\]
\end{proof}
\end{remarkHSE}

\begin{remarkHSE}[Непрерывной производной по \(y\) достаточно для условия Липшица.]
В частности, если $\dfrac{\partial f}{\partial y}\in C(\Pi)$, то на компактном прямоугольнике эта производная ограничена (по теореме Вейерштрасса), а значит условие Липшица выполнено согласно первому замечанию.
\end{remarkHSE}

\begin{exampleHSE}[Условие Липшица слабее дифференцируемости]
Функция \(f(x,y)=|y|\) не дифференцируема по $y$ при $y=0$, но
\begin{gather*}
\bigl||y_1|-|y_2|\bigr|\le1\cdot|y_1-y_2|,
\end{gather*}
то есть условие Липшица выполнено с $L=1$.

Функция липшицева всюду, но не дифференцируема в нуле.
\end{exampleHSE}

\begin{exampleHSE}[Условие Липшица существенно для единственности]
Для \(y'=3y^{2/3}\) частная производная по $y$ \textbf{не существует} при $y=0$. Имеются решения
\[
y=0,
\qquad
y=\left(x+C\right)^3,\;C\in\mathbb{R}.
\]
Через точки оси $Ox$ могут проходить разные решения, поэтому единственность нарушается.
\end{exampleHSE}

\begin{center}
\begin{tikzpicture}[
    >=Stealth,
    line cap=round,
    line join=round,
    axis/.style={->, thick},
    solution/.style={very thick},
    scale=1.15
]

% Оси
\draw[axis] (-4.4,0) -- (4.8,0)
    node[below right] {$x$};

\draw[axis] (0,-3.1) -- (0,3.2)
    node[above left] {$y$};

% Нулевое решение
\draw[
    blue!70!black,
    very thick
]
    (-4.2,0) -- (4.5,0);

\node[
    blue!70!black,
    below
] at (-4,-0.08)
    {$y=0$};

% Кубическое решение
\draw[
    red!75!black,
    very thick
]
plot[
    domain=-1.45:1.45,
    samples=150,
    variable=\x
]
    ({\x},{\x^3});

\node[
    red!75!black,
    left
] at (-1,-1)
    {$\displaystyle y=x^3$};

% Решение с задержкой
\draw[
    green!35!black,
    ultra thick,
    dash pattern=on 7pt off 3pt
]
    (-4.2,0) -- (0.8,0);

% После точки a
\draw[
    green!35!black,
    ultra thick
]
plot[
    domain=0.8:2.25,
    samples=150,
    variable=\x
]
    ({\x},{(\x-0.8)^3});

% Точка отделения
\fill[green!25!black] (0.8,0) circle (2.5pt);

\node[below] at (0.8,-0.12)
    {$a$};

% Начальная точка
\fill (0,0) circle (2.2pt);

\node[above left] at (0,0)
    {$M_0=(0,0)$};

% Подпись зелёного решения ниже оси Ox
\node[
    green!25!black,
    fill=white,
    inner sep=3pt,
    align=left
] at (2.8,-1.1)
{
    $\displaystyle
    y=
    \begin{cases}
        0, & x\le a,\\[2pt]
        (x-a)^3, & x>a
    \end{cases}$
};

% Уравнение
\node at (-2.7,2.35)
    {$y'=3y^{2/3}$};

\end{tikzpicture}
\end{center}

Через каждую точку \((x_0,0)\) проходит бесконечно много решений: они могут произвольное время совпадать с \(y=0\), а затем перейти в кубическую ветвь. Следовательно, единственность решения нарушается.

Проверим, почему Липшиц здесь не выполняется.
\begin{proof}
Пусть \(L\) существует. Берём \(y_2=0\):
\begin{gather*}
\bigl|3y_1^{2/3}-0\bigr|\le L|y_1|\\
3\le Ly_1^{1/3}\Longrightarrow y_1^{1/3}\ge\frac{3}{L}
\end{gather*}
Но \(y_1\) можно брать сколь угодно малым, например \(y_1=1/n\):
\[\left(\frac{1}{n}\right)^{1/3}\ge\frac{3}{L}\xrightarrow[n\rightarrow\infty]{}0\ge\frac{3}{L}>0,\]
противоречие, значит \(L\) не существует.
\end{proof}

\begin{exampleHSE}[Непрерывность существенна для существования]
Рассмотрим задачу
\[y'=\sgn x,\quad y(0)=14\]
Она имеет решения
\begin{gather*}
y=x+C_1,\;(x>0);\qquad y=x+C_2,\;(x<0)
\end{gather*}
Склеиваем по непрерывности и начальному условию:
\[y(x)=|x|+14\]
\(|x|\) не дифференцируема в нуле, а решение должно быть дифференцируемо всюду, следовательно \textbf{решения не существует} вообще.
\end{exampleHSE}

\begin{exampleHSE}[Непрерывность достаточна, но не необходима]
Рассмотрим уравнение, у которого правая часть разрывна в нуле (\(2x\sin(1/x)\rightarrow0\), \(\cos(1/x)\) не имеет предела):
\[
y'=
\begin{cases}
 2x\sin\dfrac1x-\cos\dfrac1x,&x\neq0,\\
 0,&x=0,
\end{cases}
\]
но решение существует:
\[
\varphi(x)=
\begin{cases}
 x^2\sin\dfrac1x,&x\neq0,\\
 0,&x=0.
\end{cases}
\]
\[\varphi'(x)=\lim_{x\rightarrow0}\frac{x^2\sin(1/x)-0}{x-0}=\lim_{x\rightarrow0}x\sin\left(\frac{1}{x}\right)=0\]
\end{exampleHSE}

\section{Теорема Пеано}

\begin{theoremHSE}[Пеано]
Пусть в задаче Коши \eqref{eq:cauchy} функция $f(x,y)$ непрерывна в области $D\subset\mathbb{R}^2$. Тогда через каждую точку $D$ проходит \textbf{хотя бы одно} решение. 

Для фиксированной точки $(x_0,y_0)$ можно указать отрезок $J$, на котором решение существует.
\end{theoremHSE}

\section{Эквивалентное интегральное уравнение}

Производная -- вещь неудобная, интеграл лучше. К тому же мы можем встроить начальное условие в само уравнение.

\begin{theoremHSE}[Эквивалентность задачи Коши и интегрального уравнения]
Пусть $f\in C(D)$. Задача Коши \eqref{eq:cauchy} эквивалентна интегральному уравнению
\begin{equation}
 y(x)=y_0+\int_{x_0}^{x} f\bigl(t,y(t)\bigr)\dd t.
 \label{eq:integral-cauchy}
\end{equation}
\begin{itemize}
\item если \(y=\varphi(x)\) --- решение задачи Коши, то \(\varphi(x)\) --- решение уравнения \eqref{eq:integral-cauchy}.
\item если \(y=\varphi(x)\) --- непрерывное решение интегрального уравнения \eqref{eq:integral-cauchy}, то \(\varphi(x)\) --- решение задачи Коши.
\end{itemize}
\end{theoremHSE}
Доказательство в три строки:
\[
\left.
\begin{aligned}
\int_{x_0}^{x}\varphi'(t)\dd t
    &= \int_{x_0}^{x}f(t,\varphi(t))\dd t,\\
\int_{x_0}^{x}\varphi'(t)\dd t
    &= \varphi(x)-\varphi(x_0),\\
\varphi(x_0) &= y_0
\end{aligned}\
\right|
\quad\Longrightarrow\quad
\varphi(x)=y_0+\int_{x_0}^{x}f(t,\varphi(t))\dd t
\]

\begin{proof}
\(\boxed{\implies}\)

Если $y=\varphi(x)$ (непрерывная на \(M\)) решает задачу Коши, то
\[
\varphi'(x)=f(x,\varphi(x)).
\]
Интегрируя от $x_0$ до $x$, получаем
\[
\varphi(x)-\varphi(x_0)=\int_{x_0}^{x}f(t,\varphi(t))\dd t,
\]
а так как $\varphi(x_0)=y_0$, получаем \eqref{eq:integral-cauchy}.

\(\boxed{\impliedby}\)

Обратно, если $\varphi(x)$ (непрерывная на \(M\)) удовлетворяет \eqref{eq:integral-cauchy}, то, по теореме о производной интеграла с переменным верхним
пределом, дифференцирование даёт
\[
\varphi'(x)=f(x,\varphi(x)),
\]
а при $x=x_0$ интеграл равен нулю, поэтому $\varphi(x_0)=y_0$.
\end{proof}

\section{Метод Пикара}

\begin{definitionHSE}[Последовательные приближения Пикара]
Начинают с $y_0(x)\equiv y_0$ и строят
\[
\begin{array}{l@{\;}l}
y_1(x)     &= y_0+\displaystyle\int_{x_0}^{x} f(t,y_0)\dd t,\\[12pt]
y_2(x)     &= y_0+\displaystyle\int_{x_0}^{x} f\bigl(t,y_1(t)\bigr)\dd t,\\
\vdots     &\\[-2pt]
y_{n+1}(x) &= y_0+\displaystyle\int_{x_0}^{x} f\bigl(t,y_n(t)\bigr)\dd t.
\end{array}
\]
\end{definitionHSE}
Последовательность \(\{y_n\}\) равномерно сходится на множестве \(M\) к функции \(\varphi\), являющейся решением задачи Коши:
\[
y_n(x)\underset{x\in M}{\overset{n\to\infty}{\rightrightarrows}}\varphi(x).
\]
\begin{exampleHSE}[Задача $y'=y$, $y(0)=3$]
\[
y_1=3+3x,
\]
\[
y_2=3+3x+\frac{3x^2}{2!},
\]
\[
y_3=3+3x+\frac{3x^2}{2!}+\frac{3x^3}{3!},
\]
и в общем виде
\[
y_n(x)=3\left(1+x+\frac{x^2}{2!}+\cdots+\frac{x^n}{n!}\right).
\]
При $n\to\infty$ получаем
\[
y(x)=3e^x.
\]
\end{exampleHSE}

\section[Продолжение решения]{Продолжение решения и предельное поведение}

\begin{definitionHSE}[Продолжение вправо]
Пусть решение $\varphi$ определено на промежутке $J$. Решение $\widetilde\varphi$ называется продолжением $\varphi$ вправо, если оно определено на большем промежутке $\widetilde J\supset J$ и
\[
\widetilde\varphi(x)=\varphi(x),\qquad x\in J.
\]
Аналогично определяется продолжение влево.
\end{definitionHSE}

\begin{definitionHSE}[Продолжаемое решение]
Решение называется продолжаемым вправо, если существует его продолжение на больший промежуток справа. Аналогично определяется продолжаемость влево.
\end{definitionHSE}

В ограниченной области решение можно продолжать, последовательно ставя новые задачи Коши в концах уже построенных отрезков, пока траектория остаётся внутри области определения и выполняются условия единственности.

\begin{theoremHSE}[Предельное поведение в правой полуплоскости]
Пусть
\[
D=\{(x,y):x_0\le x<+\infty,\ -\infty<y<+\infty\},
\]
$f\in C(D)$ и выполнены условия единственности. Тогда решение, начинающееся на прямой $x=x_0$, либо продолжается вправо неограниченно, либо существует $x_1>x_0$ такое, что
\[
\lim_{x\to x_1-0}|\varphi(x)|=+\infty.
\]
То есть при невозможности продолжения внутри такой неограниченной области возникает вертикальная асимптота.
\end{theoremHSE}

\section[Разделяющиеся переменные]{Уравнение с разделяющимися переменными}

Рассмотрим задачу
\begin{equation}
\begin{cases}
 y'=f(x)g(y),\\
 y(x_0)=y_0.
\end{cases}
\label{eq:separable}
\end{equation}

\begin{theoremHSE}[Существование и единственность]
Пусть
\[
f\in C(a,b),\qquad g\in C(c,d),\qquad g(y)\neq0\quad \forall y\in(c,d).
\]
Тогда через каждую точку прямоугольника
\[
\Pi=\{(x,y):a<x<b,\ c<y<d\}
\]
проходит единственное решение задачи Коши \eqref{eq:separable}.
\end{theoremHSE}

\begin{proof}
Пусть $y=\varphi(x)$ --- решение. Так как $g(\varphi(x))\neq0$,
\[
\frac{\dd\varphi}{g(\varphi)}=f(x)\dd x.
\]
Интегрируя от $x_0$ до $x$ и делая замену $u=\varphi(t)$, получаем
\[
\int_{y_0}^{\varphi(x)}\frac{\dd u}{g(u)}
=
\int_{x_0}^{x}f(t)\dd t.
\]
Пусть $G'(y)=1/g(y)$ и $F'(x)=f(x)$. Тогда
\[
G(\varphi(x))=F(x)+G(y_0)-F(x_0).
\]
Так как $1/g$ непрерывна и не меняет знак, $G$ строго монотонна, поэтому существует обратная функция:
\[
\varphi(x)=G^{-1}\bigl(F(x)+G(y_0)-F(x_0)\bigr).
\]
Эта формула одновременно даёт существование и единственность решения.
\end{proof}

\subsection{Частный случай: \texorpdfstring{$y'=f(x)$}{y'=f(x)}}
Если $f\in C(a,b)$, то
\[
y(x)=y_0+\int_{x_0}^{x}f(t)\dd t.
\]
Все решения отличаются вертикальным сдвигом.

Если $f(x)\to\infty$ при $x\to a_1$, то поведение определяется несобственным интегралом
\[
J_1=\int_{x_0}^{a_1-0}f(t)\dd t.
\]
Если $J_1$ сходится, решения подходят к прямой $x=a_1$ в конечных точках и единственность в расширенном поле может нарушаться. Если $J_1$ расходится, $x=a_1$ становится вертикальной асимптотой.

\subsection{Частный случай: \texorpdfstring{$y'=g(y)$}{y'=g(y)}}
При $g(y)\neq0$
\[
x(y)=x_0+\int_{y_0}^{y}\frac{\dd u}{g(u)}.
\]
Все решения отличаются горизонтальным сдвигом.

Если $g(c_1)=0$, то $y=c_1$ --- стационарное решение. Поведение соседних решений определяется интегралом
\[
J_2=\int_{y_0}^{c_1-0}\frac{\dd u}{g(u)}.
\]
При сходимости возможен контакт с прямой $y=c_1$ за конечное время; при расходимости $y=c_1$ выступает горизонтальной асимптотой.

\begin{exampleHSE}[Построение интегральных кривых]
Для
\[
y'=-\frac{x}{y}
\]
получаем изоклины $f(x,y)=k$ и общий интеграл
\[
x^2+y^2=C.
\]
Интегральные кривые --- окружности. Точка $(0,0)$ является изолированной особой точкой расширенного поля.

% TODO: таблица изоклин и рисунок со страницы 20 исходника.
\end{exampleHSE}

\section[Однородное уравнение]{Уравнение с однородной правой частью}

Рассмотрим
\begin{equation}
\begin{cases}
 y'=f\!\left(\dfrac{y}{x}\right),\quad x\neq0,\\
 y(x_0)=y_0.
\end{cases}
\label{eq:homogeneous-first}
\end{equation}

\begin{theoremHSE}[Существование и единственность]
Пусть $f$ непрерывна на $(a,b)$ и
\[
f(u)-u\neq0,\qquad u\in(a,b).
\]
Тогда через каждую точку
\[
G=\left\{(x,y):a<\frac{y}{x}<b,\ x\neq0\right\}
\]
проходит единственное решение задачи Коши \eqref{eq:homogeneous-first}.
\end{theoremHSE}

\begin{proof}
Делаем замену
\[
u=\frac{y}{x},\qquad y=ux,\qquad y'=u'x+u.
\]
Тогда
\[
u'x+u=f(u),
\]
то есть
\[
\frac{\dd u}{\dd x}=\frac{f(u)-u}{x}.
\]
Получили уравнение с разделяющимися переменными. Условия теоремы обеспечивают непрерывность множителей и отсутствие нулей $f(u)-u$, поэтому к нему применима теорема предыдущего билета.
\end{proof}

Разделяя переменные,
\[
\frac{\dd x}{x}=\frac{\dd u}{f(u)-u}.
\]
После интегрирования
\[
\ln|x|=\Phi(u)+C
=
\Phi\!\left(\frac{y}{x}\right)+C.
\]

\begin{corollaryHSE}[Подобие интегральных кривых]
Если $(x,y)$ лежит на одной интегральной кривой, то замена
\[
\widetilde x=kx,\qquad \widetilde y=ky
\]
не меняет отношение $y/x$, а константа в общем интеграле изменяется на $\ln|k|$. Поэтому все интегральные кривые подобны с центром подобия в начале координат.
\end{corollaryHSE}

\begin{definitionHSE}[Уравнение возможных касательных]
Уравнение
\[
f(u)-u=0
\]
называется уравнением возможных касательных. Если $u_1$ --- его корень, то
\[
y=u_1x
\]
является решением исходного однородного уравнения и задаёт инвариантный луч.
\end{definitionHSE}

\section[Нормальная система ОДУ]{Нормальная система. ЛОС и ЛНС}

\begin{definitionHSE}[Нормальная система ОДУ]
Нормальной системой порядка $n$ называется система
\[
\left\{
\begin{aligned}
\dot{x}^1 &= F^1(t, x^1, \dots, x^n),\\[3pt]
\dot{x}^2 &= F^2(t, x^1, \dots, x^n),\\[3pt]
&\;\;\vdots\\
\dot{x}^n &= F^n(t, x^1, \dots, x^n).
\end{aligned}
\right.
\]
\begin{equation}
\dot{x}^i=F^i(t,x^1,\ldots,x^n),\qquad i=1,\ldots,n,
\label{eq:normal-system}
\end{equation}
где \(F^i(t,x^1,\ldots,x^n)\) непрерывные в области \(G\subset R^{n+1}\) скалярные функции.

В векторной форме:
\[
\dot{\vect{x}}=\vect{F}(t,\vect{x}).
\]
\end{definitionHSE}

Неизвестные в системе --- функции \(x^1(t),\ldots, x^n(t)\), а \(t\) — их общий аргумент: вещественное число из некоторого интервала.


«Нормальная» здесь означает, что:
\[\text{кол-во уравнений} = \text{кол-во неизвестных функций } x^1(t),\ldots x^n(t)\]

Вектор \(\vect{x}(t) = (x^1(t),\ldots,x^n(t))\) удобно представлять как положение точки в \(\mathbb{R}^n\) в момент \(t\).

\begin{remarkHSE}[Промежуток с произвольными границами]
Через $\langle a; b\rangle$ будем обозначать любой из промежутков
$(a; b)$, $[a; b)$, $(a; b]$, $[a; b]$.
\end{remarkHSE}

\begin{definitionHSE}[Решение нормальной системы]
Решением системы \eqref{eq:normal-system} на промежутке $\langle a; b\rangle$ называется совокупность \(n\) функций \(\varphi_1(t),\ldots,\varphi_n(t)\) такие, что
\begin{itemize}
\item дифференцируемы на $\langle a; b\rangle$;
\item их график лежит в области определения $G$;
\item $\forall t\in \langle a; b\rangle$ выполняются все уравнения системы: \(\varphi^i(t)=F^i(t,\varphi^1(t),\ldots,\varphi^n(t))\), \(i=1,\ldots,n.\).
\end{itemize}
\end{definitionHSE}

\begin{definitionHSE}[Задача Коши для системы]
К системе добавляют начальное условие
\[\underbrace{(x^1(t_0),\ldots,x^n(t_0))}_{\textstyle\mathstrut\vect{x}(t_0)}=\underbrace{(x_0^1,\ldots,x_0^n)}_{\textstyle\mathstrut\vect{x}_0}.\]
то есть
\[
x^i(t_0)=x_0^i,\qquad i=1,\ldots,n.
\]
\end{definitionHSE}

\begin{theoremHSE}[Существование и единственность для нормальной системы]
Если в области $G\subset\R^{n+1}$ функции $F_i$ и все частные производные
\(\dfrac{\partial F_i}{\partial x_j}\) (\(i,j=1,\ldots,n\)) непрерывны, то через каждую точку $(t_0,\vect{x}_0)\in G$ проходит единственное локальное решение задачи Коши.
\end{theoremHSE}

\begin{definitionHSE}[Линейная система]
Линейная нормальная система имеет вид
\[
\left\{
\begin{aligned}
\dot{x}_1 &= a_{11}(t)x_1 + a_{12}(t)x_2 + \dots + a_{1n}(t)x_n + f_1(t),\\
\dot{x}_2 &= a_{21}(t)x_1 + a_{22}(t)x_2 + \dots + a_{2n}(t)x_n + f_2(t),\\
&\;\;\vdots\\
\dot{x}_n &= a_{n1}(t)x_1 + a_{n2}(t)x_2 + \dots + a_{nn}(t)x_n + f_n(t).
\end{aligned}
\right.
\]
\begin{equation}
\dot{\vect{x}}=A(t)\vect{x}+\vect{f}(t).
\label{eq:los}
\end{equation}
Если $\vect{f}(t)\equiv0$, система называется \textbf{линейной однородной системой} (ЛОС):
\[
\dot{\vect{x}}=A(t)\vect{x}.
\]
Если $\vect{f}\not\equiv0$, система называется \textbf{линейной неоднородной системой} (ЛНС).
\end{definitionHSE}

\begin{theoremHSE}[Линейная система с непрерывными коэффициентами]
Если все элементы $A(t)$ и координаты $\vect{f}(t)$ непрерывны на промежутке $(a,b)$, то для любого $t_0\in(a,b)$ и любого $\vect{x}_0\in\R^n$ задача
\[
\begin{cases}
\dot{\vect{x}}=A(t)\vect{x}+\vect{f}(t),\\
\vect{x}(t_0)=\vect{x}_0
\end{cases}
\]
имеет единственное решение, причём оно продолжается на весь промежуток $(a,b)$.
\end{theoremHSE}

\section[Линейная зависимость и вронскиан]{Линейная зависимость вектор-функций и вронскиан}

\begin{definitionHSE}[Линейная зависимость вектор-функций]
Вектор-функции $\vect{x}_1(t),\ldots,\vect{x}_n(t)$ называются линейно зависимыми на множестве $M$, если существуют константы
\[
(c_1,\ldots,c_n)\neq(0,\ldots,0)
\]
такие, что
\[
c_1\vect{x}_1(t)+\cdots+c_n\vect{x}_n(t)=\vect{0},
\qquad \forall t\in M.
\]
Если такого набора нет, система линейно независима.
\end{definitionHSE}

\begin{lemmaHSE}[Связь функций и их значений]
Если вектор-функции линейно зависимы на $M$, то при каждом фиксированном $t_0\in M$ векторы
\[
\vect{x}_1(t_0),\ldots,\vect{x}_n(t_0)
\]
линейно зависимы. Обратное в общем случае неверно.
\end{lemmaHSE}

\begin{definitionHSE}[Вронскиан вектор-функций]
Определителем Вронского (\textit{вронскианом}) $n$ вектор-функций называется
\[
W(t)=\det X(t),
\quad
X(t)=\bigl(\vect{x}_1(t)\ \cdots\ \vect{x}_n(t)\bigr).
\]
\end{definitionHSE}

\begin{corollaryHSE}[Линейная зависимость и вронскиан]
Если $\vect{x}_1,\ldots,\vect{x}_n$ линейно зависимы на $M$, то
\[
W(t)=0\qquad \forall t\in M.
\]
Следовательно, если существует $t_0$ такое, что $W(t_0)\neq0$, то вектор-функции линейно независимы.
\end{corollaryHSE}

\begin{lemmaHSE}[Вронскиан решений ЛОС]
Пусть $\vect{x}_1(t),\ldots,\vect{x}_n(t)$ --- решения ЛОС \eqref{eq:los}, где $A(t)$ непрерывна на $M$. Если \(\exists t_1\in M:W(t_1)=0\), то:
\begin{itemize}
\item Решения линейно зависимы на всём $M$.
\item \(W(t)=0\quad\forall t\in M.\)
\end{itemize}
\end{lemmaHSE}

\begin{proof}
Из $W(t_1)=0$ следует линейная зависимость векторов $\vect{x}_j(t_1)$. Поэтому существуют постоянные $c_j$, не все нулевые, такие что
\[
\sum_{j=1}^{n}c_j\vect{x}_j(t_1)=\vect{0}.
\]
Функция
\[
\vect{z}(t)=\sum_{j=1}^{n}c_j\vect{x}_j(t)
\]
является решением ЛОС, так как любая линейная комбинация решений с постоянными коэффициентами --- снова решение:
\[
\dot{\vect{z}}
= \sum_{j=1}^{n} c_j\,\dot{\vect{x}}_j
= \sum_{j=1}^{n} c_j A(t)\vect{x}_j
= A(t)\sum_{j=1}^{n} c_j\vect{x}_j
= A(t)\vect{z}
\]
и удовлетворяет начальному условию $\vect{z}(t_1)=0$ (проверили в начале). 

Ставим задачу Коши (\(A(t)\) непрерывна на \(M\) по условию):
\[
\dot{\vect{z}} = A(t)\vect{z}, \quad \vect{z}(t_1) = 0.
\]

Нулевое решение удовлетворяет. По единственности
\[
\vect{z}(t)\equiv0,
\]
откуда следует линейная зависимость решений на всём $M$.
\end{proof}

\section[ФСР и фундаментальная матрица]{ФСР, общее решение и фундаментальная матрица ЛОС}

\begin{definitionHSE}[Фундаментальная система решений]
Фундаментальной системой решений (ФСР) ЛОС порядка $n$ называется система из $n$ линейно независимых решений
\[
\vect{x}_1(t),\ldots,\vect{x}_n(t).
\]
\end{definitionHSE}

\begin{theoremHSE}[Общее решение ЛОС]
Пусть $A(t)$ непрерывна на промежутке $M$, а $\vect{x}_1,\ldots,\vect{x}_n$ --- ФСР системы \eqref{eq:los}.

Тогда общее решение имеет вид
\[
\vect{x}(t)=c_1\vect{x}_1(t)+\cdots+c_n\vect{x}_n(t),
\quad c_i\in\R.
\]
Иначе говоря, решения ЛОС образуют линейное пространство
\[
\operatorname{Ker}L=\bigl\{\vect{x}(t):\ L\vect{x}=\vect{0}\bigr\},\qquad L=\frac{\dd}{\dd t}-A(t),
\]
размерности $n$, и любая ФСР --- базис в нём.
\end{theoremHSE}

\begin{proof}
\textbf{Линейное пространство.} Оператор $L$ линеен, поэтому его ядро --- линейное пространство: любая линейная комбинация решений ЛОС снова является решением.

\textbf{ФСР порождает всё пространство.} Пусть $\vect{x}(t)$ --- произвольное решение и $t_0\in M$.

Так как \(W(t_0)\neq0\) (по лемме о вронскиане решений ЛОС: иначе ФСР была бы линейно зависима), векторы $\vect{x}_1(t_0),\ldots,\vect{x}_n(t_0)$ образуют базис в \(\mathbb{R}^n\).

Разложение вектора по базису единственно: найдутся единственные числа $c_1,\ldots,c_n$, для которых
\[
\vect{x}(t_0)=\sum_{j=1}^{n}c_j\vect{x}_j(t_0).
\]
Обе функции решают одну и ту же задачу Коши, поэтому по единственности совпадают на всём $M$.

\textbf{Базис и размерность.} ФСР линейно независима по определению и, как только что показано, порождает \(\operatorname{Ker}L\). Значит, это базис. ФСР существует: решения задач Коши \(\dot{\vect{x}}=A(t)\vect{x}\), \(\vect{x}_j(t_0)=\vect{e}_j\), \(j=1,\ldots,n\), определены на всём \(M\) (по теореме о линейной системе с непрерывными коэффициентами), и их вронскиан \(W(t_0)=\det E=1\neq0\). Поэтому \(\dim\operatorname{Ker}L=n\).
\end{proof}

\begin{definitionHSE}[Фундаментальная матрица]
Матрица
\[
X(t)=\bigl(\vect{x}_1(t)\ \cdots\ \vect{x}_n(t)\bigr),
\]
столбцы которой образуют ФСР, называется фундаментальной матрицей ЛОС.
\end{definitionHSE}
Другими словами, эта та матрица, столбцы которой -- это линейно независимые решения ЛОС.

\begin{theoremHSE}[Свойства фундаментальной матрицы]
Пусть $X(t)$ --- фундаментальная матрица ЛОС \eqref{eq:los} с непрерывной на промежутке $M$ матрицей $A(t)$. Тогда
\begin{enumerate}
\item \(\dot X(t)=A(t)X(t)\) \(\forall t\in M\);
\item \(\det X(t)\neq0\) \(\forall t\in M\).
\end{enumerate}
\end{theoremHSE}

\begin{proof}
\textbf{1.} \(j\)-й столбец матрицы \(\dot X(t)\) --- это \(\dot{\vect{x}}_j(t)\), а \(j\)-й столбец произведения \(A(t)X(t)\) --- это \(A(t)\vect{x}_j(t)\). Каждый столбец \(\vect{x}_j\) --- решение ЛОС: \(\dot{\vect{x}}_j=A(t)\vect{x}_j\). Матрицы, у которых совпадают все столбцы, равны, поэтому \(\dot X=AX\).

\textbf{2.} \(\det X(t)=W(t)\) --- вронскиан ФСР. Если бы \(W(t_1)=0\) хотя бы в одной точке \(t_1\in M\), то по лемме о вронскиане решений ЛОС решения были бы линейно зависимы на \(M\). Это противоречит определению ФСР.

То же видно из формулы Лиувилля--Остроградского \eqref{eq:liouville-system} (она доказана в следующем разделе):
\[
W(t)=W(t_0)\exp\left(\int_{t_0}^{t}\tr A(\tau)\dd\tau\right).
\]
Экспонента не обращается в ноль, поэтому \(W\) либо тождественно равен нулю на \(M\), либо не равен нулю нигде. Для ФСР второй случай.
\end{proof}

\begin{theoremHSE}[Переход между фундаментальными матрицами]
Если $X(t)$ --- фундаментальная матрица ЛОС, а $C$ --- постоянная невырожденная матрица, то
\[
Y(t)=X(t)C
\]
также является фундаментальной матрицей той же системы.

Обратно, любые две фундаментальные матрицы $X(t)$ и $Y(t)$ одной ЛОС связаны равенством
\[
Y(t)=X(t)C,
\quad \det C\neq0.
\]
\end{theoremHSE}

\section[Лиувилль-Остроградский для систем]{Диф-ие определителя.\\Формула Лиувилля--Остроградского}

\begin{lemmaHSE}[Дифференцирование определителя]
Пусть $X(t)=\bigl(x_{ij}(t)\bigr)$ --- матрица размера $n\times n$, все элементы которой $x_{ij}(t)$ дифференцируемы, и $D(t)=\det X(t)$. Тогда $D(t)$ дифференцируема и
\[
D'(t)=D_1(t)+\ldots+D_n(t),
\]
где $D_i(t)$ --- определитель матрицы, полученной из $X(t)$ заменой всех элементов $i$-й строки на их производные.
\end{lemmaHSE}
\begin{proof}
Обозначим элемент матрицы \(X(t)\), стоящий на пересечении \(i\)-й строки и \(j\)-го столбца, как \(x_{ij}(t)\). Формула Лейбница для определителя:
\[D(t) = \sum_{\sigma \in S_n} \operatorname{sgn}(\sigma) \cdot x_{1,\sigma(1)} \cdot x_{2,\sigma(2)} \cdots x_{n,\sigma(n)}=\sum_{\sigma \in S_n} \operatorname{sgn}(\sigma) \prod_{i=1}^{n} x_{i, \sigma(i)},\]
где \(S_n\) -- множество всех перестановок, а \(\sigma(i)\) -- индекс столбца для \(i\)-й строки.
\[D'(t)=\sum_{\sigma \in S_n} \operatorname{sgn}(\sigma)\frac{\dd}{\dd t}\left(\prod_{i=1}^{n} x_{i, \sigma(i)}\right)\]
Каждое слагаемое в этой сумме -- это произведение \(n\) функций, зависящих от \(t\). Чтобы взять производную от суммы, нужно взять производную от каждого слагаемого. А чтобы взять производную от произведения \(n\) функций, применяем правило произведения:
\[(f_1 \cdot f_2 \cdots f_n)' = f_1' \cdot f_2 \cdots f_n + f_1 \cdot f_2' \cdots f_n + \cdots + f_1 \cdot f_2 \cdots f_n'\]
\[D'(t) = \sum_{\sigma \in S_n} \operatorname{sgn}(\sigma) \left( \sum_{j=1}^n x'_{j,\sigma(j)} \prod_{\substack{i=1 \\ i \neq j}}^n x_{i,\sigma(i)} \right)\]
Меняем порядок суммирования:
\[D'(t) = \sum_{j=1}^{n} \left( \sum_{\sigma \in S_n} \operatorname{sgn}(\sigma) x'_{j,\sigma(j)} \prod_{\substack{i=1 \\ i \neq j}}^{n} x_{i,\sigma(i)} \right)= \sum_{j=1}^{n}D_j(t)=D_1(t)+\ldots+D_n(t)\]
\end{proof}
\begin{theoremHSE}[Формула Лиувилля--Остроградского]
Пусть $W(t)$ --- вронскиан $n$ решений ЛОС \eqref{eq:los}, тогда
\begin{equation}
W(t)=W(t_0)\exp\left(\int_{t_0}^{t}\tr A(\tau)\dd\tau\right),
\label{eq:liouville-system}
\end{equation}
где \(\displaystyle\tr A(t)=\sum_{i=1}^{n}a_{ii}(t)\).
\end{theoremHSE}

\begin{proof}
По формуле дифференцирования определителя
\[
W'(t)=D_1+\cdots+D_n.
\]
Cтолбцы фундаментальной матрицы являются решениями \(\dot{\vect{x}}=A(t)\vect{x}\).

\(j\)-й столбец -- это вектор \(\vect{x}_j=\begin{pmatrix}x_{1j}\\x_{2j}\\\vdots\\x_{nj}\end{pmatrix}\)

Уравнение \(\dot{\vect{x}}_j=A\vect{x}_j\) для \(i\)-й компоненты (то есть для \(i\)-й строки) выглядит так:
\[\dot{x}_{ij}(t)=\sum_{k=1}^na_{ik}(t)x_{kj}(t) = a_{i1}x_{1j} + a_{i2}x_{2j} + \ldots + a_{in}x_{nj}\]
Это означает, что вся \(i\)-я строка в матрице \(D_i\) представляет собой линейную комбинацию всех строк исходной матрицы \(X(t)\):
\begin{gather*}
D_i=\det\begin{pmatrix}\ldots\\a_{i1}R_1+a_{i2}R_2+\ldots+a_{in}R_n\\\ldots\end{pmatrix}=\\= a_{i1} \det \begin{pmatrix} \dots \\ R_1 \\ \dots \end{pmatrix} + \cdots + a_{ii} \det \begin{pmatrix} \dots \\ R_i \\ \dots \end{pmatrix} + \cdots + a_{in} \det \begin{pmatrix} \dots \\ R_n \\ \dots \end{pmatrix}=a_{ii}(t)\,W(t),
\end{gather*}
где \(R_k\) -- это k-я строка матрицы \(X\).

Поэтому
\begin{gather*}
W'(t)=\sum_{i=1}^{n}a_{ii}(t)W(t)=W(t)\sum_{i=1}^{n}a_{ii}(t)=\tr A(t)\,W(t),\\
W'(t)=\tr A(t)\,W(t).
\end{gather*}
Далее можно решить диффур или пойти через интегрирующий множитель:
\[
\Bigl(W(t)\,e^{-\int_{t_0}^{t}\operatorname{tr}A(s)\dd s}\Bigr)'
=\bigl(W'(t)-\operatorname{tr}A(t)\,W(t)\bigr)\,e^{-\int_{t_0}^{t}\operatorname{tr}A(s)\dd s}=0,
\]
поэтому выражение в скобках постоянно. При $t=t_0$ оно равно $W(t_0)$, что и даёт \eqref{eq:liouville-system}.
\end{proof}

\section[ЛНС и вариация постоянных]{Общее решение ЛНС и метод вариации постоянных}

\begin{theoremHSE}[Общее решение ЛОС]
\[\vect{x}_\text{общ. одн.}=X(t)\vect{c},\]
где \(X(t)\) --- фундаментальная матрица, а \(\vect{c}\) --- вектор-столбец из произвольных \underline{вещественных} констант.
\end{theoremHSE}

\begin{theoremHSE}[Связь общего решения ЛНС и ЛОС]
Для ЛНС
\[
\dot{\vect{x}}=A(t)\vect{x}+\vect{f}(t)
\]
и соответствующей ЛОС
\[
\dot{\vect{x}}=A(t)\vect{x}
\]
общее решение неоднородной системы имеет вид
\[
\vect{x}_{\text{общ.\,неодн.}}
=
\vect{x}_{\text{част.\,неодн.}}
+
\vect{x}_{\text{общ.\,одн.}}.
\]
\end{theoremHSE}

\begin{proof}
Пусть $\vect{x}_p$ --- фиксированное решение ЛНС, а $\vect{u}$ --- любое другое решение. Тогда
\[
\vect{v}=\vect{u}-\vect{x}_p
\]
удовлетворяет
\[
\dot{\vect{v}}
=\underbrace{A(t)\vect{u}+\vect{f}}_{\textstyle\mathstrut\dot{\vect{u}}}-\underbrace{\bigl(A(t)\vect{x}_p+\vect{f}\bigr)}_{\textstyle\mathstrut\dot{\vect{x}}_p}
=A(t)(\vect{u}-\vect{x}_p)=A(t)\vect{v},
\]
то есть является решением ЛОС. Обратная проверка проводится прямой подстановкой.
\end{proof}

\begin{theoremHSE}[Метод вариации постоянных]
Пусть $X(t)$ --- фундаментальная матрица ЛОС. Ищем решение ЛНС в виде
\[
\vect{x}=X\vect{c}(t).
\]
Метод вариации постоянных предлагает сделать безумную на первый взгляд вещь: предположить, что решение неоднородной системы имеет такой же вид, но только теперь \(\vect{c}\) — это не константа, а неизвестная функция времени \(\vect{c}(t)\).

После подстановки в исходное уравнение получаем
\[
\dot{X}\vect{c}+X\dot{\vect{c}}=A(X\vect{c})+\vect{f}.
\]
Так как каждый столбец \(\vect{x}_j\) матрицы \(X\) является решением однородного уравнения \(\dot{\vect{x}}_j=A\vect{x}_j\), то и для всей матрицы целиком выполняется матричное уравнение \(\dot{X}=AX\). Подставим:
\begin{gather*}
\cancel{(AX)\vect{c}}+X\dot{\vect{c}}=\cancel{A(X\vect{c})}+\vect{f},\quad X\dot{\vect{c}}(t)=\vect{f}(t),
\end{gather*}
Поскольку \(X(t)\) --- фундаментальная матрица, её определитель (вронскиан) никогда не равен нулю. То есть существует обратная матрица \(X^{-1}(t)\), тогда
\[
\dot{\vect{c}}(t)=X^{-1}(t)\vect{f}(t).
\]
Следовательно,
\[
\vect{c}(t)=\vect{c}_0+
\int_{t_0}^{t}X^{-1}(\tau)\vect{f}(\tau)\dd\tau,
\]
Подставляем найденный \(\vect{c}(t)\) обратно в наше самое первое предположение \(\vect{x}=X\vect{c}(t)\)
\begin{equation}
\vect{x}(t)=
X(t)\,\vect{c}_0+
X(t)\int_{t_0}^{t}X^{-1}(\tau)\vect{f}(\tau)\dd\tau.
\label{eq:variation-system}
\end{equation}
\end{theoremHSE}

\section[ЛУ \texorpdfstring{$n$}{n}-го порядка и система]{Переход от линейного уравнения к системе. Вронскиан}

\begin{definitionHSE}[Линейное уравнение \(n\)-го порядка]
Уравнение
\begin{equation}
y^{(n)}+a_{n-1}(t)y^{(n-1)}+\cdots+a_1(t)y'+a_0(t)y=f(t),
\label{eq:linear-n}
\end{equation}
где \(y=y(t)\), \(t\in M\), \(M=\langle a;b\rangle\subset\R\) --- промежуток, а коэффициенты \(a_{n-1}(t),\ldots,a_0(t)\) и правая часть \(f(t)\) непрерывны на множестве \(\tilde{M}\subset\R\), то есть \(a_{n-1},\ldots,a_0,f\in C(\tilde{M})\), называется \textit{линейным уравнением} (ЛУ) порядка \(n\).
\begin{itemize}
\item Если \(\exists t\in M:\ f(t)\neq0\), уравнение \eqref{eq:linear-n} называется \textbf{линейным неоднородным уравнением} (ЛНУ) порядка \(n\).
\item Если \(f(t)=0\ \forall t\in M\), уравнение называется \textbf{линейным однородным уравнением} (ЛОУ) порядка \(n\).
\end{itemize}
\end{definitionHSE}

Здесь \(\tilde{M}\) --- множество, на котором заданы коэффициенты и правая часть (оно может и не быть промежутком), а \(M\) --- промежуток, на котором рассматривается уравнение и его решения. Чтобы уравнение имело смысл при \(t\in M\), должно быть \(M\subset\tilde{M}\).

\begin{remarkHSE}[Почему \(M\) --- промежуток]
\begin{itemize}
\item Связные подмножества \(\R\) --- это в точности промежутки. Связность нужна, чтобы единственность решения задачи Коши действовала на всём \(M\). Без неё ломаются свойство 3 вронскиана и теорема о размерности пространства решений: контрпример --- в замечании в п.~\ref{subsec:lou-by-fsr}.
\item Промежуток невырожденный: \(-\infty\le a<b\le+\infty\). В одной точке производные \(y',\ldots,y^{(n)}\) не имеют смысла.
\item Областью (открытым связным множеством) \(M\) не называем: в \(\R\) область --- это только интервал \((a;b)\), а промежуток \(\langle a;b\rangle\) допускает и концы, где производные понимаются как односторонние. Для теории важна связность, а не открытость.
\item От \(\tilde{M}\) связности не требуется: например, у уравнения \(y''-\frac{2}{t}y'+\frac{2}{t^2}y=0\) коэффициенты заданы на \(\tilde{M}=\R\setminus\{0\}\), а решения рассматривают на \(M=(0;+\infty)\) или \(M=(-\infty;0)\).
\end{itemize}
\end{remarkHSE}

\begin{exampleHSE}
\(y''-2ty=t\e^{t}\) --- ЛНУ 2-го порядка, а \(y''-2ty=0\) --- соответствующее ему ЛОУ 2-го порядка.
\end{exampleHSE}

\begin{definitionHSE}[Решение ЛУ]
Решением ЛУ \eqref{eq:linear-n} на множестве \(M\) называется функция \(y=y(t)\in D^n(M)\) (то есть имеющая на \(M\) производные до порядка \(n\) включительно), которая обращает уравнение \eqref{eq:linear-n} в тождество на \(M\).
\end{definitionHSE}

\subsection{Переход от ЛУ к ЛС}

Делаем замену
\begin{equation}
x^1=y,\quad x^2=y',\quad\ldots,\quad x^n=y^{(n-1)}.
\label{eq:linear-n-subst}
\end{equation}
Здесь верхний индекс --- номер координаты, а не степень; нижним индексом далее нумеруются решения (\(y_j\), \(\vect{x}_j\)).

Так как \(\dot x^n=y^{(n)}\), уравнение \eqref{eq:linear-n} записывается как
\[
\dot x^n+a_{n-1}(t)x^n+\cdots+a_1(t)x^2+a_0(t)x^1=f(t),
\]
а вместе с самой заменой получаем систему
\begin{equation}
\begin{cases}
\dot x^1=x^2,\\
\dot x^2=x^3,\\
\quad\vdots\\
\dot x^{n-1}=x^n,\\
\dot x^n=-a_0(t)x^1-a_1(t)x^2-\cdots-a_{n-1}(t)x^n+f(t)
\end{cases}
\label{eq:linear-n-system}
\end{equation}
--- линейную систему (ЛС) \(n\)-го порядка. Она имеет вид
\[
\dot{\vect{x}}=A(t)\vect{x}+\vect{F}(t),
\]
где
\begin{equation}
A(t)=
\begin{pmatrix}
0&1&0&\cdots&0\\
0&0&1&\cdots&0\\
\vdots&\vdots&\vdots&\ddots&\vdots\\
0&0&0&\cdots&1\\
-a_0(t)&-a_1(t)&-a_2(t)&\cdots&-a_{n-1}(t)
\end{pmatrix},
\qquad
\vect{F}(t)=
\begin{pmatrix}0\\\vdots\\0\\f(t)\end{pmatrix}.
\label{eq:linear-n-matrix}
\end{equation}
Все элементы \(A(t)\) и координаты \(\vect{F}(t)\) непрерывны на \(M\), то есть получилась ЛС с непрерывной матрицей \(A(t)\) и непрерывной правой частью \(\vect{F}(t)\), определёнными формулами \eqref{eq:linear-n-matrix}.

\begin{exampleHSE}[\(y''-2ty=t\e^{t}\)]
Замена \(x^1=y\), \(x^2=y'\):
\[
\begin{cases}
\dot x^1=x^2,\\
\dot x^2=2t\,x^1+t\e^{t},
\end{cases}
\qquad
\begin{pmatrix}\dot x^1\\\dot x^2\end{pmatrix}
=
\begin{pmatrix}0&1\\2t&0\end{pmatrix}
\begin{pmatrix}x^1\\x^2\end{pmatrix}
+
\begin{pmatrix}0\\t\e^{t}\end{pmatrix}.
\]
Итак,
\[
A(t)=\begin{pmatrix}0&1\\2t&0\end{pmatrix},\qquad
\vect{F}(t)=\begin{pmatrix}0\\t\e^{t}\end{pmatrix}.
\]
\end{exampleHSE}

Каждой функции \(y_j(t)\) замена \eqref{eq:linear-n-subst} ставит в соответствие вектор-функцию
\[
\vect{x}_j(t)=\begin{pmatrix}y_j(t)\\y_j'(t)\\\vdots\\y_j^{(n-1)}(t)\end{pmatrix}.
\]

\begin{lemmaHSE}[О замене]
При замене \eqref{eq:linear-n-subst}:
\begin{enumerate}
\item решения ЛУ \eqref{eq:linear-n} переходят в решения ЛС \eqref{eq:linear-n-system}, и наоборот;
\item линейно зависимые на \(M\) функции \(y_1,\ldots,y_k\) класса \(D^{n-1}(M)\) переходят в линейно зависимые на \(M\) вектор-функции \(\vect{x}_1,\ldots,\vect{x}_k\), и наоборот (линейная зависимость сохраняется).
\end{enumerate}
\end{lemmaHSE}

\begin{proof}
\textbf{1.} Пусть \(y=y(t)\) --- решение ЛУ \eqref{eq:linear-n}. Тогда последнее равенство системы \eqref{eq:linear-n-system} верно: это само уравнение \eqref{eq:linear-n}, записанное в новых переменных. Остальные \(n-1\) уравнений обращаются в равенства благодаря замене: \(\dot x^k=\bigl(y^{(k-1)}\bigr)'=y^{(k)}=x^{k+1}\).

Наоборот, пусть \(\bigl(x^1(t),\ldots,x^n(t)\bigr)^{T}\) --- решение системы. Из первых \(n-1\) уравнений
\[
x^2=\dot x^1,\quad x^3=\ddot x^1,\quad\ldots,\quad x^n=(x^1)^{(n-1)}.
\]
Подставляя это в последнее равенство системы, получаем уравнение \eqref{eq:linear-n} для функции \(y=x^1(t)\), то есть \(y=x^1(t)\) --- решение ЛУ.

\textbf{2.} Пусть \(y_1,\ldots,y_k\) линейно зависимы на \(M\), то есть \(\exists(c_1,\ldots,c_k)\neq\vect{0}\):
\[
c_1y_1(t)+\cdots+c_ky_k(t)=0\qquad\forall t\in M.
\]
Дифференцируем это тождество \(n-1\) раз:
\[
\begin{cases}
c_1y_1(t)+\cdots+c_ky_k(t)=0,\\
c_1y_1'(t)+\cdots+c_ky_k'(t)=0,\\
\quad\vdots\\
c_1y_1^{(n-1)}(t)+\cdots+c_ky_k^{(n-1)}(t)=0
\end{cases}
\qquad\forall t\in M,
\]
или в векторной форме
\[
c_1\begin{pmatrix}y_1(t)\\y_1'(t)\\\vdots\\y_1^{(n-1)}(t)\end{pmatrix}+\cdots+
c_k\begin{pmatrix}y_k(t)\\y_k'(t)\\\vdots\\y_k^{(n-1)}(t)\end{pmatrix}=\vect{0}
\qquad\forall t\in M,
\]
то есть
\[
c_1\vect{x}_1(t)+\cdots+c_k\vect{x}_k(t)=\vect{0}\qquad\forall t\in M,
\]
и вектор-функции \(\vect{x}_1,\ldots,\vect{x}_k\) линейно зависимы на \(M\).

Наоборот, пусть \(\vect{x}_1,\ldots,\vect{x}_k\) линейно зависимы на \(M\): \(\exists(c_1,\ldots,c_k)\neq\vect{0}\) такой, что \(c_1\vect{x}_1(t)+\cdots+c_k\vect{x}_k(t)=\vect{0}\) \(\forall t\in M\). Взяв первую строку этого векторного равенства, получаем
\[
c_1y_1(t)+\cdots+c_ky_k(t)=0\qquad\forall t\in M,
\]
то есть \(y_1,\ldots,y_k\) линейно зависимы на \(M\).
\end{proof}

\subsection{Общее решение ЛОУ}

Рассмотрим ЛОУ
\begin{equation}
y^{(n)}+a_{n-1}(t)y^{(n-1)}+\cdots+a_1(t)y'+a_0(t)y=0,\qquad a_k(t)\in C(M).
\label{eq:lou-n}
\end{equation}

\begin{definitionHSE}[ФСР ЛОУ]
Фундаментальной системой решений (ФСР) ЛОУ \eqref{eq:lou-n} \(n\)-го порядка называются любые \(n\) его решений, линейно независимых на \(M\).
\end{definitionHSE}

\begin{theoremHSE}[Об общем решении ЛОУ]
Решения ЛОУ \eqref{eq:lou-n} порядка \(n\) образуют линейное пространство размерности \(n\). Базис в нём --- любая ФСР \(y_1,\ldots,y_n\). Общее решение имеет вид
\[
y=c_1y_1(t)+\cdots+c_ny_n(t),\qquad c_i\in\R\ \text{--- произвольные постоянные}.
\]
\end{theoremHSE}

\begin{proof}
Замена \eqref{eq:linear-n-subst} линейна: \(y\mapsto\bigl(y,y',\ldots,y^{(n-1)}\bigr)^{T}\). По лемме о замене она переводит решения ЛОУ \eqref{eq:lou-n} в решения ЛОС \(\dot{\vect{x}}=A(t)\vect{x}\) с непрерывной матрицей \eqref{eq:linear-n-matrix} и обратно, а линейно независимые функции --- в линейно независимые вектор-функции и обратно.

Решения такой ЛОС образуют линейное пространство размерности \(n\): ФСР существует (например, решения задач Коши \(\vect{x}_j(t_0)=\vect{e}_j\), \(j=1,\ldots,n\)), и по теореме об общем решении ЛОС любая ФСР --- базис. Значит, то же верно и для решений ЛОУ.
\end{proof}

\subsection{Вронскиан и его свойства}

\begin{definitionHSE}[Вронскиан функций]
Вронскианом \(n\) функций \(y_1(t),\ldots,y_n(t)\in D^{n-1}(M)\) называется
\begin{equation}
W_{y_1,\ldots,y_n}(t)=
\begin{vmatrix}
y_1(t)&\cdots&y_n(t)\\
y_1'(t)&\cdots&y_n'(t)\\
\vdots&\ddots&\vdots\\
y_1^{(n-1)}(t)&\cdots&y_n^{(n-1)}(t)
\end{vmatrix}.
\label{eq:wronskian-n}
\end{equation}
\end{definitionHSE}

\begin{remarkHSE}[Связь с вронскианом вектор-функций]
При замене \eqref{eq:linear-n-subst} столбцы определителя \eqref{eq:wronskian-n} --- это в точности вектор-функции \(\vect{x}_1,\ldots,\vect{x}_n\). Поэтому вронскиан вектор-функций переходит во вронскиан \eqref{eq:wronskian-n}:
\[
W_{\vect{x}_1,\ldots,\vect{x}_n}(t)=W_{y_1,\ldots,y_n}(t).
\]
\end{remarkHSE}

\begin{theoremHSE}[Свойства вронскиана]
\begin{enumerate}
\item Если \(y_1,\ldots,y_n\) линейно зависимы на \(M\), то \(W_{y_1,\ldots,y_n}(t)=0\) \(\forall t\in M\).
\item Если \(\exists t_0\in M:\ W_{y_1,\ldots,y_n}(t_0)\neq0\), то \(y_1,\ldots,y_n\) линейно независимы на \(M\).
\item Если \(y_1,\ldots,y_n\) --- решения на \(M\) ЛОУ \eqref{eq:lou-n} \(n\)-го порядка с непрерывными коэффициентами и \(\exists t_0\in M:\ W_{y_1,\ldots,y_n}(t_0)=0\), то \(y_1,\ldots,y_n\) линейно зависимы на \(M\) (и, значит, \(W_{y_1,\ldots,y_n}(t)=0\) \(\forall t\in M\)).
\end{enumerate}
\end{theoremHSE}

\begin{proof}
\textbf{1.} По лемме о замене вектор-функции \(\vect{x}_1,\ldots,\vect{x}_n\) линейно зависимы на \(M\), поэтому их вронскиан равен нулю на \(M\) (следствие о линейной зависимости и вронскиане вектор-функций). Остаётся заметить, что \(W_{\vect{x}_1,\ldots,\vect{x}_n}=W_{y_1,\ldots,y_n}\).

\textbf{2.} Это контрапозиция пункта 1.

\textbf{3.} По лемме о замене \(\vect{x}_1,\ldots,\vect{x}_n\) --- решения ЛОС \(\dot{\vect{x}}=A(t)\vect{x}\) с непрерывной на \(M\) матрицей \(A(t)\), причём \(W_{\vect{x}_1,\ldots,\vect{x}_n}(t_0)=W_{y_1,\ldots,y_n}(t_0)=0\). По лемме о вронскиане решений ЛОС они линейно зависимы на \(M\), а тогда по лемме о замене линейно зависимы и \(y_1,\ldots,y_n\). Равенство \(W\equiv0\) следует из пункта 1.
\end{proof}

\begin{remarkHSE}[Решений меньше, чем порядок уравнения]
Пусть \(y_1,\ldots,y_k\) --- решения ЛОУ \eqref{eq:lou-n} \(n\)-го порядка с непрерывными коэффициентами, \(k<n\), и \(\exists t_0\in M:\ W_{y_1,\ldots,y_k}(t_0)=0\). Отсюда \textbf{не} следует, что \(y_1,\ldots,y_k\) линейно зависимы на \(M\): свойство 3 работает только для \(n\) решений, вронскиан которых совпадает с вронскианом решений ЛОС.
\end{remarkHSE}

\begin{exampleHSE}[\(y_1=t,\ y_2=t^2,\ M=\R\)]
Это решения, например, уравнения \(y'''=0\) (\(k=2<n=3\)).
\[
W_{t,t^2}(t)=
\begin{vmatrix}t&t^2\\1&2t\end{vmatrix}=t^2.
\]
\(W(0)=0\), но \(t\) и \(t^2\) линейно независимы на \(\R\), так как \(W(2)=4\neq0\).
\end{exampleHSE}

\begin{theoremHSE}[Формула Лиувилля--Остроградского для ЛОУ]
Для \(n\) решений ЛОУ \eqref{eq:lou-n}
\[
W(t)=W(t_0)\exp\left(-\int_{t_0}^{t}a_{n-1}(\tau)\dd\tau\right).
\]
\end{theoremHSE}

\begin{proof}
При переходе к системе матрица \eqref{eq:linear-n-matrix} имеет след
\[
\tr A(t)=-a_{n-1}(t),
\]
а вронскиан решений ЛОУ совпадает с вронскианом соответствующих решений ЛОС. Подставляя это в формулу Лиувилля--Остроградского для систем \eqref{eq:liouville-system}, получаем требуемое равенство.
\end{proof}

\begin{remarkHSE}[Случай \(n=2\)]
Для уравнения \(y''+a_1(t)y'+a_0(t)y=0\)
\[
\begin{vmatrix}y_1&y_2\\y_1'&y_2'\end{vmatrix}=y_1y_2'-y_1'y_2=C\e^{-\int a_1(t)\dd t}.
\]
Если одно решение известно (например, \(y_1\)), то это равенство --- линейное уравнение 1-го порядка относительно второго решения \(y_2\). Решив его при \(C\neq0\), получим решение, линейно независимое с \(y_1\).
\end{remarkHSE}

\subsection{Построение ЛОУ по ФСР}
\label{subsec:lou-by-fsr}

\begin{theoremHSE}[Построение ЛОУ по ФСР]
Пусть \(y_1,\ldots,y_n\in C^n(M)\) и \(W_{y_1,\ldots,y_n}(t)\neq0\) \(\forall t\in M\). Тогда уравнение
\begin{equation}
\begin{vmatrix}
y_1&\cdots&y_n&y\\
y_1'&\cdots&y_n'&y'\\
\vdots&&\vdots&\vdots\\
y_1^{(n)}&\cdots&y_n^{(n)}&y^{(n)}
\end{vmatrix}=0
\label{eq:lou-by-fsr}
\end{equation}
является ЛОУ \(n\)-го порядка вида \eqref{eq:lou-n} с непрерывными коэффициентами, для которого \(y_1,\ldots,y_n\) --- ФСР.
\end{theoremHSE}

\begin{proof}
Разложим определитель по последнему столбцу:
\[
W_{y_1,\ldots,y_n}(t)\,y^{(n)}+b_{n-1}(t)\,y^{(n-1)}+\cdots+b_0(t)\,y=0,
\]
где \(b_k(t)\) --- алгебраические дополнения элементов последнего столбца; они составлены из \(y_j^{(i)}\), \(i\leq n\).
\begin{enumerate}
\item Определитель линеен по последнему столбцу, поэтому получилось линейное однородное уравнение.
\item Его порядок равен \(n\): старшая производная --- \(y^{(n)}\), и коэффициент при ней не равен нулю (п.~3).
\item Коэффициент при \(y^{(n)}\) равен \(W_{y_1,\ldots,y_n}(t)\neq0\) \(\forall t\in M\), поэтому на него можно поделить и привести уравнение к виду \eqref{eq:lou-n} с \(a_k=b_k/W\).
\item Каждая \(y_k\) --- решение \eqref{eq:lou-by-fsr}: при \(y=y_k\) в определителе два одинаковых столбца, и он равен нулю. Так как \(W\neq0\), функции \(y_1,\ldots,y_n\) линейно независимы, то есть образуют ФСР.
\item Все коэффициенты непрерывны: \(b_k\) --- многочлены от непрерывных \(y_j^{(i)}\) (так как \(y_j\in C^n(M)\)), а \(W\) непрерывен и не обращается в ноль (см. разложение по последнему столбцу).
\end{enumerate}
\end{proof}

\begin{exampleHSE}[ЛОУ с ФСР \(t,\ t^2\)]
\[
\begin{vmatrix}t&t^2&y\\1&2t&y'\\0&2&y''\end{vmatrix}=0
\ \Longrightarrow\
t^2y''-2ty'+2y=0
\ \Longrightarrow\
y''-\frac{2}{t}\,y'+\frac{2}{t^2}\,y=0.
\]
Коэффициенты разрывны в точке \(t=0\) --- ровно там, где \(W_{t,t^2}(t)=t^2\) обращается в ноль. Поэтому теорема применима на \(M=(0;+\infty)\) или \(M=(-\infty;0)\), а ЛОУ 2-го порядка с непрерывными на всём \(\R\) коэффициентами и решениями \(t,\ t^2\) не существует: иначе по свойству 3 из \(W(0)=0\) следовала бы их линейная зависимость.
\end{exampleHSE}

\begin{remarkHSE}[Что ломается на несвязном \(M\)]
Рассмотрим то же уравнение на несвязном множестве \(M=\R\setminus\{0\}\), где его коэффициенты непрерывны, и функции
\[
y_1=t,\qquad
y_2=\begin{cases}t^2,&t>0,\\ t,&t<0.\end{cases}
\]
Обе --- решения на \(M\): на каждой половине это решения ЛОУ. Но
\[
W_{y_1,y_2}(-1)=\begin{vmatrix}-1&-1\\1&1\end{vmatrix}=0,\qquad
W_{y_1,y_2}(1)=\begin{vmatrix}1&1\\1&2\end{vmatrix}=1\neq0,
\]
то есть \(y_1,y_2\) линейно независимы на \(M\), хотя вронскиан в точке \(t=-1\) равен нулю: свойство 3 нарушено. Пространство решений на \(M\) четырёхмерное (на каждой половине по две произвольные постоянные), а не двумерное.

Причина в том, что единственность решения задачи Коши действует только внутри одного промежутка: на разных кусках несвязного \(M\) решение можно выбирать независимо.
\end{remarkHSE}

\subsection{Понижение порядка ЛОУ при известном частном решении}

\begin{theoremHSE}[Понижение порядка ЛОУ]
Пусть известно частное решение \(y=y_1(t)\not\equiv0\) ЛОУ \eqref{eq:lou-n}. Замена \(y=y_1u\) приводит к линейному уравнению \(n\)-го порядка относительно \(u\), не содержащему самой функции \(u\). Поэтому замена \(u'=z\) понижает порядок на единицу.
\end{theoremHSE}

\begin{proof}
По формуле Лейбница
\[
(y_1u)^{(k)}=\sum_{m=0}^{k}\binom{k}{m}y_1^{(k-m)}u^{(m)}.
\]
Подставим \(y=y_1u\) в левую часть \eqref{eq:lou-n} и соберём слагаемые с \(m=0\), то есть содержащие саму \(u\) без производных:
\[
u\cdot\bigl(y_1^{(n)}+a_{n-1}(t)y_1^{(n-1)}+\cdots+a_0(t)y_1\bigr)=0,
\]
так как \(y_1\) --- решение. Остальные слагаемые содержат только \(u',\ldots,u^{(n)}\); коэффициент при \(u^{(n)}\) равен \(y_1\). Поэтому на промежутках, где \(y_1(t)\neq0\), после замены \(u'=z\) получаем ЛОУ \((n-1)\)-го порядка относительно \(z\).
\end{proof}

\section[Вариация постоянных для ЛНУ]{Метод вариации постоянных для линейного неоднородного уравнения}
\label{sec:variation}

\subsection{Структура общего решения ЛНУ}

Запишем ЛНУ \eqref{eq:linear-n} кратко: \(L(D)y=f\), где
\[
L(D)=D^n+a_{n-1}(t)D^{n-1}+\cdots+a_1(t)D+a_0(t),\quad D=\frac{\dd}{\dd t}.
\]
Причём \(L(D)\) --- линейный оператор: \(L(D)(y_1+y_2)=L(D)y_1+L(D)y_2\), \(L(D)(\alpha y)=\alpha L(D)y\).

\phantomsection\label{rem:LD-variable}
\begin{remarkHSE}[Внимание: \(L(D)\) с переменными коэффициентами]
Здесь коэффициенты \(a_i(t)\) переменные, и \(L(D)\) --- просто обозначение оператора, составленного из \(D\). Свойства многочленов от \(D\) с постоянными коэффициентами из п.~\ref{subsec:poly-D} к нему неприменимы: не работают перестановочность \(L(D)M(D)=M(D)L(D)\), формула смещения и подстановка \(D\mapsto\lambda\) (характеристический многочлен). Причина в том, что \(D\) не перестановочен с умножением на функцию:
\[
D\bigl(a(t)y\bigr)=a'(t)y+a(t)Dy,\quad\text{то есть}\quad D\,a(t)\neq a(t)\,D.
\]
Линейность при этом сохраняется, и для структуры общего решения и принципа суперпозиции её достаточно.
\end{remarkHSE}

Для системы аналогично: ЛНС записывается как \(L\vect{x}=\vect{f}\), где \({L=\frac{\dd}{\dd t}-A(t)}\) --- линейный оператор.

Как и для систем (см. теорему об общем решении ЛНС), из линейности \(L(D)\) следует
\[
y_{\text{о.н.}}=y_{\text{ч.н.}}+y_{\text{о.о.}},
\]
то есть общее решение ЛНУ --- это сумма какого-нибудь частного решения ЛНУ и общего решения соответствующего ЛОУ.

\begin{theoremHSE}[Принцип суперпозиции]
Пусть \(y_1\) --- решение уравнения \(L(D)y=f_1\), а \(y_2\) --- решение уравнения \(L(D)y=f_2\). Тогда \(y_1+y_2\) --- решение уравнения
\[
L(D)y=f_1+f_2.
\]
\end{theoremHSE}

\begin{proof}
Подставим \(y=y_1+y_2\) и воспользуемся линейностью \(L(D)\):
\[
L(D)(y_1+y_2)=L(D)y_1+L(D)y_2=f_1+f_2.
\]
\end{proof}

Суперпозиция позволяет искать частное решение по кускам, если правая часть --- сумма нескольких слагаемых.

\subsection{Метод вариации постоянных}

Пусть \(y_1,\ldots,y_n\) --- ФСР ЛОУ, соответствующего ЛНУ \eqref{eq:linear-n}, то есть общее решение ЛОУ
\[
y_{\text{о.о.}}=c_1y_1(t)+\cdots+c_ny_n(t).
\]
При переходе к системе \eqref{eq:linear-n-system} ФСР ЛОУ переходит в ФСР соответствующей ЛОС, и её фундаментальная матрица --- это матрица Вронского:
\[
X(t)=
\begin{pmatrix}
y_1&\cdots&y_n\\
y_1'&\cdots&y_n'\\
\vdots&&\vdots\\
y_1^{(n-1)}&\cdots&y_n^{(n-1)}
\end{pmatrix},
\qquad
\vect{c}(t)=
\begin{pmatrix}c_1(t)\\\vdots\\c_n(t)\end{pmatrix},
\qquad
\vect{F}(t)=
\begin{pmatrix}0\\\vdots\\0\\f(t)\end{pmatrix}.
\]
Метод вариации постоянных для системы \(\dot{\vect{x}}=A(t)\vect{x}+\vect{F}(t)\) даёт \(X(t)\dot{\vect{c}}(t)=\vect{F}(t)\) (где \(X\)(t) --- ФМ ЛОС \(\dot{\vect{x}}=A(t)\vect{x}\)). Первая строка вектора \(\vect{x}=X\vect{c}\) --- это \(y\), поэтому получаем следующий алгоритм.

\begin{theoremHSE}[Вариация постоянных для ЛНУ]
Решение ЛНУ ищем в виде
\[
y(t)=c_1(t)y_1(t)+\cdots+c_n(t)y_n(t).
\]
\begin{enumerate}
\item Находим \(\dot c_1(t),\ldots,\dot c_n(t)\) из системы
\[
\begin{cases}
y_1\dot c_1(t)+\cdots+y_n\dot c_n(t)=0,\\
y_1'\dot c_1(t)+\cdots+y_n'\dot c_n(t)=0,\\
\qquad\vdots\\
y_1^{(n-2)}\dot c_1(t)+\cdots+y_n^{(n-2)}\dot c_n(t)=0,\\
y_1^{(n-1)}\dot c_1(t)+\cdots+y_n^{(n-1)}\dot c_n(t)=f(t).
\end{cases}
\]
Её определитель --- вронскиан \(W_{y_1,\ldots,y_n}(t)\neq0\) на \(M\), поэтому решение единственно (например, по формулам Крамера).
\item Интегрированием находим \(c_1(t),\ldots,c_n(t)\), каждую с точностью до своей произвольной постоянной \(\tilde c_i\).
\item Итак, общее решение ЛНУ:
\[
y(t)=c_1(t)y_1(t)+\cdots+c_n(t)y_n(t).
\]
\end{enumerate}
\end{theoremHSE}

\begin{remarkHSE}
Здесь ЛНУ записано с коэффициентом \(1\) при \(y^{(n)}\). Если старший коэффициент равен \(a_n\neq0\), то в последней строке системы стоит \(f(t)/a_n\), а не \(f(t)\).
\end{remarkHSE}

\begin{exampleHSE}[№\,578: \(\ddot y+4y=\operatorname{tg}t\)]
\textbf{I. ЛОУ.} \(\ddot y+4y=0\), характеристическое уравнение \(\lambda^2+4=0\), \(\lambda_{1,2}=\pm2i\), кратности \(k_{1,2}=1\). Значит,
\[
\text{ФСР}=\{\cos2t,\ \sin2t\},\qquad
y_{\text{о.о.}}=c_1\cos2t+c_2\sin2t,\quad c_1,c_2\in\R,
\]
где \(c_1,c_2\) --- произвольные постоянные.

\textbf{II. ЛНУ, вариация постоянных.} Ищем решение в виде \(y=c_1(t)\cos2t+c_2(t)\sin2t\). Тогда
\[
\begin{cases}
\dot c_1(t)\cos2t+\dot c_2(t)\sin2t=0,\\
-2\dot c_1(t)\sin2t+2\dot c_2(t)\cos2t=\operatorname{tg}t,
\end{cases}
\qquad
\begin{pmatrix}\cos2t&\sin2t\\-2\sin2t&2\cos2t\end{pmatrix}
\begin{pmatrix}\dot c_1\\\dot c_2\end{pmatrix}
=\begin{pmatrix}0\\\operatorname{tg}t\end{pmatrix}.
\]
Определители:
\[
\Delta=\begin{vmatrix}\cos2t&\sin2t\\-2\sin2t&2\cos2t\end{vmatrix}=2=W_{\cos2t,\sin2t}(t),
\]
\[
\Delta_1=\begin{vmatrix}0&\sin2t\\\operatorname{tg}t&2\cos2t\end{vmatrix}=-\operatorname{tg}t\,\sin2t=-2\sin^2t,
\qquad
\Delta_2=\begin{vmatrix}\cos2t&0\\-2\sin2t&\operatorname{tg}t\end{vmatrix}=\cos2t\,\operatorname{tg}t.
\]
Отсюда
\[
\dot c_1(t)=\frac{\Delta_1}{\Delta}=-\sin^2t,\qquad
\dot c_2(t)=\frac{\Delta_2}{\Delta}=\frac12\cos2t\,\operatorname{tg}t.
\]
Интегрируем:
\[
c_1(t)=-\int\sin^2t\dd t=-\int\frac{1-\cos2t}{2}\dd t=-\frac t2+\frac{\sin2t}{4}+\tilde c_1,
\]
\[
c_2(t)=\frac12\int\frac{(2\cos^2t-1)\sin t}{\cos t}\dd t
=-\frac12\int\frac{2\cos^2t-1}{\cos t}\dd(\cos t)
=-\frac{\cos^2t}{2}+\frac12\ln|\cos t|+\tilde c_2.
\]
Общее решение:
\begin{gather*}
y_{\text{о.н.}}=c_1(t)\cos 2t+c_2(t)\sin 2t=\\
=\underbrace{\tilde c_1\cos2t+\tilde c_2\sin2t}_{y_{\text{о.о.}}}
-\frac t2\cos2t+\frac{\sin2t\cos2t}{4}-\frac{\cos^2t}{2}\sin2t+\frac12\ln|\cos t|\sin2t,\\
\tilde c_1,\tilde c_2\in\R-\text{произвольные постоянные}.
\end{gather*}
Ответ можно упростить: \(\dfrac{\sin2t}{4}\bigl(\cos2t-2\cos^2t\bigr)=-\dfrac{\sin2t}{4}\) пропорционально \(\sin2t\) и поглощается слагаемым \(\tilde c_2\sin2t\). Итак,
\[
y=\tilde c_1\cos2t+\tilde c_2\sin2t-\frac t2\cos2t+\frac12\sin2t\,\ln|\cos t|,
\]
на любом промежутке \(\bigl(-\frac\pi2+\pi k;\frac\pi2+\pi k\bigr)\), \(k\in\Z\), где \(\operatorname{tg}t\) непрерывен.
\end{exampleHSE}

\section[ЛУ с постоянными коэффициентами]{Линейные уравнения с постоянными коэффициентами. Общая теория}

Здесь неизвестная функция \(y=y(x)\). Рассмотрим уравнение
\begin{equation}
a_ny^{(n)}+a_{n-1}y^{(n-1)}+\cdots+a_1y'+a_0y=0,
\qquad a_n\neq0,\quad a_i=\mathrm{const}\in\R,
\label{eq:lou-const}
\end{equation}
--- это ЛОУ \(n\)-го порядка с постоянными коэффициентами.

Поделив на \(a_n\), получаем уравнение вида \eqref{eq:lou-n} с коэффициентами, непрерывными на всём \(\R\), так что вся общая теория ЛОУ применима на \(M=\R\).

\subsection{Оператор дифференцирования и характеристический многочлен}

Пусть \(D=\dfrac{\dd}{\dd x}\) --- оператор взятия производной:
\[
Dy=y',\qquad D^2y=y'',\qquad\ldots,\qquad D^ky=y^{(k)}.
\]
Тогда \eqref{eq:lou-const} переписывается в виде
\[
a_nD^ny+a_{n-1}D^{n-1}y+\cdots+a_1Dy+a_0y=0,
\qquad
\underbrace{\bigl(a_nD^n+a_{n-1}D^{n-1}+\cdots+a_1D+a_0\bigr)}_{L(D)}y=0,
\]
где
\[
L(D)=a_nD^n+a_{n-1}D^{n-1}+\cdots+a_1D+a_0
\]
--- многочлен от оператора \(D\) (свободный член \(a_0\) понимается как \(a_0I\), \(I=\mathrm{id}\) --- тождественный оператор).

\begin{definitionHSE}[Характеристический многочлен]
Многочлен
\[
L(\lambda)=a_n\lambda^n+\cdots+a_1\lambda+a_0
\]
называется \textit{характеристическим многочленом} ЛОУ \eqref{eq:lou-const}, а уравнение \(L(\lambda)=0\) --- его \textit{характеристическим уравнением} (\(\deg L(\lambda)=n\)).
\end{definitionHSE}

\subsection{Свойства многочленов от \texorpdfstring{$D$}{D}}
\label{subsec:poly-D}

Пусть \(L(D)\), \(M(D)\) --- многочлены степеней \(m\) и \(k\) с постоянными коэффициентами, \(y\), \(y_1(x)\), \(y_2(x)\) --- функции, \(\alpha=\mathrm{const}\). Тогда:

\begin{enumerate}
\item \(L(D)\) --- линейный оператор: для \(y_1,y_2\in D^m(M)\)
\[
L(D)(y_1+y_2)=L(D)y_1+L(D)y_2,\quad
L(D)(\alpha y_1)=\alpha L(D)y_1.
\]
\item Сумма и умножение на число: для \(y\in D^{\max(m,k)}(M)\)
\[
\bigl(L(D)+M(D)\bigr)y=L(D)y+M(D)y,\quad
\bigl(\alpha M(D)\bigr)y=\alpha M(D)y.
\]
\item Произведение: для \(y\in D^{m+k}(M)\)
\[
\bigl(M(D)L(D)\bigr)y=M(D)\bigl(L(D)y\bigr).
\]
\end{enumerate}

Свойства 1--3 проверяются непосредственно: \(D^k\) линеен и \(D^iD^j=D^{i+j}\). Именно из-за постоянства коэффициентов многочлены от \(D\) можно складывать, перемножать и раскладывать на множители как обычные многочлены; в частности, \(L(D)M(D)=M(D)L(D)\).

\begin{theoremHSE}[Свойство 4: экспонента как решение]
Пусть \(L(D)\) --- многочлен с постоянными коэффициентами. Функция \(y=\e^{\lambda x}\) --- решение ЛОУ \eqref{eq:lou-const} \(\iff\) \(\lambda\) --- корень характеристического уравнения.
\end{theoremHSE}

\begin{proof}
Применим \(L(D)\) к \(\e^{\lambda x}\). Так как \(D\e^{\lambda x}=\lambda\e^{\lambda x}\) и \(D^k\e^{\lambda x}=\lambda^k\e^{\lambda x}\),
\[
L(D)\e^{\lambda x}=\bigl(a_nD^n+\cdots+a_0\bigr)\e^{\lambda x}
=a_n\lambda^n\e^{\lambda x}+\cdots+a_1\lambda\e^{\lambda x}+a_0\e^{\lambda x}
=\underbrace{\bigl(a_n\lambda^n+\cdots+a_1\lambda+a_0\bigr)}_{L(\lambda)}\e^{\lambda x}.
\]
Итак, \(L(D)\e^{\lambda x}=L(\lambda)\e^{\lambda x}=0\iff L(\lambda)=0\), поскольку \(\e^{\lambda x}\neq0\).
\end{proof}

\begin{theoremHSE}[Свойство 5: формула смещения]
Пусть \(L(D)\) --- многочлен степени \(n\) с постоянными коэффициентами, \(f\in D^n(M)\), \(\lambda=\mathrm{const}\in\mathbb C\). Тогда
\begin{equation}
L(D)\bigl(\e^{\lambda x}f(x)\bigr)=\e^{\lambda x}L(D+\lambda)f(x).
\label{eq:shift}
\end{equation}
\end{theoremHSE}

\begin{exampleHSE}
Для \(y''+4y=L(D)y\): \(L(D)\) --- это просто формальная запись левой части уравнения.

\(L(D)=D^2+4I=D^2+4\). Пусть \(\lambda=3\):
\begin{gather*}
L(D)\bigl(\e^{3x}f(x)\bigr)=(D^2+4I)\bigl(\e^{3x}f(x)\bigr)=\frac{\dd^2}{\dd x^2}\bigl(\e^{3x}f(x)\bigr)+4\e^{3x}f(x)=\\
=\frac{\dd}{\dd x}\bigl(3\e^{3x}f(x)+\e^{3x}f'(x)\bigr)+4\e^{3x}f(x)=9\e^{3x}f+3\e^{3x}f'+3\e^{3x}f'+\e^{3x}f''+4\e^{3x}f=\\
=\e^{3x}\bigl(\underbrace{f''+6f'+9f}_{(D+3)^2f}+4f\bigr)=\e^{3x}\bigl((D+3)^2+4\bigr)f=\e^{3x}L(D+3)f(x).
\end{gather*}
\end{exampleHSE}

\begin{proof}
Индукция по \(n\).

\textbf{База} \(n=1\): \(L(D)=aD+b\), многочлен первой степени.
\[
L(D)\bigl(\e^{\lambda x}f(x)\bigr)=(aD+b)\bigl(\e^{\lambda x}f(x)\bigr)
=a\lambda\e^{\lambda x}f(x)+a\e^{\lambda x}f'(x)+b\e^{\lambda x}f(x)=
\]
\[
=\e^{\lambda x}\bigl(a\lambda f+af'+bf\bigr)
=\e^{\lambda x}\bigl(\underbrace{aD+a\lambda+b}_{a(D+\lambda)+b}\bigr)f
=\e^{\lambda x}L(D+\lambda)f.
\]

\textbf{Шаг.} Пусть формула верна для всех многочленов \(L(D)\) степени \(\le n-1\), и \(\deg L(D)=n\). Над \(\mathbb C\) многочлен раскладывается на множители:
\[
L(D)=L_1(D)\,L_2(D),\qquad \deg L_1=1,\quad \deg L_2=n-1.
\]
По свойству 3 и предположению индукции
\[
L(D)\bigl(\e^{\lambda x}f(x)\bigr)
=(L_1(D)L_2)\bigl(\e^{\lambda x}f(x)\bigr)
\underset{3}{=}L_1(D)\Bigl(L_2(D)\bigl(\e^{\lambda x}f(x)\bigr)\Bigr)
\]
\[
\underset{\text{пр.инд.}}{=}L_1(D)\Bigl(\e^{\lambda x}\underbrace{L_2(D+\lambda)f(x)}_{g(x)}\Bigr)
\underset{\text{пр.инд.}}{=}\e^{\lambda x}L_1(D+\lambda)\bigl(L_2(D+\lambda)f(x)\bigr).
\]
Тогда
\[
L(D)\bigl(\e^{\lambda x}f(x)\bigr)=\e^{\lambda x}L_1(D+\lambda)L_2(D+\lambda)f(x)=\e^{\lambda x}L(D+\lambda)f(x).
\]
\end{proof}

\begin{theoremHSE}[Свойство 6: действие на квазимногочлен]
Пусть \(L(D)\) --- многочлен степени \(n\) с постоянными коэффициентами, \(\deg P(x)=k\), и \(\lambda\) --- не корень \(L\), то есть \(L(\lambda)\neq0\). Тогда
\[
L(D)\bigl(\e^{\lambda x}P(x)\bigr)=\e^{\lambda x}\tilde P(x),\quad\deg\tilde P=k.
\]
\end{theoremHSE}
То есть после применения оператора степень нового многочлена не изменяется, если \(L(\lambda)\neq0\).

\begin{proof}
По формуле смещения \(L(D)\bigl(\e^{\lambda x}P(x)\bigr)=\e^{\lambda x}L(D+\lambda)P(x)\). Разложим \(L(D+\lambda)\) по степеням \(D\):
\[
L(D+\lambda)=b_nD^n+\cdots+b_1D+b_0\Rightarrow b_0=L(\lambda).
\]
Лямбда никуда не исчезает --- она «прячется» внутри коэффициентов.

Каждое \(D^j\), \(j\ge1\), понижает степень многочлена, поэтому старший член \(\tilde P=L(D+\lambda)P\) получается только из \(b_0P\). Если \(L(\lambda)\neq0\), то \(b_0\neq0\) \(\Rightarrow\) \(\deg\bigl(L(D+\lambda)P(x)\bigr)=\deg P(x)=k\).
\end{proof}

\section[Различные корни]{Решение ЛОУ в случае некратных корней характеристического многочлена}
\label{sec:roots-simple}

\begin{theoremHSE}[Об общем решении ЛОУ в случае некратных корней]
Рассмотрим ЛОУ \eqref{eq:lou-const} \(n\)-го порядка. Пусть \(\lambda_1,\ldots,\lambda_n\in\mathbb C\) --- \(n\) корней характеристического многочлена, все кратности \(1\). Тогда ФСР ЛОУ имеет вид
\begin{equation}
\bigl\{\e^{\lambda_1x},\ldots,\e^{\lambda_nx}\bigr\},
\label{eq:fsr-simple}
\end{equation}
а общее решение ЛОУ \eqref{eq:lou-const}
\begin{equation}
y=c_1\e^{\lambda_1x}+\cdots+c_n\e^{\lambda_nx},\qquad x\in\R.
\label{eq:gen-simple}
\end{equation}
\end{theoremHSE}

\begin{proof}
Докажем, что функции \eqref{eq:fsr-simple} образуют ФСР. Тогда по теореме об общем решении ЛОУ получим формулу \eqref{eq:gen-simple}.
\begin{enumerate}[start=0]
\item Функций ровно \(n\) --- столько, сколько нужно.
\item Каждая \(y_k=\e^{\lambda_kx}\) --- решение по свойству 4.
\item Они линейно независимы на \(\R\). Достаточно проверить, что \(\exists x_0\in\R:\ W_{y_1,\ldots,y_n}(x_0)\neq0\). Имеем
\[
W_{\e^{\lambda_1x},\ldots,\e^{\lambda_nx}}(x)=
\begin{vmatrix}
\e^{\lambda_1x}&\cdots&\e^{\lambda_nx}\\
\lambda_1\e^{\lambda_1x}&\cdots&\lambda_n\e^{\lambda_nx}\\
\vdots&&\vdots\\
\lambda_1^{n-1}\e^{\lambda_1x}&\cdots&\lambda_n^{n-1}\e^{\lambda_nx}
\end{vmatrix},
\]
и при \(x_0=0\) получаем определитель Вандермонда:
\[
W_{\e^{\lambda_1x},\ldots,\e^{\lambda_nx}}(0)=
\begin{vmatrix}
1&\cdots&1\\
\lambda_1&\cdots&\lambda_n\\
\vdots&&\vdots\\
\lambda_1^{n-1}&\cdots&\lambda_n^{n-1}
\end{vmatrix}
=\prod_{\substack{i>j\\ i,j=1,\ldots,n}}(\lambda_i-\lambda_j)\neq0,
\]
так как \(\lambda_i\neq\lambda_j\) при \(i\neq j\). Значит, \(y_1,\ldots,y_n\) линейно независимы на \(\R\).
\end{enumerate}
\end{proof}

\begin{remarkHSE}
Корни \(\lambda_k\) могут быть комплексными, тогда \eqref{eq:gen-simple} --- общее комплексное решение (\(c_k\in\mathbb C\)). Так как коэффициенты вещественны, комплексные корни идут парами \(\lambda=\alpha\pm i\beta\), и пару \(\e^{(\alpha\pm i\beta)x}\) в ФСР можно заменить вещественной парой \(\e^{\alpha x}\cos\beta x,\ \e^{\alpha x}\sin\beta x\) --- это их линейные комбинации \(\frac12(\cdot+\cdot)\) и \(\frac1{2i}(\cdot-\cdot)\). Например, в №\,578 корни \(\pm2i\) дают ФСР \(\{\cos2t,\sin2t\}\).
\end{remarkHSE}

\section[Кратные корни]{Решение ЛОУ в случае кратных корней характеристического многочлена}
\label{sec:roots-multiple}

\begin{theoremHSE}[Об общем решении ЛОУ в случае кратных корней]
Пусть \(\lambda_1,\ldots,\lambda_m\) --- корни характеристического уравнения ЛОУ \eqref{eq:lou-const} порядка \(n\) кратностей \(k_1,\ldots,k_m\) соответственно, \(\lambda_i\neq\lambda_j\) при \(i\neq j\) (\(k_1+\cdots+k_m=n\)). Тогда функции
\begin{equation}
\begin{aligned}
\lambda_1&:\quad y_1=\e^{\lambda_1x},\quad y_2=x\e^{\lambda_1x},\quad\ldots,\quad y_{k_1}=x^{k_1-1}\e^{\lambda_1x},\\
\lambda_2&:\quad y_{k_1+1}=\e^{\lambda_2x},\quad y_{k_1+2}=x\e^{\lambda_2x},\quad\ldots,\quad y_{k_1+k_2}=x^{k_2-1}\e^{\lambda_2x},\\
&\quad\ \vdots\\
\lambda_m&:\quad y_{k_1+\cdots+k_{m-1}+1}=\e^{\lambda_mx},\quad\ldots,\quad y_{k_1+\cdots+k_m}=y_n=x^{k_m-1}\e^{\lambda_mx}
\end{aligned}
\label{eq:fsr-multiple}
\end{equation}
(корню \(\lambda_j\) соответствуют \(k_j\) функций) образуют ФСР ЛОУ \eqref{eq:lou-const}; все они определены \(\forall x\in\R\). Общее решение имеет вид
\[
y=c_1y_1+\cdots+c_ny_n,\quad x\in\R,
\]
где \(c_1,\ldots,c_n\) --- произвольные постоянные.
\end{theoremHSE}

\begin{proof}
\begin{enumerate}[start=0]
\item Функций ровно \(k_1+\cdots+k_m=n\).

\item Докажем, что каждая функция \(y_l\) --- решение. Пусть \(\lambda\) --- корень многочлена \(L(\lambda)\) кратности \(k\). Покажем, что
\[
y(x)=x^r\e^{\lambda x},\quad r=0,1,\ldots,k-1,
\]
--- решения ЛОУ \eqref{eq:lou-const}, то есть \(L(D)y=0\). Раз \(\lambda\) --- корень кратности \(k\), то можно вынести соответствующий множитель из характеристического многочлена:
\[
L(D)=(D-\lambda)^kM(D),\quad M(\lambda)\neq0.
\]
По формуле смещения \eqref{eq:shift}
\begin{gather*}
L(D)\bigl(x^r\e^{\lambda x}\bigr)=\e^{\lambda x}L(D+\lambda)x^r
=\e^{\lambda x}\Bigl((D-\lambda)^kM(D)\Bigr)\Big|_{D\mapsto D+\lambda}x^r=\\
=\e^{\lambda x}\bigl(D^kM(D+\lambda)\bigr)x^r
\underset{\text{комм.}}{=}\e^{\lambda x}M(D+\lambda)\bigl(\underbrace{D^kx^r}_{=0}\bigr)=0,
\end{gather*}
так как \(k>r\): \(k\)-я производная многочлена степени \(r<k\) равна нулю.

\item Докажем, что \(y_1,\ldots,y_n\) линейно независимы на \(\R\). От противного: пусть найдётся набор постоянных \((c_1,\ldots,c_n)\neq\vect{0}\), для которого
\[
c_1y_1(x)+\cdots+c_ny_n(x)=0\quad\forall x\in\R,
\]
то есть
\[
c_1\e^{\lambda_1x}+c_2x\e^{\lambda_1x}+\cdots+c_nx^{k_m-1}\e^{\lambda_mx}=0.
\]
Сгруппируем слагаемые по экспонентам:
\begin{equation}
\e^{\lambda_1x}P_1(x)+\cdots+\e^{\lambda_mx}P_m(x)=0\quad\forall x\in\R,\qquad\text{где }\deg P_j(x)\le k_j-1.
\label{eq:li-quasi}
\end{equation}
Например, при \(\lambda_1=1\) кратности \(3\) и \(\lambda_2=7\) кратности \(2\) равенство \eqref{eq:li-quasi} имеет вид
\[
(c_1x^2+c_2x+c_3)\e^{x}+(c_4x+c_5)\e^{7x}=0\quad\forall x\in\R,
\]
и нужно доказать, что \(c_1=c_2=\cdots=c_5=0\).

Так как \((c_1,\ldots,c_n)\neq\vect{0}\), то хотя бы один из многочленов \(P_j\) ненулевой. Без ограничения общности это \(P_1\): \(P_1(x)\not\equiv0\).

\begin{remarkHSE}[Как дифференцирование действует на \(\e^{\lambda x}P(x)\)]
По формуле смещения для \(L(D)=D\)
\[
D\bigl(\e^{\lambda x}P(x)\bigr)=\e^{\lambda x}(D+\lambda)P(x)=\e^{\lambda x}\bigl(DP(x)+\lambda P(x)\bigr)=\e^{\lambda x}\tilde P(x),
\]
и при \(\lambda\neq0\) \(\deg\tilde P=\deg P\) (свойство 6). В частности, ненулевой многочлен остаётся ненулевым.
\end{remarkHSE}

Делим \eqref{eq:li-quasi} на \(\e^{\lambda_mx}\neq0\):
\begin{equation}
\e^{(\lambda_1-\lambda_m)x}P_1(x)+\cdots+\e^{(\lambda_{m-1}-\lambda_m)x}P_{m-1}(x)+P_m(x)=0.
\label{eq:li-quasi-div}
\end{equation}
Дифференцируем \eqref{eq:li-quasi-div} \(k_m\) раз. Слагаемое \(P_m\) исчезнет, так как \(\deg P_m\le k_m-1\). В остальных слагаемых \(\lambda_j-\lambda_m\neq0\), поэтому степени многочленов сохраняются:
\[
\e^{(\lambda_1-\lambda_m)x}\tilde P_1(x)+\cdots+\e^{(\lambda_{m-1}-\lambda_m)x}\tilde P_{m-1}(x)=0\quad\forall x\in\R,
\qquad\deg\tilde P_j=\deg P_j\le k_j-1.
\]
Делим на \(\e^{(\lambda_{m-1}-\lambda_m)x}\). Показатель в первом слагаемом станет \(\lambda_1-\lambda_m-\lambda_{m-1}+\lambda_m=\lambda_1-\lambda_{m-1}\):
\[
\e^{(\lambda_1-\lambda_{m-1})x}\tilde P_1(x)+\cdots+\tilde P_{m-1}(x)=0.
\]
Это равенство того же вида, но с \(m-1\) слагаемым. Дифференцируем его \(k_{m-1}\) раз, и так далее.

Продифференцировав в сумме \(k_m+k_{m-1}+\cdots+k_3\) раз и сокращая каждый раз на экспоненту, получим
\[
\e^{(\lambda_1-\lambda_2)x}\hat P_1(x)+\hat P_2(x)=0.
\]
Дифференцируем \(k_2\) раз:
\[
\e^{(\lambda_1-\lambda_2)x}\check P_1(x)=0\quad\Longrightarrow\quad\check P_1(x)=0\quad\forall x\in\R.
\]
Это противоречит тому, что \(P_1(x)\not\equiv0\): на каждом шаге степень многочлена при первой экспоненте сохранялась, \(\deg\check P_1=\deg P_1\). (При \(m=1\) (одном корне \(\lambda_1\) определённой кратности) противоречие сразу: \(\e^{\lambda_1x}P_1(x)=0\Rightarrow P_1\equiv0\).)

Значит, \(c_1=\cdots=c_n=0\), и функции \eqref{eq:fsr-multiple} линейно независимы на \(\R\).
\end{enumerate}
Итак, \eqref{eq:fsr-multiple} --- ФСР, и по теореме об общем решении ЛОУ получаем формулу общего решения.
\end{proof}

\begin{remarkHSE}[Сгруппированная запись и комплексные корни]
Сгруппировав слагаемые по корням, общее решение можно записать как
\[
y=P_{k_1-1}(x)\e^{\lambda_1x}+\cdots+P_{k_m-1}(x)\e^{\lambda_mx},
\]
где \(P_{k_j-1}\) --- произвольный многочлен степени не выше \(k_j-1\).

Если среди корней есть комплексные, то постоянные \(c_i\in\mathbb C\). Как и в случае некратных корней, паре корней \(\alpha\pm i\beta\) кратности \(k\) в вещественной ФСР соответствуют функции \(x^r\e^{\alpha x}\cos\beta x\), \(x^r\e^{\alpha x}\sin\beta x\), \(r=0,\ldots,k-1\).
\end{remarkHSE}

\section[Квазимногочлен]{ЛНУ с правой частью в виде квазимногочлена}
\label{sec:quasi}

Функцию вида \(\e^{\lambda x}P(x)\), где \(P(x)\) --- многочлен, \(\lambda=\mathrm{const}\in\mathbb C\), называют \textit{простейшим квазимногочленом}.

\begin{definitionHSE}[Квазимногочлен]
Конечная сумма функций вида \(\e^{\lambda x}P(x)\), где \(P(x)\) --- многочлен, \(\lambda=\mathrm{const}\in\mathbb C\), называется \textit{квазимногочленом}.

В вещественной записи это даёт комбинации функций вида
\[
P(x)\e^{ax},\quad
P(x)\e^{ax}\cos bx,\quad
P(x)\e^{ax}\sin bx.
\]
\end{definitionHSE}

\begin{exampleHSE}
\(\e^{2x}+\e^{12ix}x^2\) и \((x^2+1)\e^{x}+(x-1)\e^{2x}+(4x^3-1)\e^{7x}\) --- квазимногочлены.
\end{exampleHSE}

Общее решение ЛНУ --- это сумма частного решения ЛНУ и общего решения ЛОУ, а общее решение ЛОУ с постоянными коэффициентами мы уже умеем находить. Остаётся найти частное решение.

\begin{theoremHSE}[О частном решении ЛНУ с правой частью --- квазиодночленом]
Рассмотрим ЛНУ \(n\)-го порядка с постоянными коэффициентами
\begin{equation}
L(D)y=\e^{\mu x}P_N(x),
\label{eq:lnu-quasi}
\end{equation}
где \(P_N\) --- многочлен степени \(N\), \(\mu=\mathrm{const}\in\mathbb C\). Тогда уравнение \eqref{eq:lnu-quasi} имеет частное решение вида
\[
y=x^s\e^{\mu x}Q_N(x),
\]
где \(Q_N\) --- многочлен степени \(N\), а
\[
s=\begin{cases}
\text{кратность }\mu\text{ как корня }L(\lambda),&\text{если }\mu\text{ --- корень},\\
0,&\text{иначе}.
\end{cases}
\]
Здесь \(L(\lambda)\) --- характеристический многочлен. Удобно считать, что если \(\mu\) не является корнем \(L(\lambda)\), то его кратность равна нулю.
\end{theoremHSE}

\begin{proof}
Подставим \(y=x^s\e^{\mu x}Q_N(x)\) в ЛНУ \eqref{eq:lnu-quasi}. Запишем
\[
P_N(x)=a_Nx^N+P_{N-1}(x),\quad
Q_N(x)=b_Nx^N+Q_{N-1}(x),\qquad
\deg P_{N-1},\ \deg Q_{N-1}\le N-1,
\]
где \(a_N\neq0\) известен, а \(b_N\) и коэффициенты \(Q_{N-1}\) нужно найти.

Так как \(\mu\) --- корень \(L(\lambda)\) кратности \(s\) (при \(s=0\) это значит, что \(L(\mu)\neq0\)),
\[
L(D)=(D-\mu)^sM(D),\qquad M(\mu)\neq0.
\]
По формуле смещения \eqref{eq:shift}
\[
L(D)\bigl(\e^{\mu x}x^sQ_N(x)\bigr)=\e^{\mu x}P_N(x)
\iff
\e^{\mu x}L(D+\mu)\bigl(x^sQ_N(x)\bigr)=\e^{\mu x}P_N(x)\quad\Big|:\e^{\mu x}\neq0,
\]
а
\[
L(D+\mu)=D^sM(D+\mu)=M(D+\mu)D^s.
\]
Получаем уравнение на многочлен \(Q_N\):
\begin{equation}
M(D+\mu)D^s\bigl(b_Nx^{N+s}+x^sQ_{N-1}(x)\bigr)=a_Nx^N+P_{N-1}(x).
\label{eq:quasi-poly-eq}
\end{equation}
Левая часть --- многочлен степени не выше \(N\): \(D^s\) понижает степень \(N+s\) до \(N\), а \(M(D+\mu)=M(\mu)+c_1D+c_2D^2+\cdots\) степень не повышает.

Приравняем коэффициенты при \(x^N\). Слева \(x^N\) даёт только слагаемое \(b_Nx^{N+s}\):
\[
D^s\bigl(b_Nx^{N+s}\bigr)=b_N(N+s)(N+s-1)\cdots(N+1)\,x^N=b_N\frac{(N+s)!}{N!}\,x^N,
\]
а из \(M(D+\mu)\) на старший член действует только \(M(\mu)\): остальные слагаемые содержат \(D\) и понижают степень. Справа коэффициент равен \(a_N\). Значит,
\[
b_N\frac{(N+s)!}{N!}M(\mu)=a_N
\Longrightarrow
b_N=\frac{a_N}{M(\mu)}\cdot\frac{N!}{(N+s)!}.
\]
Делить можно, так как \(M(\mu)\neq0\). Причём \(b_N\neq0\) (т.к. \(a_N\neq0\)), так что \(\deg Q_N=N\).

Итак, нашли коэффициент \(b_N\) при \(x^N\). Подставив его в \eqref{eq:quasi-poly-eq}, получим
\[
M(D+\mu)D^s\bigl(x^sQ_{N-1}(x)\bigr)=a_Nx^N+P_{N-1}(x)-M(D+\mu)D^s\bigl(b_Nx^{N+s}\bigr)=:\hat P_{N-1}(x),
\]
где \(\deg\hat P_{N-1}\le N-1\): члены с \(x^N\) сократились. Это уравнение того же вида для многочлена \(Q_{N-1}\), то есть всё свелось к поиску старшего коэффициента \(Q_{N-1}\). Повторяя рассуждение (индукция по \(N\)), однозначно находим все коэффициенты \(Q_N\).
\end{proof}

\begin{remarkHSE}[Метод неопределённых коэффициентов]
На практике \(Q_N\) записывают с неизвестными коэффициентами, подставляют \(y=x^s\e^{\mu x}Q_N(x)\) в уравнение и приравнивают коэффициенты при одинаковых степенях \(x\). Если правая часть --- сумма нескольких простейших квазимногочленов, частное решение ищут для каждого слагаемого отдельно и складывают (принцип суперпозиции).
\end{remarkHSE}

\section[Интегральная формула Коши]{Интегральная формула Коши для частного решения ЛНУ}
\label{sec:cauchy}

В этом разделе независимая переменная обозначается \(x\). Рассмотрим ЛНУ
\begin{equation}
y^{(n)}+a_{n-1}(x)y^{(n-1)}+\ldots+a_1(x)y'+a_0(x)y=f(x),
\label{eq:lnu-cauchy}
\end{equation}
где \(a_i(x)\), \(f(x)\) непрерывны на \((\alpha;\beta)\) (допускается \(\alpha=-\infty\), \(\beta=+\infty\)). Пусть \(y_1(x),\ldots,y_n(x)\) --- ФСР ЛОУ, соответствующего ЛНУ \eqref{eq:lnu-cauchy}.

\begin{definitionHSE}[Ядро Коши]
Функция
\begin{equation}
K(x,s)=c_1(s)y_1(x)+\cdots+c_n(s)y_n(x),\qquad x,s\in(\alpha;\beta),
\label{eq:cauchy-kernel}
\end{equation}
коэффициенты которой выбраны так, что
\begin{equation}
\begin{cases}
K(s,s)=0,\\
K'_x(s,s)=0,\\
\quad\vdots\\
K^{(n-2)}_x(s,s)=0,\\
K^{(n-1)}_x(s,s)=1,
\end{cases}
\label{eq:cauchy-kernel-cond}
\end{equation}
называется \textit{ядром Коши} ЛНУ \eqref{eq:lnu-cauchy}. Здесь \(K^{(k)}_x(s,s)\) --- производная \(\dfrac{\partial^kK}{\partial x^k}\), вычисленная при \(x=s\).
\end{definitionHSE}

\begin{remarkHSE}[Корректность определения]
\begin{itemize}
\item Условия \eqref{eq:cauchy-kernel-cond} --- линейная система относительно \(c_1(s),\ldots,c_n(s)\), матрица которой --- матрица Вронского ФСР в точке \(s\). Её определитель \(W_{y_1,\ldots,y_n}(s)\neq0\), поэтому \(c_i(s)\) находятся однозначно (по формулам Крамера) и непрерывны по \(s\).
\item При каждом фиксированном \(s\in(\alpha;\beta)\) функция \(K(x,s)\) --- решение ЛОУ по переменной \(x\): это линейная комбинация решений \(y_1,\ldots,y_n\) с постоянными по \(x\) коэффициентами.
\item Ядро задано на квадрате \((\alpha;\beta)\times(\alpha;\beta)\), а условия \eqref{eq:cauchy-kernel-cond} ставятся на его диагонали \(x=s\).
\end{itemize}
\end{remarkHSE}

\begin{center}
\begin{tikzpicture}[scale=0.8, >=Stealth]
  \fill[HSElightblue!15] (1,1) rectangle (3.6,3.6);
  \draw[thick, HSEdarkgray] (1,1) rectangle (3.6,3.6);
  \draw[very thick, HSEblue] (1,1) -- (3.6,3.6);
  \node[HSEblue, right] at (3.6,3.6) {$x=s$};
  \draw[->] (-0.3,0) -- (4.6,0) node[below] {$s$};
  \draw[->] (0,-0.3) -- (0,4.6) node[left] {$x$};
  \draw[densely dashed, gray] (1,0) -- (1,1) (3.6,0) -- (3.6,1) (0,1) -- (1,1) (0,3.6) -- (1,3.6);
  \node[below] at (1,0) {$\alpha$};
  \node[below] at (3.6,0) {$\beta$};
  \node[left] at (0,1) {$\alpha$};
  \node[left] at (0,3.6) {$\beta$};
\end{tikzpicture}
\end{center}

\begin{theoremHSE}[Интегральная формула Коши]
Частное решение ЛНУ \eqref{eq:lnu-cauchy} имеет вид
\begin{equation}
Y(x)=\int_\gamma^x K(x,s)f(s)\dd s,
\label{eq:cauchy-formula}
\end{equation}
где \(\gamma\in(\alpha;\beta)\), а \(K(x,s)\) --- ядро Коши.
\end{theoremHSE}

\begin{proof}
Подставим \(Y(x)\) из \eqref{eq:cauchy-formula} в ЛНУ \eqref{eq:lnu-cauchy}. Интеграл с переменным верхним пределом дифференцируется по правилу
\[
\frac{\dd}{\dd x}\int_\gamma^x g(x,s)\dd s=\int_\gamma^x g'_x(x,s)\dd s+g(x,x)
\qquad(g,\ g'_x\text{ непрерывны}).
\]
Внеинтегральные слагаемые обнуляются условиями \eqref{eq:cauchy-kernel-cond} при \(s=x\):
\begin{gather*}
Y'(x)=\int_\gamma^x K'_x(x,s)f(s)\dd s+\underbrace{K(x,x)}_{=0}f(x)=\int_\gamma^x K'_x(x,s)f(s)\dd s,\\
Y''(x)=\int_\gamma^x K''_{xx}(x,s)f(s)\dd s+\underbrace{K'_x(x,x)}_{=0}f(x)=\int_\gamma^x K''_{xx}(x,s)f(s)\dd s,\\
\cdots\\
Y^{(n-1)}(x)=\int_\gamma^x K^{(n-1)}_x(x,s)f(s)\dd s,\\
Y^{(n)}(x)=\int_\gamma^x K^{(n)}_x(x,s)f(s)\dd s+\underbrace{K^{(n-1)}_x(x,x)}_{=1}f(x)=\int_\gamma^x K^{(n)}_x(x,s)f(s)\dd s+f(x).
\end{gather*}
Подставим в левую часть \eqref{eq:lnu-cauchy}. Коэффициенты \(a_i(x)\) не зависят от \(s\), поэтому их можно внести под интеграл (аргументы \((x,s)\) у \(K\) опущены):
\begin{gather*}
Y^{(n)}+a_{n-1}(x)Y^{(n-1)}+\cdots+a_1(x)Y'+a_0(x)Y=\\
=\int_\gamma^x\Bigl(K^{(n)}_x+a_{n-1}(x)K^{(n-1)}_x+\cdots+a_1(x)K'_x+a_0(x)K\Bigr)f(s)\dd s+f(x)=\\
=\int_\gamma^x\underbrace{\bigl(L(D)K(x,s)\bigr)}_{=0}f(s)\dd s+f(x)=f(x),
\end{gather*}
где \(L(D)=D^n+a_{n-1}(x)D^{n-1}+\ldots+a_1(x)D+a_0(x)\), \(D={\dd}/{\dd x}\). Здесь \(L(D)K(x,s)=0\), так как при каждом фиксированном \(s\) функция \(K(x,s)\) --- линейная комбинация решений однородного уравнения, значит, сама является его решением ЛОУ.
\end{proof}

\begin{remarkHSE}[\(L(D)\) с переменными коэффициентами]
Коэффициенты \(a_i(x)\) здесь переменные, поэтому к \(L(D)\) неприменимы свойства многочленов от \(D\) из п.~\ref{subsec:poly-D} (см. замечание в п.~\ref{rem:LD-variable}). В доказательстве они и не нужны: используется только то, что \(L(D)K(x,s)=0\) при каждом фиксированном \(s\).
\end{remarkHSE}

\begin{exampleHSE}[№\,578: \(y''+4y=\operatorname{tg}x\)]
ФСР \(=\{\cos2x,\ \sin2x\}\) (см. пример в методе вариации постоянных). Ищем ядро Коши:
\[
K(x,s)=c_1(s)\cos2x+c_2(s)\sin2x,\qquad
K'_x(x,s)=-2c_1(s)\sin2x+2c_2(s)\cos2x.
\]
Условия \eqref{eq:cauchy-kernel-cond} при \(n=2\):
\[
\begin{cases}K(s,s)=0,\\K'_x(s,s)=1\end{cases}
\iff
\begin{cases}c_1(s)\cos2s+c_2(s)\sin2s=0,\\-2c_1(s)\sin2s+2c_2(s)\cos2s=1.\end{cases}
\]
Определители:
\[
\Delta=\begin{vmatrix}\cos2s&\sin2s\\-2\sin2s&2\cos2s\end{vmatrix}=2,
\]
\[
\Delta_1=\begin{vmatrix}0&\sin2s\\1&2\cos2s\end{vmatrix}=-\sin2s,\qquad
\Delta_2=\begin{vmatrix}\cos2s&0\\-2\sin2s&1\end{vmatrix}=\cos2s.
\]
Отсюда
\[
c_1(s)=\frac{\Delta_1}{\Delta}=-\frac{\sin2s}{2},\qquad
c_2(s)=\frac{\Delta_2}{\Delta}=\frac{\cos2s}{2},
\]
\[
K(x,s)=-\frac{\sin2s}{2}\cos2x+\frac{\cos2s}{2}\sin2x=\frac{\sin2(x-s)}{2}.
\]
Правая часть \(\operatorname{tg}x\) непрерывна на \(\bigl(-\frac\pi2;\frac\pi2\bigr)\); возьмём \(\gamma=0\):
\[
y_{\text{ч.н.}}=\int_0^x\frac{\sin2(x-s)}{2}\operatorname{tg}s\dd s,\quad x\in\Bigl(-\frac\pi2;\frac\pi2\Bigr).
\]
Интеграл берётся: раскрыв \(\sin2(x-s)=\sin2x\cos2s-\cos2x\sin2s\) и учитывая \(\sin2s\operatorname{tg}s=1-\cos2s\), \(\cos2s\operatorname{tg}s=\sin2s-\operatorname{tg}s\), получаем
\[
y_{\text{ч.н.}}=-\frac x2\cos2x+\frac12\sin2x\,\ln\cos x+\frac14\sin2x.
\]
Слагаемое \(\frac14\sin2x\) --- решение ЛОУ, так что с точностью до него это то же частное решение, что и найденное методом вариации постоянных.
\end{exampleHSE}

\section[ЛОС с постоянными коэффициентами]{ЛОС с постоянными коэффициентами}

Рассмотрим систему
\[
\dot{\vect{x}}=A\vect{x},
\qquad A=\text{const}.
\]

\begin{definitionHSE}[Характеристическое уравнение системы]
Уравнение
\[
\det(A-\lambda E)=0
\]
называется характеристическим уравнением системы. Его корни $\lambda$ --- собственные значения матрицы $A$.
\end{definitionHSE}

\begin{definitionHSE}[Собственный вектор]
Ненулевой вектор $\vect{v}$ называется собственным вектором, отвечающим собственному значению $\lambda$, если
\[
A\vect{v}=\lambda\vect{v}.
\]
\end{definitionHSE}

\begin{definitionHSE}[Присоединённые векторы]
В исходном конспекте вводится цепочка присоединённых векторов. В стандартной матричной записи она имеет вид
\[
(A-\lambda E)\vect{h}_1=\vect{v},
\qquad
(A-\lambda E)\vect{h}_2=\vect{h}_1,
\quad\ldots
\]
\end{definitionHSE}

\begin{theoremHSE}[Решение по собственному вектору]
Если $\vect{v}\neq0$ --- собственный вектор матрицы $A$, соответствующий собственному значению $\lambda$, то
\[
\vect{x}(t)=e^{\lambda t}\vect{v}
\]
является решением ЛОС.
\end{theoremHSE}

\begin{proof}
\[
\frac{\dd}{\dd t}\bigl(e^{\lambda t}\vect{v}\bigr)
=\lambda e^{\lambda t}\vect{v},
\]
а
\[
A\bigl(e^{\lambda t}\vect{v}\bigr)
=e^{\lambda t}A\vect{v}
=\lambda e^{\lambda t}\vect{v}.
\]
Обе части системы совпадают.
\end{proof}

\begin{remarkHSE}[Объём исходного конспекта]
В приложенном вопроснике для этого билета подробно выписаны определения характеристического уравнения, собственных и присоединённых векторов и доказана формула $e^{\lambda t}\vect v$. Полная формула общего решения через все жордановы цепочки в тексте источника отдельно не расписана, поэтому здесь она не добавлялась.
\end{remarkHSE}

\section[Векторный квазимногочлен]{Частное решение ЛНС с векторным квазимногочленом}

\begin{theoremHSE}[Форма частного решения]
Рассмотрим ЛНС с постоянной матрицей
\[
\dot{\vect{x}}=A\vect{x}+e^{\mu t}\vect{P}_N(t),
\]
где $\vect{P}_N(t)$ --- векторный многочлен степени $N$.

Частное решение ищется в виде
\[
x_i(t)=G_i(t)e^{\mu t},
\]
где $G_i(t)$ --- многочлен степени не выше $N+S$. Здесь
\begin{itemize}
  \item $S=0$, если $\mu$ не является собственным значением $A$;
  \item если $\mu$ --- собственное значение $A$, то $S$ берётся равным максимальному размеру жордановой клетки, соответствующей $\mu$.
\end{itemize}
\end{theoremHSE}

Если правая часть является суммой нескольких векторных квазимногочленов, частные решения для каждого слагаемого находят отдельно и складывают.

\begin{exampleHSE}[Векторный квазимногочлен]
Для $n=2$, $N=3$ в источнике приведён пример
\[
e^{2t}
\begin{pmatrix}
17t+4\\
2t^3+4t+5
\end{pmatrix}.
\]
\end{exampleHSE}

\section[Автономные системы]{Автономная система и фазовый портрет}

\begin{definitionHSE}[Автономная система]
Система ОДУ называется автономной, если независимая переменная $t$ не входит в неё явно:
\[
\dot{\vect{x}}=\vect{f}(\vect{x}).
\]
В координатах:
\[
\begin{cases}
\dot x_1=f_1(x_1,\ldots,x_n),\\
\qquad\vdots\\
\dot x_n=f_n(x_1,\ldots,x_n).
\end{cases}
\]
\end{definitionHSE}

Кинематически $t$ понимается как время, $\vect{x}(t)$ --- положение точки, а $\dot{\vect{x}}(t)$ --- её скорость. В автономной системе скорость зависит только от положения точки.

\begin{definitionHSE}[Фазовое пространство]
Фазовым пространством называется множество $\widetilde G\subset\R^n$, в котором определены координаты системы и её правая часть.
\end{definitionHSE}

\begin{definitionHSE}[Фазовая траектория]
Решение $\vect{x}(t)=\vect{\varphi}(t)$ задаёт в фазовом пространстве линию. Эта линия называется фазовой траекторией.
\end{definitionHSE}

\begin{definitionHSE}[Положение равновесия]
Точка $\vect{a}$ является стационарной точкой, положением равновесия или особой точкой, если
\[
\vect{f}(\vect{a})=\vect{0}.
\]
В таком случае $\vect{x}(t)\equiv\vect a$ --- решение.
\end{definitionHSE}

\begin{definitionHSE}[Фазовый портрет]
Совокупность всех фазовых траекторий автономной системы называется её фазовым портретом.
\end{definitionHSE}

\begin{theoremHSE}[Сдвиг времени]
Если $\vect{\varphi}(t)$ --- решение автономной системы, то для любого $c\in\R$
\[
\vect{\psi}(t)=\vect{\varphi}(t+c)
\]
также является решением. Эти решения задают одну и ту же фазовую траекторию.
\end{theoremHSE}

\begin{proof}
\[
\dot{\vect{\psi}}(t)
=\dot{\vect{\varphi}}(t+c)
=\vect{f}(\vect{\varphi}(t+c))
=\vect{f}(\vect{\psi}(t)).
\]
\end{proof}

\begin{theoremHSE}[Фазовые траектории не пересекаются]
Две фазовые траектории автономной системы либо совпадают, либо не пересекаются.
\end{theoremHSE}

\begin{proof}
Пусть $\vect{\varphi}(t_1)=\vect{\psi}(t_2)$. Сдвинем первое решение по времени на $c=t_1-t_2$:
\[
\vect{\chi}(t)=\vect{\varphi}(t+c).
\]
Тогда $\vect{\chi}(t_2)=\vect{\psi}(t_2)$. По единственности решения одной задачи Коши $\vect{\chi}$ и $\vect{\psi}$ совпадают, а значит совпадают и их фазовые траектории.
\end{proof}

\begin{corollaryHSE}[Следствия]
\begin{itemize}
  \item Решение не может войти в особую точку за конечное время.
  \item Если нестационарное решение возвращается в ту же фазовую точку, оно периодично.
  \item Фазовая траектория автономной системы является либо точкой, либо замкнутой кривой без самопересечений, либо незамкнутой кривой без самопересечений.
\end{itemize}
\end{corollaryHSE}

\section[Классификация особых точек]{Классификация изолированных особых точек}

Рассматривается уравнение
\[
\frac{\dd y}{\dd x}=\frac{cx+dy}{ax+by},
\]
которому соответствует линейная автономная система
\[
\begin{cases}
\dot x=ax+by,\\
\dot y=cx+dy.
\end{cases}
\]
Пусть
\[
A=\begin{pmatrix}a&b\\c&d\end{pmatrix},
\qquad \det A\neq0.
\]
Тогда $(0,0)$ --- единственная особая точка.

\begin{theoremHSE}[Классификационная теорема Пуанкаре]
Тип точки $(0,0)$ определяется собственными значениями $\lambda_1,\lambda_2$ матрицы $A$:
\begin{enumerate}
  \item $\lambda_1,\lambda_2\in\R$ разных знаков --- \textbf{седло};
  \item $\lambda_1,\lambda_2\in\R$ различны и одного знака --- \textbf{узел};
  \item $\lambda_1=\lambda_2\in\R$ --- \textbf{вырожденный} или \textbf{дикритический узел};
  \item $\lambda_{1,2}=\pm i\beta$ --- \textbf{центр};
  \item $\lambda_{1,2}=\alpha\pm i\beta$, $\alpha\neq0$ --- \textbf{фокус}.
\end{enumerate}
Знак вещественных частей определяет направление движения к особой точке или от неё.
\end{theoremHSE}

\begin{exampleHSE}[Типовые примеры из конспекта]
\begin{enumerate}
  \item Центр:
  \[
y'=-\frac{x}{y},\qquad
A=\begin{pmatrix}0&1\\-1&0\end{pmatrix},
\qquad \lambda_{1,2}=\pm i.
  \]
  \item Седло:
  \[
y'=\frac{x}{y},\qquad
A=\begin{pmatrix}0&1\\1&0\end{pmatrix},
\qquad \lambda_{1,2}=\pm1.
  \]
  \item Дикритический узел:
  \[
y'=\frac{y}{x},\qquad
A=\begin{pmatrix}1&0\\0&1\end{pmatrix},
\qquad \lambda_{1,2}=1.
  \]
  \item Узел:
  \[
y'=\frac{2y}{x},\qquad
A=\begin{pmatrix}1&0\\0&2\end{pmatrix},
\qquad \lambda_1=1,\ \lambda_2=2.
  \]
  \item Вырожденный узел:
  \[
y'=\frac{x+y}{x},\qquad
A=\begin{pmatrix}1&0\\1&1\end{pmatrix}.
  \]
  \item Фокус:
  \[
y'=\frac{x+y}{x-y},\qquad
A=\begin{pmatrix}1&-1\\1&1\end{pmatrix},
\qquad \lambda_{1,2}=1\pm i.
  \]
\end{enumerate}

% TODO: фазовые портреты из таблицы на странице 50 и рисунки со страницы 52.
\end{exampleHSE}

\section[Пуанкаре--Перрон]{Исключительные направления. Теорема Пуанкаре--Перрона}

\begin{definitionHSE}[Цикл]
Циклом уравнения $y'=f(x,y)$ называется замкнутая интегральная кривая, не проходящая через особую точку.
\end{definitionHSE}

\begin{definitionHSE}[Исключительное направление]
Исключительным называется направление, по которому в особую точку входит по крайней мере одна интегральная кривая.
\end{definitionHSE}

\subsection{Топологические типы особых точек}
\begin{itemize}
  \item \textbf{Центр}: в некоторой окрестности все интегральные кривые --- вложенные циклы; исключительных направлений нет.
  \item \textbf{Фокус}: интегральные кривые --- спирали, асимптотически приближающиеся к особой точке; исключительных направлений нет.
  \item \textbf{Седло}: две пары исключительных направлений; по ним входят сепаратрисы, остальные кривые в точку не входят.
  \item \textbf{Простой узел}: две пары исключительных направлений; одна пара кривых входит по одному направлению, остальные --- по другому.
  \item \textbf{Вырожденный узел}: одна пара исключительных направлений; все интегральные кривые входят в особую точку.
  \item \textbf{Дикритический узел}: любое направление исключительное; по каждому направлению входит ровно одна интегральная кривая.
  \item \textbf{Центро-фокус}: в любой достаточно малой окрестности встречаются и циклы, и спирали.
\end{itemize}

Рассмотрим автономную систему
\[
\begin{cases}
\dot x=P(x,y),\\
\dot y=Q(x,y),
\end{cases}
\]
и изолированное положение равновесия $(0,0)$. Линеаризация имеет вид
\[
\begin{cases}
P(x,y)=ax+by+p(x,y),\\
Q(x,y)=cx+dy+q(x,y),
\end{cases}
\]
где $p,q$ --- члены более высокого порядка.

\begin{theoremHSE}[Пуанкаре--Перрона]
Пусть
\[
\frac{\dd y}{\dd x}
=
\frac{cx+dy+q(x,y)}{ax+by+p(x,y)}
\]
--- полное уравнение, а
\[
\frac{\dd y}{\dd x}
=
\frac{cx+dy}{ax+by}
\]
--- укороченное. Пусть $\det A\neq0$, $(0,0)$ --- изолированная особая точка, выполнена локальная единственность, а $p,q$ имеют более высокий порядок малости.

Если собственные значения $A$ удовлетворяют
\[
\Re\lambda_{1,2}\neq0,
\]
то тип особой точки полного уравнения совпадает с типом особой точки укороченного уравнения; исключительные направления также совпадают.

Если
\[
\Re\lambda_{1,2}=0,
\]
то при центре линейной системы полная система может иметь центр, фокус или центро-фокус.
\end{theoremHSE}

\begin{exampleHSE}[Система $\dot x=y$, $\dot y=-x+x^2$]
Положения равновесия:
\[
(0,0),\qquad(1,0).
\]
В точке $(0,0)$ линейная часть
\[
A=\begin{pmatrix}0&1\\-1&0\end{pmatrix},
\qquad \lambda_{1,2}=\pm i.
\]
По теореме Пуанкаре--Перрона остаются варианты центр/фокус/центро-фокус; в исходнике дальнейшая симметрия используется для вывода о центре.

Для точки $(1,0)$ после замены
\[
u=x-1,\qquad v=y
\]
линейная часть имеет матрицу
\[
A=\begin{pmatrix}0&1\\1&0\end{pmatrix},
\qquad \lambda_{1,2}=\pm1,
\]
поэтому это седло.
\end{exampleHSE}

\section[Лотка--Вольтерра]{Модель Лотки--Вольтерры}

Пусть $x(t)$ --- количество жертв, $y(t)$ --- количество хищников. Модель:
\[
\begin{cases}
\dot x=\alpha x-\gamma xy,\\
\dot y=-\beta y+\delta xy,
\end{cases}
\qquad \alpha,\beta>0.
\]
В конспекте берётся частный случай
\[
\alpha=\beta=\gamma=\delta=1,
\]
то есть
\begin{equation}
\begin{cases}
\dot x=x-xy=x(1-y),\\
\dot y=-y+xy=y(x-1).
\end{cases}
\label{eq:lotka}
\end{equation}

Положения равновесия:
\[
(0,0),\qquad(1,1).
\]

\subsection{Точка \texorpdfstring{$(0,0)$}{(0,0)}}
Линеаризация:
\[
\begin{cases}
\dot x=x,\\
\dot y=-y,
\end{cases}
\qquad
A=\begin{pmatrix}1&0\\0&-1\end{pmatrix}.
\]
Собственные значения $1$ и $-1$, поэтому $(0,0)$ --- седло и неустойчивая точка.

\subsection{Точка \texorpdfstring{$(1,1)$}{(1,1)}}
Делаем сдвиг
\[
\widetilde x=x-1,
\qquad
\widetilde y=y-1.
\]
Тогда
\[
\begin{cases}
\dot{\widetilde x}=-\widetilde y-\widetilde x\widetilde y,\\
\dot{\widetilde y}=\widetilde x+\widetilde x\widetilde y.
\end{cases}
\]
Линейная часть:
\[
A=\begin{pmatrix}0&-1\\1&0\end{pmatrix},
\qquad \lambda_{1,2}=\pm i.
\]
Для укороченной системы это центр. В конспекте далее исследуется симметрия и знаки $\dot x$, $\dot y$, после чего строится фазовый портрет с замкнутыми траекториями вокруг $(1,1)$.

Знаки производных определяются прямыми $x=0$, $x=1$, $y=0$, $y=1$:
\[
\dot x=x(1-y),
\qquad
\dot y=y(x-1).
\]

% TODO: фазовый портрет Лотки--Вольтерры со страницы 59 исходника.

\section[Устойчивость по Ляпунову]{Устойчивость (и асимптотическая) по Ляпунову}

Рассмотрим систему
\[
\dot{\vect{x}}=\vect{f}(t,\vect{x})
\]
и её частное решение
\[
\vect{x}=\vect{\varphi}(t),
\qquad t\ge t_0.
\]

\begin{definitionHSE}[Устойчивость по Ляпунову]
Решение $\vect{\varphi}(t)$ называется устойчивым по Ляпунову, если для любого $\varepsilon>0$ существует $\delta=\delta(\varepsilon)>0$ такое, что всякое решение $\widetilde{\vect{x}}(t)$ с
\[
|\widetilde{\vect{x}}(t_0)-\vect{\varphi}(t_0)|<\delta
\]
продолжается вправо и удовлетворяет
\[
|\widetilde{\vect{x}}(t)-\vect{\varphi}(t)|<\varepsilon,
\qquad \forall t\ge t_0.
\]
\end{definitionHSE}

\begin{definitionHSE}[Асимптотическая устойчивость]
Решение $\vect{\varphi}(t)$ называется асимптотически устойчивым, если оно устойчиво по Ляпунову и существует $\delta_0>0$ такое, что для всех решений с начальными данными из $\delta_0$-окрестности
\[
\lim_{t\to+\infty}
|\widetilde{\vect{x}}(t)-\vect{\varphi}(t)|=0.
\]
\end{definitionHSE}

\begin{exampleHSE}[Проверка по определению]
Рассмотрим
\[
\dot x=1+t-x,
\qquad x(0)=0.
\]
Общее решение:
\[
x(t)=t+Ce^{-t}.
\]
Начальному условию $x(0)=0$ соответствует $C=0$, поэтому
\[
\varphi(t)=t.
\]
Для соседнего решения
\[
|\widetilde x(t)-\varphi(t)|=|C|e^{-t}.
\]
Если $|C|<\delta$ и выбрать $\delta(\varepsilon)=\varepsilon$, то
\[
|C|e^{-t}<\varepsilon
\qquad \forall t\ge0.
\]
Кроме того,
\[
\lim_{t\to+\infty}|C|e^{-t}=0.
\]
Следовательно, $\varphi(t)=t$ асимптотически устойчива.
\end{exampleHSE}

\begin{exampleHSE}[Одновременно устойчивое и неустойчивое стационарные решения]
Для нелинейного уравнения
\[
\dot x=1-x^2
\]
стационарные решения равны
\[
x\equiv-1,
\qquad x\equiv1.
\]
По фазовой прямой видно: $x=-1$ неустойчиво, а $x=1$ устойчиво.

% TODO: графический разбор со страницы 62 исходника.
\end{exampleHSE}

\section[Устойчивость линейных систем]{Устойчивость ЛНС и нулевого решения ЛОС}

Пусть
\begin{equation}
\dot{\vect{x}}=A(t)\vect{x}
\label{eq:los-stab}
\end{equation}
--- ЛОС, а
\begin{equation}
\dot{\vect{y}}=A(t)\vect{y}+\vect{F}(t)
\label{eq:lns-stab}
\end{equation}
--- соответствующая ЛНС.

\begin{theoremHSE}[Сведение устойчивости ЛНС к ЛОС]
При непрерывных $A(t)$ и $\vect{F}(t)$ на $[t_0,+\infty)$:
\begin{enumerate}
  \item все решения ЛНС устойчивы тогда и только тогда, когда нулевое решение ЛОС устойчиво;
  \item все решения ЛНС асимптотически устойчивы тогда и только тогда, когда нулевое решение ЛОС асимптотически устойчиво.
\end{enumerate}
\end{theoremHSE}

\begin{proof}
Пусть $\vect{\psi}(t)$ --- фиксированное решение ЛНС. Любое другое решение ЛНС имеет вид
\[
\vect{y}(t)=\vect{\psi}(t)+\vect{x}(t),
\]
где $\vect{x}(t)$ --- решение ЛОС. Поэтому
\[
|\vect{y}(t)-\vect{\psi}(t)|=|\vect{x}(t)|.
\]
Все условия устойчивости и асимптотической устойчивости для $\vect{\psi}$ буквально превращаются в соответствующие условия для нулевого решения ЛОС. Обратное рассуждение проводится тем же равенством.
\end{proof}

\begin{theoremHSE}[Критерий для ЛОС с постоянными коэффициентами]
Пусть
\[
\dot{\vect{x}}=A\vect{x},
\qquad A=\text{const},
\]
и $\lambda_1,\ldots,\lambda_k$ --- собственные значения $A$.
\begin{enumerate}
  \item Если для всех $j$
  \[
  \Re\lambda_j<0,
  \]
  то нулевое решение асимптотически устойчиво.
  \item Если
  \[
  \Re\lambda_j\le0
  \]
  для всех $j$, а каждому собственному значению с нулевой вещественной частью соответствуют только жордановы клетки размера $1\times1$, то нулевое решение устойчиво.
  \item В остальных случаях нулевое решение неустойчиво: в частности, если существует $\Re\lambda_j>0$ или при $\Re\lambda_j=0$ имеется жорданова клетка размера не меньше $2\times2$.
\end{enumerate}
\end{theoremHSE}

\section[Первое приближение]{Исследование устойчивости по первому приближению}

Рассмотрим линейную систему
\[
\dot{\vect{x}}=A\vect{x},
\qquad A=\text{const},
\]
и возмущённую систему
\begin{equation}
\dot{\vect{x}}=A\vect{x}+\vect{\varphi}(t,\vect{x}),
\label{eq:first-approx}
\end{equation}
где $\vect{\varphi}$ и её частные производные непрерывны.

Пусть при $\vect{x}\to0$
\[
|\vect{\varphi}(t,\vect{x})|
\le
\gamma(\vect{x})|\vect{x}|,
\qquad
\gamma(\vect{x})\to0.
\]
То есть нелинейная часть имеет более высокий порядок малости по сравнению с $|\vect{x}|$.

\begin{theoremHSE}[Устойчивость по первому приближению]
Пусть $\vect{x}=0$ --- решение \eqref{eq:first-approx}, а $\lambda$ --- собственные значения матрицы $A$.
\begin{enumerate}
  \item Если для всех собственных значений
  \[
  \Re\lambda<0,
  \]
  то нулевое решение полной системы асимптотически устойчиво.
  \item Если существует собственное значение
  \[
  \Re\lambda>0,
  \]
  то нулевое решение полной системы неустойчиво.
\end{enumerate}
\end{theoremHSE}

\begin{exampleHSE}[Система $\dot x=\sin x$, $\dot y=\cos y$]
Рассмотрим положение равновесия
\[
\left(0,\frac{\pi}{2}\right).
\]
Сдвигаем координаты:
\[
\widetilde y=y-\frac{\pi}{2}.
\]
Тогда
\[
\begin{cases}
\dot x=\sin x,\\
\dot{\widetilde y}=-\sin\widetilde y.
\end{cases}
\]
Разложение в окрестности нуля:
\[
\dot x=x-\frac{x^3}{3!}+O(x^5),
\]
\[
\dot{\widetilde y}
=-\widetilde y+\frac{\widetilde y^3}{3!}+O(\widetilde y^5).
\]
Матрица первого приближения
\[
A=\begin{pmatrix}1&0\\0&-1\end{pmatrix},
\qquad \lambda_{1,2}=\pm1.
\]
Поскольку есть собственное значение с положительной вещественной частью, положение равновесия неустойчиво. Линейная часть имеет тип седла.
\end{exampleHSE}

\section{Алгоритмы решения ОДУ}

\begin{definitionHSE}[Семейство кривых]
Семейство кривых --- все кривые, получаемые из уравнения
\[\Phi(x,y,C_1,\ldots,C_n)=0\]
при различных значениях произвольных постоянных \(C_1,\ldots,C_n\).
\end{definitionHSE}

Пусть дано семейство кривых
\[\Phi(x, y, C_1,\ldots, C_n) = 0.\]

Чтобы получить дифференциальное уравнение этого семейства, дифференцируем соотношение по \(x\) столько раз, сколько имеется независимых постоянных, а затем исключаем \(C_1,\ldots,C_n\).

\subsection{Линейные и однородные уравнения}

\phantomsection\label{box:type0}
\begin{definitionHSE}[Уравнения типа 0]
Если \(y' = f(x)\), то семейство решений имеет вид:
\begin{gather*}
y'=f(x),\quad\frac{\dd y}{\dd x}=f(x),\quad\int\dd y = \int f(x)\dd x,\\
\boxed{\ y = \int f(x)\dd x + C, \quad C \in \mathbb{R}.\ }
\end{gather*}
\end{definitionHSE}

\phantomsection\label{box:type1}
\begin{definitionHSE}[Уравнения типа 1 (с разделяющимися переменными)]
Рассмотрим уравнение
\[
y' = f(x)g(y).
\]
Отыщем стационарные решения. Стационарное решение должно обращать ОДУ в тождество:
\begin{itemize}
\item \(y=y_0-\text{const}: f(x)g(y_0)=0\);
\item \(x=x_0-\text{const}: \dfrac{1}{f(x_0)g(y)}=0\).
\end{itemize}
Далее разделяем переменные:
\[
\frac{y'}{g(y)} = f(x),\quad\frac{\frac{\dd y}{\dd x}}{g(y)}=f(x)
\]
и после интегрирования
\[
\boxed{\ \int \frac{\dd y}{g(y)} = \int f(x)\dd x + C, \quad C \in \mathbb{R}.\ }
\]
То есть изначально мы фактически интегрируем по \(x\), но в левой части делаем замену переменной:
\[
\int \frac{y'(x)}{g(y(x))}\dd x = \int f(x)\dd x + C, \quad C \in \mathbb{R}.
\]

Получаем, что решение описывается однопараметрическим семейством и набором стационарных решений.

\end{definitionHSE}

\phantomsection\label{box:type2}
\begin{definitionHSE}[Уравнения типа 2 (с линейной правой частю)]
Пусть
\[y'=f(ax+by+c)\]
Вводим \(z=ax+by+c\). Тогда
\[z'=a+by'=a+bf(z),\quad y=\frac{z-ax-c}{b}\]
Получаем уравнение на \(z\) с разделяющимися переменными (\hyperref[box:type1]{тип 1}):
\[\boxed{\ z'=a+bf(z) \ }\]
\end{definitionHSE}

\begin{exampleHSE}[\(y'-y=2x-3\)]
\begin{gather*}
y'-y=2x-3,\quad y'=\underbrace{2x+y-3}_{=z}\quad\Longrightarrow\quad z'=(2x+y-3)'=2+y'=2+z
\end{gather*}
Стационарное решение перед делением:
\[z'=0\Rightarrow z+2=0\Rightarrow z=-2-\text{решение}\]
Делим:
\begin{gather*}
\frac{\dd z}{z+2}=\dd x,\quad\int\frac{\dd z}{z+2}=x+C,\ C\in\R\\
\ln|z+2|=x+C,\quad z+2=\pm\e^{x+C}\\
z+2=\pm C_1\e^x\quad\xRightarrow[]{\text{стац. реш.}}\quad z=C_1\e^x-2
\end{gather*}
\end{exampleHSE}

\begin{remarkHSE}
Часто произвольную постоянную при интегрировании удобнее записать как \(\ln C, C > 0\). Тогда можно чуть сократить количество переобозначений произвольных постоянных.
\end{remarkHSE}

Cтационарные решения по \(x\) (\(x=x_0=\text{const}\)) не являются решениями в классическом смысле, поэтому в дальнейшем мы ищем их только в следующих случаях:

\begin{itemize}
\item \(M(x,y)\dd x+N(x,y)\dd y=0\).
\item Когда ищем интегральные кривые.
\item \(y'=f(x,y)\).
\end{itemize}

\phantomsection\label{box:type3}
\begin{definitionHSE}[Уравнения типа 3 (с однородной правой частью)]
Рассмотрим уравнение
\[y'=f\left(\frac{y}{x}\right)\]
\begin{definitionHSE}[Определение однородной функции]
Функция \(f(x,y)\) называется \textit{однородной} порядка \(k\), если
\[\forall\lambda\neq0\quad f(\lambda x,\lambda y)=\lambda^kf(x,y)\]
\end{definitionHSE}
Сделаем замену:
\vspace{-\baselineskip}
\begin{gather*}
y=u(x)\cdot x,\quad y'=\boxed{u'x+u=f(u)}=f\left(\frac{y}{x}\right)\\
u'=\frac{f(u)-u}{x},\quad\frac{\dd u}{f(u)-u}=\frac{\dd x}{x}
\end{gather*}
Далее решаем уравнение с \hyperref[box:type1]{разделяющимися перменными}.
\end{definitionHSE}
Уравнение
\[M(x,y)\dd x+N(x,y)\dd y=0\]
называется однородным, если \(M(x,y)\) и \(N(x,y)\) --- однородные функции одного и того же порядка. Оно решается таким же способом.

\begin{exampleHSE}[\(y'=\dfrac{y}{x}\cos\ln\left(\dfrac{y}{x}\right)\)]
\[y'=\frac{y}{x}\cos\ln\left(\frac{y}{x}\right),\quad y=ux,\quad\frac{y}{x}>0\]
\[
\begin{aligned}
y'&=(ux)'=u'x+u,\\
y'&=u\cos\ln u
\end{aligned}\quad\Longrightarrow\quad u'x+u=u\cos\ln u,\quad u'=\frac{u\cos\ln u-u}{x}\]

Отыщем стационарные решения:
\[u\cos\ln u-u=0,\quad\cos\ln u=1,\quad \ln u=2\pi k,\ k\in\Z,\quad u=\e^{2\pi k}\]
Продолжим решать:
\begin{gather*}
\frac{\dd u}{\dd x}=\frac{u\cos\ln u-u}{x},\quad\frac{\dd u}{u\cos\ln u-u}=\frac{\dd x}{x}\\
\int\frac{\dd u}{u\cos\ln u-u}=|\ln u=t|=\int\frac{\dd t}{\cos t-1}=\int\frac{\dd t}{\cos t-1}=\int\frac{\dd t}{-2\sin^2(t/2)}=\\=-\frac{1}{2}\cdot\left(-2\ctg\left(\frac{t}{2}\right)\right)+C=\ctg\left(\frac{t}{2}\right)+C=\ctg\left(\frac{\ln u}{2}\right)+C,\ C\in\R\\
\ctg\left(\frac{\ln u}{2}\right)=\ln|x|+C_1,\quad\boxed{\ctg\left(\frac{\ln (x/y)}{2}\right)=\ln|C_2x|\ (C_2\neq0)}
\end{gather*}
Сделаем обратную замену в стационарном решении:
\[\boxed{y=e^{2\pi k}x,\ k\in\Z}\]
\end{exampleHSE}

\begin{exampleHSE}[\((x+2y)\dd x-x\dd y=0\)]
$M=x+2y$ и $N=-x$ --- однородные функции первого порядка, значит, уравнение однородное.

$x=0$ --- решение ($\dd x=0$, и уравнение обращается в тождество). Делаем замену $y=ux$, $\dd y=x\dd u+u\dd x$:
\begin{gather*}
(x+2ux)\dd x-x(x\dd u+u\dd x)=0,\quad (x+ux)\dd x-x^2\dd u=0,\\
x(1+u)\dd x=x^2\dd u\quad\bigl|\,:(1+u)x^2
\end{gather*}
Когда делим, смотрим, не потеряли ли стационарные решения: $1+u=0$ --- решение.
\begin{gather*}
\frac{\dd x}{x}=\frac{\dd u}{1+u},\qquad
\ln|x|=\ln|1+u|+\ln C_1,\ C_1>0,\\
|x|=|1+u|\cdot C_1,\qquad x=\pm(1+u)\,C_1.
\end{gather*}
\[
\left[\begin{aligned}
&x=(1+u)\,C,\ C\in\R,\\
&u=-1
\end{aligned}\right.
\ \longrightarrow\
\left[\begin{aligned}
&x=\Bigl(1+\frac yx\Bigr)C,\ C\in\R,\\
&\frac yx=-1
\end{aligned}\right.
\ \longrightarrow\
\boxed{\left[\begin{aligned}
&x^2=Cx+Cy,\ C\in\R,\\
&y=-x
\end{aligned}\right.}
\]
Решение $x=0$ входит в семейство при $C=0$.
\end{exampleHSE}

\phantomsection\label{box:type4}
\begin{definitionHSE}[Уравнения типа 4 (уравнением с дробно-линейной правой частью)]
Рассмотрим уравнение
\[y'=f\left(\frac{ax+by+c}{a_1x+b_1y+c_1}\right),\quad\begin{vmatrix}a&b\\a_1&b_1\end{vmatrix}\neq0\]
Идея --- убрать свободные члены, чтобы уравнение стало однородным. Сделаем замену
\[
\begin{aligned}
x&=x_0+t,\\y&=y_0+w,
\end{aligned}
\text{ где }x_0,y_0\text{ находим из }\begin{cases}
ax_0+by_0+c=0\\a_1x_0+b_1y_0+c_1=0
\end{cases}
\]
Тогда
\[ax+by+c=a(x_0+t)+b(y_0+w)+c=(ax_0+by_0+c)+at+bw=at+bw,\]
Аналогично
\[a_1x+b_1y+c_1=a_1t+b_1w\]
Таким образом
\[w'=f\left(\frac{at+bw}{a_1t+b_1w}\right)=\left(\frac{a+b(w/t)}{a_1+b_1(w/t)}\right)\]
\end{definitionHSE}

\begin{exampleHSE}[\((x+4y)y'=2x+3y-5\)]
\[
y'=\frac{2x+3y-5}{x+4y},\qquad y=-\frac{x}{4}\ \text{--- не решение.}
\]
Определитель $\begin{vmatrix}2&3\\1&4\end{vmatrix}=5\neq0$. Центр сдвига:
\[
\left( \begin{array}{cc|c} 
2 & 3 & 5 \\ 
1 & 4 & 0 
\end{array} \right) 
\rightarrow 
\left( \begin{array}{cc|c} 
0 & 1 & -1 \\ 
1 & 0 & 4 
\end{array} \right) 
\rightarrow 
\begin{array}{l} 
y_0 = -1 \\ 
x_0 = 4 
\end{array}
\]
\[
\begin{cases}2x_0+3y_0-5=0\\x_0+4y_0=0\end{cases}\quad\Longrightarrow\quad x_0=4,\ y_0=-1.
\]
Замена $x=4+t$, $y=-1+w$, затем $w=u(t)\cdot t$:
\begin{gather*}
w'=\frac{2t+3w}{t+4w}=\frac{2+3\frac{w}{t}}{1+4\frac{w}{t}},\qquad
u't+u=\frac{2+3u}{1+4u},\\
\frac{\dd u}{\dd t}=\frac{-4u^2+2u+2}{(1+4u)\,t}=\frac{-2(u-1)(2u+1)}{(1+4u)\,t}.
\end{gather*}
Стационарные решения: $u=1$, $u=-1/2$. Разделяем переменные:
\begin{gather*}
\frac{(1+4u)\dd u}{-4(u-1)\left(u+\frac12\right)}=\frac{\dd t}{t},\qquad
\frac{1+4u}{-4(u-1)\left(u+\frac12\right)}=-\frac{5/6}{u-1}-\frac{1/6}{u+\frac12},\\
\frac56\ln|u-1|+\frac16\ln\Bigl|u+\frac12\Bigr|=-\ln|t|+\ln C_1\ (C_1>0),\\
|u-1|^5\Bigl|u+\frac12\Bigr|\,|t|^6=C_2\ (C_2>0)
\ \xRightarrow[]{\text{стац. реш.}}\
(u-1)^5\Bigl(u+\frac12\Bigr)t^6=C_3,\ (C_3\in\R).
\end{gather*}
Обратная замена $u=w/t=(y+1)/(x-4)$:
\[
\boxed{\ (y-x+5)^5\Bigl(y+\frac{x}{2}-1\Bigr)=C\ (C\in\R)\ }
\]
\end{exampleHSE}

\phantomsection\label{box:type5}
\begin{definitionHSE}[Уравнения типа 5 (сводящиеся к однородным степенной заменой)]
Если уравнение не однородно, но становится однородным после замены
\[
y=z^m\Longrightarrow y'=m\,z^{m-1}z',
\]
то показатель $m$ подбираем так, чтобы все слагаемые получили одинаковую степень, и уравнение стало \hyperref[box:type3]{однородным}.
\end{definitionHSE}
\begin{remarkHSE}[Как распознать применение метода?]
Что-то должно «мешать» быть уравнению однородным. Пример:
\begin{gather*}
y'=\frac{y^3+xy}{2x^2}=\frac{y}{2x}\left(1+\frac{y^2}{x}\right)\Longrightarrow z=y^2\\
z'=2yy'=2y\cdot\frac{y}{2x}\left(1+\frac{y^2}{x}\right)=\frac{z}{x}\left(1+\frac{z}{x}\right)
\end{gather*}
\end{remarkHSE}
\begin{exampleHSE}[\(y'=\dfrac{y^3+xy}{2x^2}\)]
$y=0$ --- решение. Замена $y=z^m$:
\[
m\,z^{m-1}z'=\frac{z^{3m}+x\,z^m}{2x^2}\ \Longrightarrow\
z'=\frac{z^{2m+1}+x\,z}{2m\,x^2},\qquad 2m+1=2\ \Longrightarrow\ m=\frac12.
\]
Итак, $y^2=z$ и
\[
z'=\frac{z^2+xz}{x^2}\ \text{--- однородное.}
\]
Замена $z=ux$:
\begin{gather*}
u'x+u=u^2+u,\quad x\,\frac{\dd u}{\dd x}=u^2,\quad u=0\ \text{--- решение},\\
\frac{\dd u}{u^2}=\frac{\dd x}{x},\quad -\frac1u=\ln|x|+C,\quad u=-\frac{1}{\ln|x|+C}.
\end{gather*}
Обратная замена $z=ux$, $y^2=z$:
\[
\boxed{\ y^2=-\frac{x}{\ln|x|+C},\ C\in\R;\quad y=0\ }
\]
\end{exampleHSE}

\phantomsection\label{box:type6}
\begin{definitionHSE}[Уравнения типа 6 (линейные)]
\textit{Лниейнное однородное} уравнение:
\[y'+a(x)y=0.\]
Всегда есть стационарное решение $y=0$; при $y\neq0$
\begin{gather*}
\frac{\dd y}{y}=-a(x)\dd x,\qquad \ln|y|=-\int a(x)\dd x+\ln C_0,\ C_0>0,\\
|y|=C_0\,\e^{-\int a(x)\dd x},\quad y=\pm C_0\,\e^{-\int a(x)\dd x}
\end{gather*}
Объединяем со стационарным решением:
\[\boxed{\ y=C\,\e^{-\int a(x)\dd x},\quad C\in\R\ }\]
Представим \(y\) в следующем виде: \(y=C\cdot u(x)\), где \(u(x)\) --- функция, удовлетворяющая однородному уравнению

\textit{Лниейнное неоднородное} уравнение:
\[y'+a(x)y=b(x).\]
\textcircled{1} Сначала решаем соответствующее однородное уравнение. Выражаем \(y=C\cdot u(x)\).

\textcircled{2} Используем \textit{метод вариации постоянной} $C=C(x)$, $y=C(x)u(x)$:
\begin{gather*}
C'u+Cu'+a(x)Cu=b(x),\quad C'u+C\underbrace{\bigl(u'+a(x)u\bigr)}_{=0}=b(x),\\
\underline{C'(x)=\frac{b(x)}{u(x)}},\quad C(x)=\int\frac{b(x)}{u(x)}\dd x+\widetilde C,
\end{gather*}
Подставим \(C(x)\) в \(y=C(x)u(x)\) из первого пункта:
\[
\boxed{\ y=\left(\int\frac{b(x)}{u(x)}\dd x+\widetilde C\right)u(x),\quad\widetilde C\in\R\ }
\]
\end{definitionHSE}

\begin{exampleHSE}[\((2x+1)y'=4x+2y\)]
\begin{gather*}
y'=\frac{4x+2y}{2x+1}=\frac{2}{2x+1}\,y+\frac{4x}{2x+1},\qquad
y'-\frac{2}{2x+1}\,y=\frac{4x}{2x+1}.
\end{gather*}
\textcircled{1} Однородное уравнение: $y=0$ --- решение;
\[
\frac{\dd y}{y}=\frac{2\dd x}{2x+1},\qquad \ln|y|=\ln|2x+1|+\ln C_0,\qquad y=C(2x+1),\ C\in\R.
\]
\textcircled{2} Вариация постоянной, $y=C(x)(2x+1)$:
\begin{gather*}
C'(2x+1)=\frac{4x}{2x+1},\qquad C'=\frac{4x}{(2x+1)^2}=\frac{2}{2x+1}-\frac{2}{(2x+1)^2},\\
C=\ln|2x+1|+\frac{1}{2x+1}+\widetilde C.
\end{gather*}
\[
\boxed{\ y=(2x+1)\ln|2x+1|+1+\widetilde C(2x+1),\quad \widetilde C\in\R\ }
\]
\end{exampleHSE}

\phantomsection\label{box:type6-1}
\begin{definitionHSE}[Уравнения типа 6.1 (линейные относительно \(x\))]
Если уравнение линейно относительно функции $x(y)$, переходим к ней, используя
\[
y'=\frac{1}{x'_y},
\]
и решаем как \hyperref[box:type6]{уравнение типа 6} с независимой переменной $y$.
\end{definitionHSE}
\begin{remarkHSE}[Как распознать подобное уравнение?]
\(y'\) умножается на выражение, линейное по \(x\), а в другой части стоит функция от \(y\).

Практический признак для таких примеров:
\[\bigl(a(y)x+b(y)\bigr)y'=c(y)\]
\end{remarkHSE}

\begin{exampleHSE}[\((\sin^2y+x\ctg y)\,y'=1\)]
\[
(\sin^2y+x\ctg y)\,\frac{1}{x'_y}=1\quad\Longrightarrow\quad x'_y-x\ctg y=\sin^2y.
\]
\textcircled{1} Однородное: $x=0$ --- решение;
\[
\frac{\dd x}{x}=\ctg y\dd y,\quad \ln|x|=\ln|\sin y|+\ln C_1,\quad x=C\sin y,\ C\in\R.
\]
\textcircled{2} $C=C(y)$, $x=C(y)\sin y$:
\[
C'\sin y+C\cos y-\underbrace{x\ctg y}_{=C\cos y}=\sin^2y,\quad C'=\sin y,\quad C=-\cos y+\widetilde C.
\]
\[
\boxed{\ x=(\widetilde C-\cos y)\sin y,\quad\widetilde C\in\R\ }
\]
\end{exampleHSE}

\phantomsection\label{box:type8}
\begin{definitionHSE}[Уравнения типа 8 (уравнение Бернулли)]
Это линейное неоднородное уравнение, свободный член которого домножили на \(y^n\):
\[
y'+a(x)y=b(x)\,y^n,\qquad n\neq0,1.
\]
При $n>0$ будет решение $y=0$. Делим на $y^n$ и делаем замену $z=y^{1-n}\Rightarrow z'=(1-n)y^{-n}y'$:
\begin{gather*}
y'y^{-n}+a(x)\,\underbrace{y^{1-n}}_{=z}=b(x)\quad\Longrightarrow\quad\frac{z'}{1-n}+a(x)z=b(x)\\
\boxed{\ z'+(1-n)\,a(x)\,z=(1-n)\,b(x)\ }
\end{gather*}
--- линейное уравнение (\hyperref[box:type6]{тип 6}).
\end{definitionHSE}

\begin{exampleHSE}[\(y'=y^4\cos x+y\tg x\)]
\[
y'-y\tg x=y^4\cos x,\quad y=0\ \text{--- решение},\quad \frac{y'}{y^4}-\frac{\tg x}{y^3}=\cos x.
\]
Замена $z=y^{-3}$, $z'=-3y^{-4}y'$:
\[
z'+3\tg x\cdot z=-3\cos x.
\]
\textcircled{1} Однородное: $\dfrac{\dd z}{z}=-3\tg x\dd x$, $\ln|z|=3\ln|\cos x|+\ln C_1$, $z=C\cos^3x$.

\textcircled{2} Вариация, $C=C(x)$:
\[
C'\cos^3x-\cancel{3C\cos^2x\sin x}+\cancel{3C\cos^3x\tg x}=-3\cos x,
\]
\[
C'=-\frac{3}{\cos^2x},\qquad C=-3\tg x+\widetilde C.
\]
\[
\boxed{\ y^{-3}=(\widetilde C-3\tg x)\cos^3x,\ \widetilde C\in\R;\quad y=0\ }
\]
\end{exampleHSE}

\phantomsection\label{box:type9}
\begin{definitionHSE}[Уравнения типа 9 (в полных дифференциалах)]
Мы предполагаем, что левая часть этого уравнения — это полный дифференциал \(\dd u\) некоторой функции \(u(x,y)\) (которую мы называем потеницалом).
\[
\begin{array}{r@{}c@{}l@{}c@{}l}
          & M(x,y)                         & \,\dd x+{} & N(x,y)                         & \,\dd y=0\\
          & \Updownarrow                   &            & \Updownarrow                   &          \\
\dd u={}  & \dfrac{\partial u}{\partial x} & \,\dd x+{} & \dfrac{\partial u}{\partial y} & \,\dd y=0
\end{array}
\]
то есть $\dfrac{\partial u}{\partial x}=M$, $\dfrac{\partial u}{\partial y}=N$. Тогда $\dd u=0$ и
\[
\boxed{\ u(x,y)=C,\quad C\in\R\ }
\]
--- интегральные кривые.

\textbf{Необходимое условие:} $M_y=N_x$.

\textbf{Достаточные условия:} $M,N,M_y,N_x$ непрерывны в \emph{односвязной} области («без дыр») и $M_y=N_x$.

Восстановим $u$. При интегрировании по \(x\) считаем \(y\) постоянной:
\[u=\int M(x,y)\dd x+C(y).\]

Мы знаем, что \(\dfrac{\partial u}{\partial y}=N\), воспользуемся этим, для того чтобы найти \(C'(y)\) и затем \(C(y)\):
\begin{gather*}
\bigl(u\bigr)'_y=\left(\int M(x,y)\dd x\right)'_y+C'(y)=N(x,y)\ \Longrightarrow\\ \Longrightarrow C'(y)\ \Longrightarrow\ C(y).
\end{gather*}
Подставим \(C(y)\) в восстановленный \(u\).
\end{definitionHSE}

\begin{remarkHSE}
Неодносвязную область можно покрыть односвязными. Например, $\R^2\setminus\{(0,0)\}=G_1\cup G_2$, где $G_1$ --- плоскость с разрезом по лучу $\{y=0,\ x\le0\}$, а $G_2$ --- с разрезом по лучу $\{y=0,\ x\ge0\}$.
\end{remarkHSE}

\begin{remarkHSE}
Любое уравнение с разделяющимися переменными представимо как уравнение в полных дифференциалах, но не наоборот.
\end{remarkHSE}

\begin{exampleHSE}[\(\dfrac{y}{x}\dd x+(y^3+\ln x)\dd y=0\)]
\[
M=\frac yx,\quad N=y^3+\ln x,\quad M_y=\frac1x=N_x.
\]
$M,N,M_y,N_x$ непрерывны в односвязной области $G=\{(x,y):x>0\}$.
\begin{gather*}
u=\int\frac yx\dd x=y\ln x+C(y),\qquad
u'_y=\underline{\ln x+C'(y)=y^3+\ln x}=N,\\
C'(y)=y^3,\quad C(y)=\frac{y^4}{4}\quad(\text{можем не писать константу, она поглотится}).
\end{gather*}
\[
\boxed{\ y\ln x+\frac{y^4}{4}=C,\quad C\in\R\ }
\]
\end{exampleHSE}

\phantomsection\label{box:type10}
\begin{definitionHSE}[Уравнения типа 10 (введение под знак дифференциала)]
Группируем слагаемые так, чтобы выделить полные дифференциалы, и делаем замену. Полезные формулы:
\[
\dd(uv)=u\dd v+v\dd u,\qquad
\dd\Bigl(\frac uv\Bigr)=\frac{v\dd u-u\dd v}{v^2},\qquad
\dd y^2=2y\dd y,\qquad
\dd\ln y=\frac{\dd y}{y}.
\]
\end{definitionHSE}

\begin{exampleHSE}[\((x^2+y^2+x)\dd x+y\dd y=0\)]
\[
(x^2+y^2)\dd x+\frac{\dd(x^2)}{2}+\frac{\dd(y^2)}{2}=0,\qquad u=x^2+y^2:\quad u\dd x+\frac{\dd u}{2}=0.
\]
$u=0$, то есть $x^2+y^2=0$, --- не решение: это точка, а не дифференцируемая функция. При $u\neq0$:
\[
\frac{\dd u}{u}=-2\dd x,\qquad 2x=-\ln|u|-\ln C,\ C>0,\qquad \e^{2x}=\frac{1}{C(x^2+y^2)}.
\]
\[
\boxed{\ \e^{2x}(x^2+y^2)=C,\quad C>0\ }
\]
\end{exampleHSE}

\subsection{Уравнения, не разрешённые относительно производной}

\[
F(x,y,p)=0,\qquad p=y'=\frac{\dd y}{\dd x}.
\]
Если удаётся разрешить относительно $p$, получаем совокупность уравнений
\[
p=f_1(x,y),\quad\ldots,\quad p=f_n(x,y).
\]
Например, $y'^2-x^2=0\ \Longrightarrow\ y'=x$ или $y'=-x$.

Иначе используем параметр $p$: из уравнения выражаем $y$ (или $x$), берём дифференциал и подставляем $\dd y=p\dd x$.

Подобным методом решается 4 типа уравнений:
\begin{itemize}
\item уравнение Клеро;
\item уравнение Лагранжа;
\item \(y=f(y')\);
\item \(x=f(y')\).
\end{itemize}

\subsubsection{Уравнения на замену \texorpdfstring{\(y'=p\)}{y'=p}}

\phantomsection\label{box:type11}
\begin{definitionHSE}[Уравнения типа 11 (уравнение Клеро)]
\[
y=xy'+\varphi(y').
\]
$\boxed{y'=p}\Rightarrow y=xp+\varphi(p)$. Дифференцируем по $x$:
\[
p=p+xp'+\varphi'(p)p'\quad\Longrightarrow\quad p'\bigl(x+\varphi'(p)\bigr)=0.
\]
\textcircled{\small I} $p'=0$, $p=C$ --- общее решение (семейство прямых):
\[
\boxed{\ y=Cx+\varphi(C),\quad C\in\R\ }
\]
\textcircled{\small II} $x+\varphi'(p)=0$ --- особое решение (огибающая семейства прямых):
\[
\boxed{\ \begin{cases}x=-\varphi'(p),\\ y=-p\,\varphi'(p)+\varphi(p).\end{cases}}
\]
\end{definitionHSE}

\begin{exampleHSE}[\(y=xy'-y'^2\)]
\[
y'=p,\quad y=xp-p^2,\quad p=p+xp'-2pp',\quad p'(x-2p)=0.
\]
\textcircled{\small I} $p'=0\Rightarrow p=C\Rightarrow$ Подставляем в исходное уравнение: $y=Cx-C^2$.

\textcircled{\small II} $x=2p\Rightarrow p=\dfrac x2\Rightarrow$ Подставляем в исходное уравнение: \(y=\dfrac{x^2}{2}-\dfrac{x^2}{4}=\dfrac{x^2}{4}\).
\[
\boxed{\ y=Cx-C^2,\ C\in\R;\quad y=\frac{x^2}{4}\ }
\]
\end{exampleHSE}

\phantomsection\label{box:type12}
\begin{definitionHSE}[Уравнения типа 12 (уравнение Лагранжа)]
\[
y=x\,\varphi(y')+\psi(y').
\]
$\boxed{y'=p}\Rightarrow y=x\varphi(p)+\psi(p)$. Дифференцируем по $x$:
\begin{gather*}
p=\varphi(p)+x\varphi'(p)p'+\psi'(p)p',\qquad p'\bigl(x\varphi'(p)+\psi'(p)\bigr)=p-\varphi(p),\\[0.6em]
\bigl(x\varphi'(p)+\psi'(p)\bigr)\dd p=(p-\varphi(p))\dd x\\[0.6em]
\dd p=0\Longrightarrow p-\varphi(p)=0
\end{gather*}
Корни уравнения $p-\varphi(p)=0$ дают решения $y=x\varphi(p_i)+\psi(p_i)$. Если \(p-\varphi(p)\neq0\), получаем \hyperref[box:type6]{линейное} по $x$ уравнение:
\[
\boxed{\ x'_p-\frac{\varphi'(p)}{p-\varphi(p)}\,x=\frac{\psi'(p)}{p-\varphi(p)}\ }
\]
Ответ в параметрической форме: $\begin{cases}x=x(p),\\ y=x(p)\,\varphi(p)+\psi(p).\end{cases}$
\end{definitionHSE}

\begin{exampleHSE}[\(y+xy'=4\sqrt{y'}\)]
$y'=p\ge0$, $y=-xp+4\sqrt p$. Дифференцируем:
\[
p=-p-xp'+\frac{2p'}{\sqrt p}\quad\Longrightarrow\quad 2p=p'\Bigl(\frac{2}{\sqrt p}-x\Bigr).
\]
$p=0\ \Longrightarrow\ y=0$ --- решение. При $p\neq0$, $p'=\frac{1}{x'_p}$:
\[
2p\,x'_p=\frac{2}{\sqrt p}-x\quad\Longrightarrow\quad x'_p+\frac{x}{2p}=\frac{1}{p\sqrt p}.
\]
Однородное:
\[\frac{\dd x}{x}=-\frac{\dd p}{2p},\quad x=\frac{C}{\sqrt p}.\]
Вариация $C=C(p)$:
\[
\frac{C'}{\sqrt p}=\frac{1}{p\sqrt p},\quad C'=\frac1p,\quad C=\ln p+\widetilde C.
\]
\[
\boxed{\ \begin{cases}x=\dfrac{\ln p+C}{\sqrt p},\\[1.5ex] y=(4-\ln p-C)\sqrt p,\end{cases}\ C\in\R;\qquad y=0\ }
\]
\end{exampleHSE}

\phantomsection\label{box:type13}
\begin{definitionHSE}[Уравнения типа 13: \(y=f(y')\)]
$\boxed{y'=p}\Rightarrow y=f(p)$, $\dd y=p\dd x$:
\[
p\dd x=f'(p)\dd p,\quad \frac{\dd x}{\dd p}=\frac{f'(p)}{p},\ p\neq0\quad\Longrightarrow\quad x=\int\frac{f'(p)}{p}\dd p+C.
\]
\[
\boxed{\ \begin{cases}x=\displaystyle\int\frac{f'(p)}{p}\dd p+C,\\ y=f(p);\end{cases}\quad y=f(0)\ \text{--- тоже решение.}}
\]
\end{definitionHSE}

\begin{exampleHSE}[\(3y'^4=y'+y\)]
\[
y=3y'^4-y',\qquad y=0\ \text{--- решение}\ (p=0).
\]
$y'=p\Rightarrow y=3p^4-p$, $p\dd x=\dd y=(12p^3-1)\dd p$:
\[
\frac{\dd x}{\dd p}=12p^2-\frac1p\quad\Longrightarrow\quad x=4p^3-\ln|p|+C.
\]
\[
\boxed{\ \begin{cases}x=4p^3-\ln|p|+C,\\ y=3p^4-p,\end{cases}\ C\in\R;\qquad y=0\ }
\]
\end{exampleHSE}

\phantomsection\label{box:type14}
\begin{definitionHSE}[Уравнения типа 14: \(x=f(y')\)]
$\boxed{y'=p}\Rightarrow x=f(p)$, $\dd x=f'(p)\dd p$; из $\dd y=p\dd x$:
\[
\dd y=p\,f'(p)\dd p\quad\Longrightarrow\quad
\boxed{\ \begin{cases}y=\displaystyle\int p\,f'(p)\dd p+C,\ C\in\R,\\ x=f(p).\end{cases}}
\]
\end{definitionHSE}

\begin{exampleHSE}[\(2y'=x+\ln y'\) (№360)]
$x=2y'-\ln y'$; $y'=p>0$, $x=2p-\ln p$:
\[
\dd y=p\dd x=p\Bigl(2-\frac1p\Bigr)\dd p=(2p-1)\dd p,\quad y=p^2-p+C.
\]
\[
\boxed{\ \begin{cases}x=2p-\ln p,\\ y=p^2-p+C,\end{cases}\quad C\in\R\ }
\]
\end{exampleHSE}

\subsubsection{Уравнения старшего порядка}
\phantomsection\label{box:type15}
\begin{definitionHSE}[Уравнения типа 15 (понижение порядка: нет младших производных)]
\[
F\bigl(y^{(n)},y^{(n-1)},\ldots,y^{(n-k)},x\bigr)=0.
\]
Замена \textbf{младшей} из существующих производных $z=y^{(n-k)}$ понижает порядок уравнения на её порядок $n-k$:
\[
F\bigl(z^{(k)},z^{(k-1)},\ldots,z,x\bigr)=0,\quad z=\varphi(x),\quad y^{(n-k)}=\varphi(x).
\]
\end{definitionHSE}

\begin{exampleHSE}[\(y'''=y''^2\)]
Замена $z=y''$: $z'=z^2$.

$z=0$: $y''=0$, $y=C_1x+C_2$.

$z\neq0$:
\[\frac{\dd z}{\dd x}=z^2,\quad\frac{\dd z}{z^2}=\dd x,\quad-\frac1z=x+C,\]
то есть $y''=-\dfrac{1}{x+C}$.

Интегрируем дважды:
\[
y'=-\ln|x+C|+C_1,\qquad y=-(x+C)\ln|x+C|+(x+C)+C_1x+C_2.
\]
\[
\boxed{\ y=-(x+C)\ln|x+C|+C_1x+C_2;\qquad y=C_1x+C_2,\qquad C,C_1,C_2\in\R\ }
\]
\end{exampleHSE}

\phantomsection\label{box:type16}
\begin{definitionHSE}[Уравнения типа 16 (не содержащие явно \(x\))]
\[
F\bigl(y^{(n)},y^{(n-1)},\ldots,y',y\bigr)=0.
\]
Замена $y'=p(y)$ понижает порядок на единицу:
\[
y''=\frac{\dd}{\dd x}\,p=\frac{\dd p}{\dd y}\cdot\underbrace{\frac{\dd y}{\dd x}}_{=p}=p'_y\,p.
\]
\end{definitionHSE}

\begin{exampleHSE}[\(y^3y''=1\)]
$y'=p(y)$, $y''=p'_yp$:
\[
y^3p\,p'_y=1,\quad p\dd p=\frac{\dd y}{y^3},\quad p^2=-\frac{1}{y^2}+C_1,\quad
y'=\pm\sqrt{\frac{C_1y^2-1}{y^2}}.
\]
Стационарных решений нет, иначе левая часть бы обращалась в нуль.
\[
\frac{y\dd y}{\sqrt{C_1y^2-1}}=\pm\dd x,\quad \frac{\sqrt{C_1y^2-1}}{C_1}=\pm(x+C_2).
\]
\[
\boxed{\ C_1^2(x+C_2)^2=C_1y^2-1,\quad C_1>0,\ C_2\in\R\ }
\]
\end{exampleHSE}

\phantomsection\label{box:thm-n}
\begin{theoremHSE}[Существование и единственность для уравнения \(n\)-го порядка]
Рассмотрим задачу Коши
\[
\begin{cases}
y^{(n)}=F\bigl(x,y,y',\ldots,y^{(n-1)}\bigr),\\
y(x_0)=y_0,\\
y'(x_0)=y_1,\\
\ldots,\\
y^{(n-1)}(x_0)=y_{n-1}.
\end{cases}
\]
\textcircled{\small 1} \(F,\ F'_y,\ F'_{y'},\ \ldots,\ F'_{y^{(n-1)}}\in C(\Omega),\ \Omega-\text{область}\subset\R^{n+1}.\)

\textcircled{\small 2} \((x_0,y_0,y_1,\ldots,y_{n-1})\in\Omega.\)

Если существует подобная область $\Omega$, и выполнены указанные условия, то задача Коши имеет, причём единственное решение, определённое на \([x_0-h,x_0+h]\).
\end{theoremHSE}

\begin{exampleHSE}[\(y''\cos y+y'^2\sin y=y',\ y(-1)=\dfrac\pi6,\ y'(-1)=2\)]
\textbf{Существование и единственность.} Необходимо разрешить уравнение относительно старшей производной. При $\cos y\neq0$
\[
y''=F=\frac{y'-y'^2\sin y}{\cos y},\qquad
F'_y=\frac{y'\sin y-y'^2}{\cos^2y},\qquad
F'_{y'}=\frac{1-2y'\sin y}{\cos y}.
\]ы
В области $\Omega=\bigl\{(x,y,y')\mid 0<y<\frac\pi2\bigr\}$, содержащей $\bigl(-1,\frac\pi6,2\bigr)$, функции $F,F'_y,F'_{y'}$ непрерывны $\Longrightarrow$ решение существует и единственно.

\textbf{Решение.} Так как $\Bigl(\dfrac{y'}{\cos y}\Bigr)'=\dfrac{y''\cos y+y'^2\sin y}{\cos^2y}$, после деления на $\cos^2y$:
\[
\Bigl(\frac{y'}{\cos y}\Bigr)'=\frac{y'}{\cos^2y}=(\tg y)'\quad\Longrightarrow\quad\frac{y'}{\cos y}=\tg y+C.
\]
Из начальных условий $C=\sqrt3$, поэтому $y'=\sin y+\sqrt3\cos y=2\sin\bigl(y+\frac\pi3\bigr)$ и
\[
\boxed{\ y=2\arctg\e^{2(x+1)}-\frac\pi3\ }
\]
\end{exampleHSE}

\phantomsection\label{box:type17}
\begin{definitionHSE}[Уравнения типа 17 (однородные относительно \(y\) и её производных)]
Если уравнение не меняется при замене $y\to ky,\ y'\to ky',\ \ldots,\ y^{(n)}\to ky^{(n)}$, то порядок понижается подстановкой
\[
y'=y\cdot z(x)\Longrightarrow y''=y'z+yz'=y\,(z^2+z').
\]
\end{definitionHSE}

\begin{exampleHSE}[\(y'^2+2yy''=0\)]
$y=0$ --- решение. Замена $y'=yz$:
\[
y^2z^2+2y^2(z^2+z')=0\quad\Longrightarrow\quad 3z^2+2z'=0.
\]
$z=0$: $y'=0$, $y=C$. При $z\neq0$:
\[
3\dd x=-\frac{2\dd z}{z^2},\quad 3x=\frac2z-C_1,\quad \frac{y'}{y}=z=\frac{2}{3x+C_1},\quad \ln|y|=\frac23\ln|3x+C_1|+\widetilde C.
\]
\[
\boxed{\ y=C_2\,(3x+C_1)^{2/3},\ C_1,C_2\in\R;\qquad y=C,\ C\in\R\ }
\]
\end{exampleHSE}

\phantomsection\label{box:type18}
\begin{definitionHSE}[Уравнения типа 18 (выделение полной производной)]
Левую часть приводим к виду $\bigl(\Phi(x,y,y',\ldots)\bigr)'=0$, откуда $\Phi=C$, и порядок понижается на единицу. Например,
\[
y'y''+yy'''=(yy'')',\qquad 2y'y''=(y'^2)'.
\]
\end{definitionHSE}

\begin{exampleHSE}[\(yy'''+3y'y''=0\)]
\[(yy'')'=y'y''+yy''',\quad yy'''=(yy'')'-y'y''\]
\begin{gather*}
(yy'')'+2y'y''=0,\quad (yy'')'+(y'^2)'=0,\quad (yy''+y'^2)'=0,\\
yy''+y'^2=C,\quad (yy')'=C,\quad yy'=Cx+C_1,\quad \frac12(y^2)'=Cx+C_1.
\end{gather*}
\[
\boxed{\ y^2=C_1x^2+C_2x+C_3,\quad C_1,C_2,C_3\in\R\ }
\]
\end{exampleHSE}

\phantomsection\label{box:type19}
\begin{definitionHSE}[Уравнения типа 19 (обобщённо-однородные)]
Пусть уравнение не меняется при
\[
x\to kx,\quad y\to k^my,\quad y'\to k^{m-1}y',\quad \ldots,\quad y^{(l)}\to k^{m-l}y^{(l)}.
\]
Тогда делаем замену
\[
\begin{cases}x=\e^t,\\ y=z\,\e^{mt},\end{cases}\quad z=z(t).
\]
Например, при $m=2$:
\[
y'=\e^t(z'_t+2z),\qquad y''=z''_{tt}+3z'_t+2z.
\]
После замены уравнение не содержит $t$ явно (\hyperref[box:type16]{тип 16}).
\end{definitionHSE}

\begin{exampleHSE}[\(x^2y''-3xy'=\dfrac{6y^2}{x^2}-4y,\ y(1)=1,\ y'(1)=4\)]
Подбираем $m$: $x^2y''\sim k^m$, $xy'\sim k^m$, $y\sim k^m$, $\frac{y^2}{x^2}\sim k^{2m-2}$, откуда $2m-2=m$, $m=2$.

Замена $x=\e^t$, $y=z\e^{2t}$:
\[
\e^{2t}(z''+3z'+2z)-3\e^{2t}(z'+2z)=6z^2\e^{2t}-4z\e^{2t}\quad\Longrightarrow\quad z''=6z^2.
\]
Начальные условия: $t=0$, $z(0)=1$, $z'(0)=y'(1)-2z(0)=2$. Умножаем на $z'$:
\[
(z'^2)'=(4z^3)',\qquad z'^2=4z^3+C,\quad C=0,\qquad z'=2z^{3/2},\qquad z^{-1/2}=1-t.
\]
\[
z=\frac{1}{(1-t)^2}\quad\Longrightarrow\quad
\boxed{\ y=\frac{x^2}{(1-\ln x)^2}\ }
\]
\end{exampleHSE}

\clearpage
\phantomsection
\subsection{Схема решения произвольного ОДУ}
\label{sec:ode-scheme}

\begin{center}
\begin{tikzpicture}[
    font=\small,
    >={Stealth[length=2.2mm]},
    odestep/.style={draw=HSEblue, thick, rounded corners=4pt, fill=HSEblue!6,
                 align=center, inner sep=5pt},
    odeask/.style={draw=HSEorange!85!black, thick, rounded corners=10pt, fill=HSEorange!10,
                align=center, inner sep=5pt},
    odeleaf/.style={draw=HSEgreen!70!black, thick, rounded corners=4pt, fill=HSEgreen!6,
                 align=flush left, inner sep=5pt, text width=4.95cm, anchor=north},
    odefin/.style={draw=HSEred!80!black, thick, rounded corners=4pt, fill=HSEred!5,
                align=flush left, inner sep=6pt, text width=16.2cm},
    odelab/.style={font=\footnotesize\itshape, fill=white, inner sep=1.5pt}
  ]
  % --- шаг 0 ---
  \node[odestep, text width=15cm] (start) at (0,0)
    {\textbf{Шаг 0.} Привести уравнение к удобному виду, определить порядок $n$.\\
     Сразу выписать всё, на что придётся делить, и проверить, не решения ли это
     (стац. решения $y=\text{const}$, а для $M\dd x+N\dd y=0$ также $x=\text{const}$).};

  % --- вопросы ---
  \node[odeask] (q1) at (0,-1.75) {Порядок $n=1$?};
  \node[odeask, text width=5.4cm] (q2) at (-2.95,-3.35)
    {Разрешено относительно $y'$\\ (или вид $M\dd x+N\dd y=0$)?};

  % --- листья ---
  \node[odeleaf] (A) at (-5.6,-4.75) {%
    \textbf{Первый порядок, $y'=f(x,y)$}\\[2pt]
    1) $y'=f(x)$ --- \hyperref[box:type0]{тип~0};\\
    2) $y'=f(x)g(y)$ --- \hyperref[box:type1]{тип~1};\\
    3) $y'=f(ax+by+c)$ --- \hyperref[box:type2]{тип~2};\\
    4) $y'=f\bigl(\frac yx\bigr)$ или $M,N$ однородны одного порядка --- \hyperref[box:type3]{тип~3};\\
    5) $y'=f\Bigl(\frac{ax+by+c}{a_1x+b_1y+c_1}\Bigr)$: при $\Delta\neq0$ --- \hyperref[box:type4]{тип~4}, при $\Delta=0$ --- \hyperref[box:type2]{тип~2};\\
    6) «почти» однородно: однородно после $y=z^m$ --- \hyperref[box:type5]{тип~5};\\
    7) линейно по $y$ --- \hyperref[box:type6]{тип~6}; линейно по $x(y)$ --- \hyperref[box:type6-1]{тип~6.1};\\
    8) $y'+a(x)y=b(x)y^n$ --- \hyperref[box:type8]{тип~8};\\
    9) $M_y=N_x$ --- \hyperref[box:type9]{тип~9};\\
    10) видны $\dd(x^2+y^2)$, $\dd\bigl(\frac xy\bigr)$, \ldots\ --- \hyperref[box:type10]{тип~10}.};

  \node[odeleaf] (B) at (0,-4.75) {%
    \textbf{Первый порядок, $F(x,y,y')=0$}\\[2pt]
    0) Разрешить относительно $y'$: получится совокупность \(y'=f(x,y)\);\\
    1) $y=xy'+\varphi(y')$ --- \hyperref[box:type11]{тип~11} (Клеро);\\
    2) $y=x\varphi(y')+\psi(y')$ --- \hyperref[box:type12]{тип~12} (Лагранж);\\
    3) $y=f(y')$ --- \hyperref[box:type13]{тип~13};\\
    4) $x=f(y')$ --- \hyperref[box:type14]{тип~14}.\\[3pt]
    \textit{В общем случае:} $y'=p$, выразить $y$ или $x$, взять дифференциал и подставить $\dd y=p\dd x$.};

  \node[odeleaf] (C) at (5.6,-4.75) {%
    \textbf{Высший порядок, $n\ge2$}\\[2pt]
    0) Задача Коши --- проверить \hyperref[box:thm-n]{теорему $\exists!$};\\
    1) нет мл. производных --- $z=y^{(k)},k=\min\{i:y^{(i)}\}$, \hyperref[box:type15]{тип~15};\\
    2) нет $x$: $y'=p(y)$, \hyperref[box:type16]{тип~16};\\
    3) однородно по $y,y',\ldots$ --- $y'=yz(x)$, \hyperref[box:type17]{тип~17};\\
    4) обобщённо-однородно --- $x=\e^t$, $y=z\e^{mt}$, \hyperref[box:type19]{тип~19};\\
    5) левая часть --- полная производная --- \hyperref[box:type18]{тип~18}.\\[3pt]
    \textit{После понижения порядка --- снова по схеме.}};

  % --- стрелки ---
  \draw[->, thick, HSEgray] (start) -- (q1);
  \draw[->, thick, HSEgray] (q1.west) -| node[odelab, pos=0.25] {да} (q2.north);
  \draw[->, thick, HSEgray] (q1.east) -| node[odelab, pos=0.25] {нет} (C.north);
  \draw[thick, HSEgray] (q2.south) -- ++(0,-0.3) coordinate (br);
  \draw[->, thick, HSEgray] (br) -| node[odelab, pos=0.3] {да} (A.north);
  \draw[->, thick, HSEgray] (br) -| node[odelab, pos=0.3] {нет} (B.north);

  % --- итог ---
  \path let \p1 = (A.south), \p2 = (B.south), \p3 = (C.south)
        in coordinate (low) at (0, {min(\y1,\y2,\y3)});
  \node[odefin, anchor=north] (F) at ($(low)+(0,-0.8)$) {%
    \textbf{Итог.} 1) Сделать обратную замену. 2) Добавить стационарные и особые решения; проверить, что ничего не потеряли при делении. 3) Где можно, объединить семейства (часто стац. решение получается при $C=0$). 4) Для задачи Коши найти константы из начальных условий.};
  \foreach \n in {A,B,C} \draw[->, thick, HSEgray] (\n.south) -- (\n.south |- F.north);
\end{tikzpicture}
\end{center}

\begin{remarkHSE}[Если ни один тип не подошёл]
\begin{itemize}
  \item Поменять ролями $x$ и $y$: уравнение для $x(y)$ может оказаться линейным, Бернулли и т.\,д.
  \item Попробовать замену, подсказанную видом уравнения: $u=\frac yx$, $z=ax+by$, $z=y^m$, $z=\varphi(y)$.
  \item Для $M\dd x+N\dd y=0$ при $M_y\neq N_x$ искать интегрирующий множитель $\mu(x)$ или $\mu(y)$.
\end{itemize}
\end{remarkHSE}

\section{Алгоритмы решения ЛУ с пост. коэф-ами}
\label{sec:lu-const-scheme}
Решаем задачи где правая часть --- вещественная функция и коэффициенты уравнения тоже вещественные. Это позволит нам сделать небольшой трюк:
\[\boxed{L(D)\Re Y=\Re(L(D)Y),}\qquad L(D)\Im Y=\Im(L(D)Y)\]
По формуле Эйлера:
\[e^{ibx}=\cos bx+i\sin bx\]
Разрешим многочлену в \(\e^{\mu x}P_N(x)\) иметь комплексные коэффициенты и запишем его как \(P(x)-iR(x)\). Тогда квазиодночлен в правой части принимает следующий вид:
\begin{gather*}
\e^{\mu x}P_N(x)=\e^{(a+ib)x}P_N(x)=\e^{(a+ib)x}(P(x)-iR(x))\\
\Re\bigl[\e^{(a+bi)x}(P(x)-iR(x)\bigr]=\e^{ax}(P(x)\cos bx+R(x)\sin bx)
\end{gather*}
Комплексное частное решение имеет вид:
\begin{gather*}
Y=x^s\e^{\mu x}Q_N(x)=x^se^{(a+ib)x}(Q_N(x)-iS_N(x))\\
y=\Re Y=x^se^{ax}(Q_N(x)\cos bx+S_N(x)\sin bx)
\end{gather*}

Временно решаем вспомогательное уравнение
\[L(D) Y=f(x)+ig(x)\]
\(\Re Y\) и \(\Im Y\) --- действительные функции. Подставим в уравнение:
\begin{gather*}
L(D)Y=L(D)(\Re Y+i\Im Y)=L(D)\Re Y+iL(D)\Im Y\\[1em]
L(D)Y=f(x)+ig(x)\Rightarrow L(D)\Re Y+iL(D)\Im Y=f(x)+ig(x)\Rightarrow\\[1em]
\Rightarrow\begin{cases}\boxed{L(D)\Re Y=f(x)}\\ L(D)\Im Y=g(x)\end{cases}
\end{gather*}

Схема покрывает ЛНУ с постоянными коэффициентами и правой частью --- квазимногочленом и уравнения Эйлера, которые сводятся к ним.

\begin{center}
\begin{tikzpicture}[
    font=\small,
    >={Stealth[length=2.2mm]},
    odestep/.style={draw=HSEblue, thick, rounded corners=4pt, fill=HSEblue!6,
                 align=flush left, inner sep=5pt, text width=15.6cm},
    odeask/.style={draw=HSEorange!85!black, thick, rounded corners=10pt, fill=HSEorange!10,
                align=center, inner sep=5pt},
    odeleaf/.style={draw=HSEgreen!70!black, thick, rounded corners=4pt, fill=HSEgreen!6,
                 align=flush left, inner sep=5pt, anchor=north},
    odefin/.style={draw=HSEred!80!black, thick, rounded corners=4pt, fill=HSEred!5,
                align=flush left, inner sep=6pt, text width=15.6cm},
    odelab/.style={font=\footnotesize\itshape, fill=white, inner sep=1.5pt}
  ]
  % --- шаг 0 ---
  \node[odestep] (s0) at (0,0) {%
    \textbf{Шаг 0.} Записать уравнение в виде $L(D)y=f(x)$, где $L(D)=a_nD^n+\cdots+a_1D+a_0$, $a_i=\mathrm{const}$, $D=\frac{\dd}{\dd x}$.\\
    Уравнение Эйлера $a_nx^ny^{(n)}+\cdots+a_1xy'+a_0y=f(x)$ сначала свести к такому виду заменой $x=\e^t$ (см. \hyperref[box:euler]{коробочку под схемой}).};

  % --- шаг 1 ---
  \node[odestep, anchor=north] (s1) at ($(s0.south)+(0,-0.55)$) {%
    \textbf{Шаг 1. ЛОУ $L(D)y=0$} (разделы~\ref{sec:roots-simple}--\ref{sec:roots-multiple}). Характеристический многочлен $L(\lambda)$ $\to$ корни и кратности $\to$ ФСР $\to$ $y_{\text{о.о.}}=c_1y_1+\cdots+c_ny_n$:
    \[e^{\lambda x}=e^{(\alpha+i\beta)x}=e^{\alpha x}\cdot e^{i\beta x}=e^{\alpha x}(\cos \beta x+i\sin\beta x)\]
    \begin{itemize}[nosep, leftmargin=1.6em, label=$\bullet$]
    \item вещественный корень $\lambda$ кратности $k$: \ $\e^{\lambda x},\ x\e^{\lambda x},\ \ldots,\ x^{k-1}\e^{\lambda x}$;
    \item так как \(L(\lambda)\in\mathbb{R}[\lambda]\), комплексные корни идут сопряжёнными парами: для $\alpha\pm i\beta$ кратности $k$: \ $x^r\e^{\alpha x}\cos\beta x,\ x^r\e^{\alpha x}\sin\beta x$, \ $r=0,\ldots,k-1$.
    \end{itemize}};

  % --- шаг 2 ---
  \node[odestep, anchor=north] (s2) at ($(s1.south)+(0,-0.55)$) {%
    \textbf{Шаг 2. Правая часть.} Преобразовать $f$ в сумму слагаемых и искать частное решение для каждого отдельно (принцип суперпозиции):
    \begin{itemize}[nosep, leftmargin=1.6em, label=$\bullet$]
    \item произведения $\sin\cdot\cos$, $\cos\cdot\cos$, $\sin\cdot\sin$ и степени $\sin^2$, $\cos^2$ --- в суммы по тригонометрическим формулам;
    \item $\operatorname{sh}x=\frac{\e^x-\e^{-x}}{2}$, \ $\operatorname{ch}x=\frac{\e^x+\e^{-x}}{2}$;
    \item многочлен $P(x)$ --- это $\e^{0\cdot x}P(x)$, то есть $\mu=0$.
    \end{itemize}};

  % --- вопрос ---
  \node[odeask, anchor=north] (q) at ($(s2.south)+(0,-0.6)$)
    {Каждое слагаемое имеет вид\\ $\e^{ax}\bigl(P_m(x)\cos bx+Q_n(x)\sin bx\bigr)$?};

  % --- листья ---
  \node[odeleaf, text width=9.3cm] (A) at ($(-3.2,0 |- q.south)+(0,-0.9)$) {%
    \textbf{Квазимногочлен} (раздел~\ref{sec:quasi}). Для каждого слагаемого знаем \(a,b,P(x),R(x)\):
    \begin{enumerate}[nosep, leftmargin=1.6em, label=\arabic*)]
    \item $N=\max(m, n)$;
    \item $s$ = кратность $\mu=a+ib$ как корня $L(\lambda)$;
    \item вид частного решения:\\
      $y_j=x^s\e^{ax}\bigl(\tilde{P}_N(x)\cos bx+\tilde{Q}_N(x)\sin bx\bigr)$,\\
      при $b=0$: \ $y_j=x^s\e^{ax}Q_N(x)$.\\
      Если в слагаемом только $\cos bx$ (только $\sin bx$), можно решить $L(D)Y=\e^{\mu x}P(x)$ и взять $y_j=\Re Y$ ($y_j=\Im Y$);
    \item подставить и приравнять коэффициенты; удобно через формулу смещения $L(D)\bigl(\e^{\mu x}u\bigr)=\e^{\mu x}L(D+\mu)u$.
    \end{enumerate}};

  \node[odeleaf, text width=5.6cm] (B) at ($(5.1,0 |- q.south)+(0,-0.9)$) {%
    \textbf{Другая $f$} (например, $\operatorname{tg}x$, $\frac1{\cos x}$, $\ln x$):
    \begin{enumerate}[nosep, leftmargin=1.6em, label=\arabic*)]
    \item вариация постоянных (раздел~\ref{sec:variation});
    \item или интегральная формула Коши (раздел~\ref{sec:cauchy}):\\
      $Y(x)=\int_\gamma^x K(x,s)f(s)\dd s$.
    \end{enumerate}
    \vspace{3pt}
    Ответ --- на промежутках непрерывности $f$.};

  % --- стрелки ---
  \draw[->, thick, HSEgray] (s0) -- (s1);
  \draw[->, thick, HSEgray] (s1) -- (s2);
  \draw[->, thick, HSEgray] (s2) -- (q);
  \draw[->, thick, HSEgray] (q.west) -| node[odelab, pos=0.75, left, fill=none] {да} (A.north);
  \draw[->, thick, HSEgray] (q.east) -| node[odelab, pos=0.75, right, fill=none] {нет} (B.north);

  % --- итог ---
  \path let \p1 = (A.south), \p2 = (B.south)
        in coordinate (low) at (0, {min(\y1,\y2)});
  \node[odefin, anchor=north] (F) at ($(low)+(0,-0.8)$) {%
    \textbf{Итог.} 1) $y_{\text{о.н.}}=y_{\text{о.о.}}+y_1+\cdots+y_m$. 2) Для уравнения Эйлера --- обратная замена $t=\ln|x|$; ответ отдельно на $x>0$ и $x<0$, где правая часть определена. 3) Проверить частное решение подстановкой в уравнение (хотя бы старшие члены).};
  \foreach \n in {A,B} \draw[->, thick, HSEgray] (\n.south) -- (\n.south |- F.north);
\end{tikzpicture}
\end{center}

\phantomsection\label{box:euler}
\begin{definitionHSE}[Уравнение Эйлера]
Уравнение
\[
a_nx^ny^{(n)}+a_{n-1}x^{n-1}y^{(n-1)}+\cdots+a_1xy'+a_0y=f(x),\qquad a_i=\mathrm{const},
\]
называется \textit{уравнением Эйлера}. При \(x>0\) замена \(x=\e^t\) (при \(x<0\) --- \(x=-\e^t\)), то есть \(t=\ln|x|\), сводит его к ЛУ с постоянными коэффициентами относительно \(y(t)\):
\[
xy'=Dy,\quad x^2y''=D(D-1)y,\quad\ldots,\quad x^ky^{(k)}=D(D-1)\cdots(D-k+1)y,\qquad D=\frac{\dd}{\dd t}.
\]
\end{definitionHSE}

\begin{proof}[Вывод формул]
При \(x=\pm\e^t\) имеем \(\frac{\dd x}{\dd t}=x\), поэтому \(\frac{\dd}{\dd x}=\frac1x\frac{\dd}{\dd t}\). Отсюда \(y'=\frac1xDy\), то есть \(xy'=Dy\), и
\[
y''=\frac1x\frac{\dd}{\dd t}\Bigl(\frac1xDy\Bigr)=\frac1x\Bigl(\frac1xD^2y-\frac1{x^2}\cdot x\cdot Dy\Bigr)=\frac1{x^2}\bigl(D^2-D\bigr)y,
\]
то есть \(x^2y''=D(D-1)y\). Общая формула доказывается по индукции тем же способом.
\end{proof}

\begin{remarkHSE}[Короткий путь к характеристическому многочлену]
Для уравнения Эйлера характеристический многочлен можно получить сразу, подставив \(y=x^\lambda\) в левую часть: \(x^k\bigl(x^\lambda\bigr)^{(k)}=\lambda(\lambda-1)\cdots(\lambda-k+1)\,x^\lambda\). Это то же самое, что заменить \(D\) на \(\lambda\) в формулах выше.
\end{remarkHSE}

\begin{remarkHSE}[Частые ошибки]
\begin{itemize}
\item \(s\) считается для \(\mu=a+ib\) целиком: для \(\e^{x}\cos3x\) это \(\mu=1+3i\), а не \(1\) и не \(3i\).
\item Если \(b\neq0\), в виде частного решения нужны и \(\cos bx\), и \(\sin bx\) с многочленами одной степени \(N\), даже если в правой части только косинус или только синус.
\item Многочлены \(Q_N\), \(S_N\) --- полные: неизвестны все коэффициенты при \(x^N,\ldots,x^0\), даже если в \(P\) каких-то степеней нет.
\item При суперпозиции \(\mu\), \(N\) и \(s\) находятся для каждого слагаемого отдельно.
\item Многочлен или константа в правой части --- это \(\mu=0\). Если в уравнении нет члена с самой \(y\), то \(L(0)=0\), и \(s\ge1\).
\end{itemize}
\end{remarkHSE}

\begin{exampleHSE}[Задачи Филиппова: что применить]
\renewcommand{\arraystretch}{1.3}
\begin{tabular}{@{}p{1.1cm}p{3.7cm}p{10.1cm}@{}}
\textbf{№} & \textbf{Правая часть} & \textbf{Приёмы}\\
540 & \(x\cos x\) & квазимногочлен с комплексным \(\mu\); \(\Re\)\\
541 & \(\e^{-4x}+x\e^{-x}\) & суперпозиция двух слагаемых\\
542 & \(x^2\e^{x}\) & простейший квазимногочлен\\
543 & \(\e^{2x}+\sin2x\) & суперпозиция; для \(\sin2x\) --- \(\Im\)\\
544 & \(\e^{3x}\cos x\) & комплексное \(\mu\); \(\Re\)\\
545 & \(6x\e^{x}\) & простейший квазимногочлен\\
546 & \(x\sin x\) & комплексное \(\mu\); \(\Im\)\\
547 & \(x\e^{2x}\) & простейший квазимногочлен\\
548 & \(3x^2+\sin5x\) & суперпозиция; многочлен --- это \(\mu=0\)\\
569 & \(\sin x\cos2x\) & порядок 4; произведение \(\to\) сумма; суперпозиция\\
573 & \(\operatorname{sh}x\sin2x\) & \(\operatorname{sh}x\) через экспоненты; суперпозиция\\
596 & \(3x^2\) (Эйлер) & замена \(x=\e^t\), затем квазимногочлен\\
598 & \(\sin\ln x\) (Эйлер) & замена \(x=\e^t\): \(\sin\ln x=\sin t\); \(x>0\)\\
\end{tabular}

Во всех задачах \(s\) проверяется для каждого слагаемого отдельно, сравнением \(\mu\) с корнями \(L(\lambda)\).
\end{exampleHSE}


\end{document}
